    \chapter{Set Theory and Logic}

    \relatedcourses{(MATH202) Set Theory}

    \relatedbooks{Introduction to Set Theory (Karel Hrbacek, Thomas Jech)}

    \begin{definition}[Set and Element]
        A set is a collection of objects, called its elements. We write
        $x\in A$ if $x$ is an element of $A$, and $x\notin A$ otherwise.
        The relation $A\subseteq B$ means
        \[
            \forall x,\qquad x\in A\Rightarrow x\in B.
        \]
    \end{definition}

    \begin{formula}[Basic Set Operations]
        For sets $A$ and $B$ inside an ambient set $U$, we use
        \[
            A\cup B=\{x\in U\mid (x\in A)\lor(x\in B)\},
        \]
        \[
            A\cap B=\{x\in U\mid (x\in A)\land(x\in B)\},
        \]
        and
        \[
            A\setminus B=\{x\in A\mid x\notin B\}.
        \]
    \end{formula}

    \begin{formula}[Set Algebra]
        For subsets $A,B,C\subseteq U$, the following identities are often useful:
        \[
            A\cup\varnothing=A,
            \qquad
            A\cap U=A,
            \qquad
            A\cup A=A,
            \qquad
            A\cap A=A.
        \]
        The distributive laws are
        \[
            A\cap(B\cup C)=(A\cap B)\cup(A\cap C),
        \]
        \[
            A\cup(B\cap C)=(A\cup B)\cap(A\cup C).
        \]
        The absorption laws are
        \[
            A\cup(A\cap B)=A,
            \qquad
            A\cap(A\cup B)=A.
        \]
    \end{formula}

    \begin{formula}[Basic Logical Equivalences]
        The most useful logical equivalences are
        \[
            (P\Rightarrow Q)\Longleftrightarrow(\neg Q\Rightarrow\neg P),
        \]
        \[
            \neg(P\Rightarrow Q)\Longleftrightarrow P\land\neg Q,
        \]
        \[
            \neg(P\land Q)\Longleftrightarrow(\neg P)\lor(\neg Q),
            \qquad
            \neg(P\lor Q)\Longleftrightarrow(\neg P)\land(\neg Q),
        \]
        and
        \[
            \neg(\forall x\in A,\ P(x))
            \Longleftrightarrow
            \exists x\in A\ \text{s.t.}\ \neg P(x),
        \]
        \[
            \neg(\exists x\in A\ \text{s.t.}\ P(x))
            \Longleftrightarrow
            \forall x\in A,\ \neg P(x).
        \]
    \end{formula}

    \begin{strategy}
        To prove $P\Rightarrow Q$, it is often enough to prove the contrapositive
        $\neg Q\Rightarrow\neg P$. To disprove a universal statement, find a
        counterexample.
    \end{strategy}

    \begin{formula}[De Morgan's Laws]
        For subsets $A,B\subseteq U$, we have
        \[
            (A\cup B)^{\mathsf c}=A^{\mathsf c}\cap B^{\mathsf c},
            \qquad
            (A\cap B)^{\mathsf c}=A^{\mathsf c}\cup B^{\mathsf c}.
        \]
        More generally,
        \[
            \left(\bigcup_{i\in I}A_i\right)^{\mathsf c}
            =
            \bigcap_{i\in I}A_i^{\mathsf c},
            \qquad
            \left(\bigcap_{i\in I}A_i\right)^{\mathsf c}
            =
            \bigcup_{i\in I}A_i^{\mathsf c}.
        \]
    \end{formula}

    \begin{definition}[Power Set]
        The power set of a set $A$ is
        \[
            \Pow(A)=\{B\mid B\subseteq A\}.
        \]
        If $|A|=n<\infty$, then
        \[
            |\Pow(A)|=2^n.
        \]
    \end{definition}

    \begin{definition}[Function]
        A function $f:X\to Y$ assigns to each element $x\in X$ a unique element
        $f(x)\in Y$. The set $X$ is the domain, and $Y$ is the codomain. The
        notation $x\mapsto f(x)$ describes the rule on elements.
    \end{definition}

    \begin{definition}[Injection, Surjection, and Bijection]
        A function $f:X\to Y$ is injective if distinct elements of $X$ have
        distinct images. Equivalently,
        \[
            \forall x_1,x_2\in X,
            \qquad
            f(x_1)=f(x_2)\implies x_1=x_2.
        \]
        Once injectivity has been established, one may write
        $f:X\hookrightarrow Y$.

        The function is surjective if every element of $Y$ has a preimage:
        \[
            \forall y\in Y,
            \qquad
            \exists x\in X\ \text{s.t.}\ f(x)=y.
        \]
        Once surjectivity has been established, one may write
        $f:X\twoheadrightarrow Y$. The function is bijective if it is both
        injective and surjective.
    \end{definition}

    \begin{proposition}[Finite Injective-Surjective Criterion]
        If $X$ and $Y$ are finite sets with $|X|=|Y|$, then for a function
        $f:X\to Y$, the following are equivalent:
        \[
            f\text{ is injective}
            \quad\Longleftrightarrow\quad
            f\text{ is surjective}
            \quad\Longleftrightarrow\quad
            f\text{ is bijective}.
        \]
    \end{proposition}

    \begin{proposition}[Inverse Image Preserves Set Operations]
        If $f:X\to Y$ and $A,B\subseteq Y$, then
        \[
            f^{-1}(A\cup B)=f^{-1}(A)\cup f^{-1}(B),
        \]
        \[
            f^{-1}(A\cap B)=f^{-1}(A)\cap f^{-1}(B),
        \]
        and
        \[
            f^{-1}(Y\setminus A)=X\setminus f^{-1}(A).
        \]
    \end{proposition}

    \begin{definition}[Relation and Equivalence Relation]
        A binary relation on a set $A$ is a subset $R\subseteq A\times A$.
        An equivalence relation $\sim$ is reflexive, symmetric, and transitive;
        explicitly,
        \[
            \forall x\in A,
            \qquad x\sim x,
        \]
        \[
            \forall x,y\in A,
            \qquad x\sim y\implies y\sim x,
        \]
        and
        \[
            \forall x,y,z\in A,
            \qquad
            (x\sim y\land y\sim z)\implies x\sim z.
        \]
        The equivalence class of $a\in A$ is
        \[
            [a]=\{x\in A\mid x\sim a\}.
        \]
    \end{definition}

    \begin{proposition}[Equivalence Relations and Partitions]
        Equivalence relations on a set $A$ are the same structure as partitions
        of $A$: each equivalence relation partitions $A$ into equivalence
        classes, and each partition of $A$ determines an equivalence relation.
    \end{proposition}

    \begin{definition}[Countability]
        A set $A$ is countable if it is finite or if there exists a bijection
        $f:A\to\N$. In the countably infinite case,
        \[
            \exists f:A\xrightarrow{\sim}\N,
            \qquad
            |A|=|\N|=\aleph_0.
        \]
        A set is uncountable if it is not countable.
    \end{definition}

    \begin{theorem}[Cantor's Diagonal Argument]
        The set $\R$ of real numbers is uncountable. In particular,
        \[
            |\N|<|\R|.
        \]
    \end{theorem}

    \begin{theorem}[Cantor's Theorem]
        For every set $A$, there is no surjection from $A$ onto $\Pow(A)$.
        Equivalently,
        \[
            |A|<|\Pow(A)|.
        \]
    \end{theorem}

    \begin{theorem}[Schr\"oder--Bernstein Theorem; Cantor--Bernstein Theorem]
        If there exists an injection $f:A\hookrightarrow B$ and an injection
        $g:B\hookrightarrow A$, then there exists a bijection
        \[
            h:A\xrightarrow{\sim}B.
        \]
        Thus two sets have the same cardinality once each can be embedded into
        the other.
    \end{theorem}

    \begin{algorithm}[Mathematical Induction]
        To prove a statement $P(n)$ for all integers $n\geq n_0$:
        \begin{enumerate}
            \item Prove the base case $P(n_0)$.
            \item Prove the induction step: for every $k\geq n_0$, show that
                  $P(k)$ implies $P(k+1)$.
        \end{enumerate}
        Then $P(n)$ holds for all integers $n\geq n_0$.
    \end{algorithm}

    \begin{theorem}[Well-Ordering Principle]
        Every nonempty subset of $\N$ has a least element.
    \end{theorem}

    \begin{fact}[Russell's Paradox]
        The expression ``the set of all sets that do not contain themselves''
        leads to a contradiction in naive set theory. This is one reason modern
        set theory is formulated axiomatically.
    \end{fact}

    \begin{fact}[Axiom of Choice and Zorn's Lemma]
        The Axiom of Choice, Zorn's Lemma, and the Well-Ordering Theorem are
        equivalent over the usual Zermelo--Fraenkel axioms. They are standard
        foundational tools in advanced algebra, analysis, and topology.
    \end{fact}

    \chapter{Calculus and Analysis}

    \relatedcourses{(MATH101) Calculus I, (MATH102) Calculus II,
        (MATH210) Applied Complex Variables, (MATH311) Analysis I,
        (MATH312) Analysis II}

    \relatedbooks{Calculus With Applications (Peter D. Lax, Maria Shea Terrell),
        Multivariable Calculus with Applications (Peter D. Lax, Maria Shea Terrell),
        Principles of Mathematical Analysis (Walter Rudin)}

    \section*{One-Variable Calculus}

    \begin{formula}[Arithmetic and Geometric Progressions]
        For an arithmetic progression with first term $a$ and common difference
        $d$, we have
        \[
            a_n=a+(n-1)d,
            \qquad
            \sum_{k=1}^{n}a_k=\frac{n}{2}(a_1+a_n).
        \]
        For a geometric progression with first term $a$ and common ratio $r$, we
        have
        \[
            a_n=ar^{n-1},
            \qquad
            \sum_{k=0}^{n-1}ar^k=a\frac{1-r^n}{1-r}\quad(r\neq 1).
        \]
        If $|r|<1$, then
        \[
            \sum_{k=0}^{\infty}ar^k=\frac{a}{1-r}.
        \]
    \end{formula}

    \begin{formula}[Sums of Powers; Faulhaber's Formulas]
        The first power sums are
        \[
            \sum_{k=1}^{n}k
            =
            \frac{n(n+1)}{2},
        \]
        \[
            \sum_{k=1}^{n}k^2
            =
            \frac{n(n+1)(2n+1)}{6},
            \qquad
            \sum_{k=1}^{n}k^3
            =
            \left(\frac{n(n+1)}{2}\right)^2,
        \]
        and
        \[
            \sum_{k=1}^{n}k^4
            =
            \frac{n(n+1)(2n+1)(3n^2+3n-1)}{30}.
        \]
        More generally, for each fixed nonnegative integer $m$,
        \[
            \sum_{k=1}^{n}k^m
        \]
        is a polynomial in $n$ of degree $m+1$ with leading coefficient
        $1/(m+1)$.
    \end{formula}

    \begin{formula}[Finite Differences; Newton's Forward Difference Formula]
        For a sequence $(a_n)$, define the forward difference by
        \[
            \Delta a_n=a_{n+1}-a_n.
        \]
        If $a_n=P(n)$ for a polynomial $P$ of degree $d$ with leading
        coefficient $c$, then
        \[
            \Delta^d a_n=d!c,
            \qquad
            \Delta^{d+1}a_n=0.
        \]
        Conversely, a sequence with constant $d$th differences is polynomial
        in $n$ of degree at most $d$. For every polynomial $P$ of degree at most
        $d$,
        \[
            P(n)
            =
            \sum_{k=0}^{d}\binom{n}{k}\Delta^kP(0).
        \]
    \end{formula}

    \begin{strategy}
        A table of successive differences is one of the fastest ways to detect a
        hidden polynomial sequence. Constant first, second, and third differences
        indicate linear, quadratic, and cubic behavior, respectively.
    \end{strategy}

    \begin{formula}[Linear Recurrences with Constant Coefficients]
        Consider
        \[
            a_n=c_1a_{n-1}+c_2a_{n-2}+\cdots+c_ma_{n-m}.
        \]
        Its characteristic polynomial is
        \[
            r^m-c_1r^{m-1}-c_2r^{m-2}-\cdots-c_m.
        \]
        If the roots $r_1,\dots,r_m$ are distinct, then
        \[
            a_n=A_1r_1^n+\cdots+A_mr_m^n.
        \]
        A root $r$ of multiplicity $q$ contributes terms
        \[
            r^n,\ nr^n,\ \dots,\ n^{q-1}r^n.
        \]
        The constants are determined from the initial values.
    \end{formula}

    \begin{formula}[First-Order Linear Recurrence]
        If
        \[
            a_n=ra_{n-1}+b_n
            \qquad(n\geq 1),
        \]
        then iteration gives
        \[
            a_n
            =
            r^na_0
            +
            \sum_{j=1}^{n}r^{\,n-j}b_j.
        \]
        In particular, if $b_n=b$ and $r\neq 1$, then
        \[
            a_n
            =
            r^na_0+b\frac{r^n-1}{r-1}.
        \]
    \end{formula}

    \begin{formula}[Nonhomogeneous Linear Recurrences]
        For a linear recurrence with constant coefficients, the general solution
        is
        \[
            a_n=a_n^{(h)}+a_n^{(p)},
        \]
        where $a_n^{(h)}$ solves the associated homogeneous recurrence and
        $a_n^{(p)}$ is one particular solution. If the forcing term has the form
        \[
            P(n)\alpha^n,
        \]
        try a particular solution of the form
        \[
            n^sQ(n)\alpha^n,
        \]
        where $\deg Q=\deg P$ and $s$ is the multiplicity of $\alpha$ as a root
        of the characteristic polynomial.
    \end{formula}

    \begin{strategy}
        For recurrences, first identify whether the quickest language is
        characteristic roots, generating functions, a transfer matrix, or
        telescoping after a suitable normalization.
    \end{strategy}

    \begin{definition}[Limit of a Function; Epsilon--Delta Definition]
        Let $D\subseteq\R$, let $a$ be an accumulation point of $D$, and let
        $f:D\to\R$. The statement
        \[
            \lim_{x\to a}f(x)=L
        \]
        means that $f(x)$ can be made arbitrarily close to $L$ by taking $x$
        sufficiently close to, but different from, $a$. Formally,
        \[
            \forall\eps>0,
            \quad
            \exists\delta>0,
            \quad
            \forall x\in D,
            \qquad
            0<|x-a|<\delta
            \implies
            |f(x)-L|<\eps.
        \]
    \end{definition}

    \begin{definition}[One-Sided Limits and Limits at Infinity]
        The right-hand limit $\lim_{x\to a^+}f(x)=L$ is defined by
        \[
            \forall\eps>0,
            \quad
            \exists\delta>0,
            \quad
            \forall x\in D,
            \qquad
            0<x-a<\delta
            \implies
            |f(x)-L|<\eps.
        \]
        The left-hand definition is analogous, with $0<a-x<\delta$.
        Moreover,
        \[
            \lim_{x\to\infty}f(x)=L
        \]
        means
        \[
            \forall\eps>0,
            \quad
            \exists M>0,
            \quad
            \forall x\in D,
            \qquad
            x>M\implies |f(x)-L|<\eps.
        \]
        Finally, $\lim_{x\to a}f(x)=\infty$ means
        \[
            \forall M>0,
            \quad
            \exists\delta>0,
            \quad
            \forall x\in D,
            \qquad
            0<|x-a|<\delta\implies f(x)>M.
        \]
    \end{definition}

    \begin{theorem}[Squeeze Theorem; Sandwich Theorem]
        Suppose that $g(x)\leq f(x)\leq h(x)$ for all $x$ sufficiently close
        to $a$, except possibly at $a$, and that
        \[
            \lim_{x\to a}g(x)
            =
            \lim_{x\to a}h(x)
            =L.
        \]
        Then
        \[
            \lim_{x\to a}f(x)=L.
        \]
    \end{theorem}

    \begin{definition}[Continuity at a Point]
        Let $f:D\to\R$ and $a\in D$. The function $f$ is continuous at $a$ if
        \[
            \lim_{x\to a}f(x)=f(a).
        \]
        Equivalently,
        \[
            \forall\eps>0,
            \quad
            \exists\delta>0,
            \quad
            \forall x\in D,
            \qquad
            |x-a|<\delta
            \implies
            |f(x)-f(a)|<\eps.
        \]
        The function is continuous on $D$ if it is continuous at every point of
        $D$.
    \end{definition}

    \begin{definition}[Derivative]
        Let $f$ be defined near $a$. The derivative of $f$ at $a$ is
        \[
            f'(a)
            =
            \lim_{h\to 0}\frac{f(a+h)-f(a)}{h},
        \]
        provided this limit exists as a finite real number. Thus $f'(a)=L$ if
        and only if
        \[
            \forall\eps>0,
            \quad
            \exists\delta>0,
            \quad
            \forall h\in\R,
            \qquad
            0<|h|<\delta
            \implies
            \left|
                \frac{f(a+h)-f(a)}{h}-L
            \right|<\eps.
        \]
    \end{definition}

    \begin{formula}[Standard Trigonometric and Exponential Limits]
        The following limits are fundamental:
        \[
            \lim_{x\to 0}\frac{\sin x}{x}=1,
            \qquad
            \lim_{x\to 0}\frac{1-\cos x}{x^2}=\frac12,
        \]
        \[
            \lim_{x\to 0}\frac{e^x-1}{x}=1,
            \qquad
            \lim_{x\to 0}\frac{\ln(1+x)}{x}=1,
        \]
        and
        \[
            \lim_{x\to 0}(1+x)^{1/x}=e.
        \]
    \end{formula}

    \begin{strategy}
        These limits are often faster than L'H\^opital's Rule. The
        expressions $\sin x$, $1-\cos x$, $e^x-1$, $\ln(1+x)$, and
        $(1+x)^{1/x}$ are strong signals for these standard patterns.
    \end{strategy}

    \begin{formula}[Trigonometric Identity Table]
        The principal trigonometric identities are

        {\small
        \renewcommand{\arraystretch}{1.18}
        \begin{longtable}{@{}L{0.25\textwidth}L{0.67\textwidth}@{}}
            \toprule
            Type & Identities \\
            \midrule
            \endfirsthead

            \toprule
            Type & Identities \\
            \midrule
            \endhead

            \bottomrule
            \endlastfoot

            Reciprocal &
            $\displaystyle \sec x=\frac{1}{\cos x},\quad
            \csc x=\frac{1}{\sin x},\quad
            \cot x=\frac{1}{\tan x}$ \\
            Quotient &
            $\displaystyle \tan x=\frac{\sin x}{\cos x},\quad
            \cot x=\frac{\cos x}{\sin x}$ \\
            Pythagorean &
            $\sin^2x+\cos^2x=1$, $1+\tan^2x=\sec^2x$,
            $1+\cot^2x=\csc^2x$ \\
            Angle addition &
            $\sin(x\pm y)=\sin x\cos y\pm\cos x\sin y$ \\
            Angle addition &
            $\cos(x\pm y)=\cos x\cos y\mp\sin x\sin y$ \\
            Angle addition &
            $\displaystyle\tan(x\pm y)=
            \frac{\tan x\pm\tan y}{1\mp\tan x\tan y}$ \\
            Double angle &
            $\sin 2x=2\sin x\cos x$ \\
            Double angle &
            $\cos 2x=\cos^2x-\sin^2x=2\cos^2x-1=1-2\sin^2x$ \\
            Double angle &
            $\displaystyle\tan 2x=\frac{2\tan x}{1-\tan^2x}$ \\
            Half angle &
            $\displaystyle\sin^2\frac{x}{2}=\frac{1-\cos x}{2},\quad
            \cos^2\frac{x}{2}=\frac{1+\cos x}{2}$ \\
            Product-to-sum &
            $\displaystyle2\sin x\sin y=\cos(x-y)-\cos(x+y)$ \\
            Product-to-sum &
            $\displaystyle2\cos x\cos y=\cos(x-y)+\cos(x+y)$ \\
            Product-to-sum &
            $\displaystyle2\sin x\cos y=\sin(x+y)+\sin(x-y)$ \\
            Sum-to-product &
            $\displaystyle\sin x+\sin y=
            2\sin\frac{x+y}{2}\cos\frac{x-y}{2}$ \\
            Sum-to-product &
            $\displaystyle\cos x+\cos y=
            2\cos\frac{x+y}{2}\cos\frac{x-y}{2}$ \\
        \end{longtable}
        }
    \end{formula}

    \begin{theorem}[L'H\^opital's Rule]
        Let $f$ and $g$ be differentiable on a deleted neighborhood of $a$, and
        suppose that $g'(x)\neq 0$ there. Assume either
        \[
            f(x)\to 0,
            \qquad
            g(x)\to 0,
        \]
        or
        \[
            |f(x)|\to\infty,
            \qquad
            |g(x)|\to\infty,
        \]
        as $x\to a$. If
        \[
            \lim_{x\to a}\frac{f'(x)}{g'(x)}=L
        \]
        exists in the extended real numbers, then
        \[
            \lim_{x\to a}\frac{f(x)}{g(x)}=L.
        \]
    \end{theorem}

    \begin{strategy}
        Apply L'H\^opital's Rule only after identifying a $0/0$ or
        $\infty/\infty$ form. Rewrite $0\cdot\infty$ and $\infty-\infty$ as a
        quotient first. For $0^0$, $1^\infty$, or $\infty^0$, take logarithms
        before differentiating.
    \end{strategy}

    \begin{formula}[Derivative and Antiderivative Table]
        The following table records the basic one-variable forms. Affine changes
        of variable are handled by the chain rule and substitution. The constant
        of integration is omitted from the right column.

        {\small
        \renewcommand{\arraystretch}{1.18}
        \begin{longtable}{@{}L{0.26\textwidth}L{0.29\textwidth}L{0.35\textwidth}@{}}
            \toprule
            $f(x)$ & $f'(x)$ & $\displaystyle\int f(x)\dd x$ \\
            \midrule
            \endfirsthead

            \toprule
            $f(x)$ & $f'(x)$ & $\displaystyle\int f(x)\dd x$ \\
            \midrule
            \endhead

            \bottomrule
            \endlastfoot

            $c$ & $0$ & $cx$ \\
            $x$ & $1$ & $\displaystyle\frac{x^2}{2}$ \\
            $x^a$ & $ax^{a-1}$ & $\displaystyle\frac{x^{a+1}}{a+1}\quad(a\neq -1)$ \\
            $\displaystyle\frac{1}{x}$ & $\displaystyle -\frac{1}{x^2}$ & $\ln|x|$ \\
            $e^x$ & $e^x$ & $e^x$ \\
            $a^x$ & $a^x\ln a$ & $\displaystyle\frac{a^x}{\ln a}\quad(a>0,\ a\neq 1)$ \\
            $\ln x$ & $\displaystyle\frac{1}{x}$ & $x\ln x-x$ \\
            $\lg x$ & $\displaystyle\frac{1}{x\ln 10}$ & $\displaystyle\frac{x\ln x-x}{\ln 10}$ \\
            $\lb x$ & $\displaystyle\frac{1}{x\ln 2}$ & $\displaystyle\frac{x\ln x-x}{\ln 2}$ \\
            $\log_a x$ & $\displaystyle\frac{1}{x\ln a}$ & $\displaystyle\frac{x\ln x-x}{\ln a}\quad(a>0,\ a\neq 1)$ \\
            $x\ln x$ & $\ln x+1$ & $\displaystyle\frac{x^2}{2}\ln x-\frac{x^2}{4}$ \\
            $\displaystyle\frac{\ln x}{x}$ & $\displaystyle\frac{1-\ln x}{x^2}$ & $\displaystyle\frac{1}{2}(\ln x)^2$ \\
            $\displaystyle\frac{1}{x\ln x}$ & $\displaystyle-\frac{\ln x+1}{x^2(\ln x)^2}$ & $\ln|\ln x|$ \\
            $\sin x$ & $\cos x$ & $-\cos x$ \\
            $\cos x$ & $-\sin x$ & $\sin x$ \\
            $\tan x$ & $\sec^2x$ & $-\ln|\cos x|$ \\
            $\cot x$ & $-\csc^2x$ & $\ln|\sin x|$ \\
            $\sec x$ & $\sec x\tan x$ & $\ln|\sec x+\tan x|$ \\
            $\csc x$ & $-\csc x\cot x$ & $\ln|\csc x-\cot x|$ \\
            $\sec^2x$ & $2\sec^2x\tan x$ & $\tan x$ \\
            $\csc^2x$ & $-2\csc^2x\cot x$ & $-\cot x$ \\
            $\sec x\tan x$ & $\sec x\tan^2x+\sec^3x$ & $\sec x$ \\
            $\csc x\cot x$ & $-\csc x\cot^2x-\csc^3x$ & $-\csc x$ \\
            $\arcsin x$ & $\displaystyle\frac{1}{\sqrt{1-x^2}}$ & $x\arcsin x+\sqrt{1-x^2}$ \\
            $\arccos x$ & $\displaystyle-\frac{1}{\sqrt{1-x^2}}$ & $x\arccos x-\sqrt{1-x^2}$ \\
            $\arctan x$ & $\displaystyle\frac{1}{1+x^2}$ & $x\arctan x-\displaystyle\frac{1}{2}\ln(1+x^2)$ \\
            $\operatorname{arccot}x$ & $\displaystyle-\frac{1}{1+x^2}$ & $x\operatorname{arccot}x+\displaystyle\frac{1}{2}\ln(1+x^2)$ \\
            $\operatorname{arcsec}x$ & $\displaystyle\frac{1}{|x|\sqrt{x^2-1}}$ & $x\operatorname{arcsec}x-\ln\left|x+\sqrt{x^2-1}\right|$ \\
            $\operatorname{arccsc}x$ & $\displaystyle-\frac{1}{|x|\sqrt{x^2-1}}$ & $x\operatorname{arccsc}x+\ln\left|x+\sqrt{x^2-1}\right|$ \\
            $\sinh x$ & $\cosh x$ & $\cosh x$ \\
            $\cosh x$ & $\sinh x$ & $\sinh x$ \\
            $\tanh x$ & $\operatorname{sech}^2x$ & $\ln(\cosh x)$ \\
            $\coth x$ & $-\operatorname{csch}^2x$ & $\ln|\sinh x|$ \\
            $\operatorname{sech}x$ & $-\operatorname{sech}x\tanh x$ & $\arctan(\sinh x)$ \\
            $\operatorname{csch}x$ & $-\operatorname{csch}x\coth x$ & $\ln\left|\tanh\frac{x}{2}\right|$ \\
            $\operatorname{arsinh}x$ & $\displaystyle\frac{1}{\sqrt{x^2+1}}$ & $x\operatorname{arsinh}x-\sqrt{x^2+1}$ \\
            $\operatorname{arcosh}x$ & $\displaystyle\frac{1}{\sqrt{x-1}\sqrt{x+1}}$ & $x\operatorname{arcosh}x-\sqrt{x^2-1}$ \\
            $\operatorname{artanh}x$ & $\displaystyle\frac{1}{1-x^2}$ & $x\operatorname{artanh}x+\displaystyle\frac{1}{2}\ln(1-x^2)$ \\
            $\displaystyle\frac{1}{x^2+a^2}$ & $\displaystyle-\frac{2x}{(x^2+a^2)^2}$ & $\displaystyle\frac{1}{a}\arctan\frac{x}{a}\quad(a>0)$ \\
            $\displaystyle\frac{1}{x^2-a^2}$ & $\displaystyle-\frac{2x}{(x^2-a^2)^2}$ & $\displaystyle\frac{1}{2a}\ln\left|\frac{x-a}{x+a}\right|\quad(a>0)$ \\
            $\displaystyle\frac{1}{a^2-x^2}$ & $\displaystyle\frac{2x}{(a^2-x^2)^2}$ & $\displaystyle\frac{1}{2a}\ln\left|\frac{a+x}{a-x}\right|\quad(a>0)$ \\
            $\displaystyle\frac{1}{\sqrt{a^2-x^2}}$ & $\displaystyle\frac{x}{(a^2-x^2)^{3/2}}$ & $\displaystyle\arcsin\frac{x}{a}\quad(a>0)$ \\
            $\displaystyle\frac{1}{\sqrt{x^2+a^2}}$ & $\displaystyle-\frac{x}{(x^2+a^2)^{3/2}}$ & $\ln\left|x+\sqrt{x^2+a^2}\right|$ \\
            $\displaystyle\frac{1}{\sqrt{x^2-a^2}}$ & $\displaystyle-\frac{x}{(x^2-a^2)^{3/2}}$ & $\ln\left|x+\sqrt{x^2-a^2}\right|\quad(a>0)$ \\
            $\displaystyle\frac{x}{\sqrt{x^2+a^2}}$ & $\displaystyle\frac{a^2}{(x^2+a^2)^{3/2}}$ & $\sqrt{x^2+a^2}$ \\
            $\displaystyle\frac{x}{\sqrt{a^2-x^2}}$ & $\displaystyle\frac{a^2}{(a^2-x^2)^{3/2}}$ & $-\sqrt{a^2-x^2}$ \\
            $\sqrt{a^2-x^2}$ & $\displaystyle-\frac{x}{\sqrt{a^2-x^2}}$ & $\displaystyle\frac{x}{2}\sqrt{a^2-x^2}+\frac{a^2}{2}\arcsin\frac{x}{a}$ \\
            $\sqrt{x^2+a^2}$ & $\displaystyle\frac{x}{\sqrt{x^2+a^2}}$ & $\displaystyle\frac{x}{2}\sqrt{x^2+a^2}+\frac{a^2}{2}\ln\left|x+\sqrt{x^2+a^2}\right|$ \\
            $\sqrt{x^2-a^2}$ & $\displaystyle\frac{x}{\sqrt{x^2-a^2}}$ & $\displaystyle\frac{x}{2}\sqrt{x^2-a^2}-\frac{a^2}{2}\ln\left|x+\sqrt{x^2-a^2}\right|$ \\
        \end{longtable}
        }
    \end{formula}

    \begin{tactic}
        An antiderivative is often exposed by a simple substitution, symmetry,
        completing the square, or integration by parts. Try those before expanding into a long calculation.
    \end{tactic}

    \begin{theorem}[Intermediate Value Theorem]
        Let $f$ be continuous on $[a,b]$. If $N$ is any number between
        $f(a)$ and $f(b)$, then there exists $c\in[a,b]$ such that
        \[
            f(c)=N.
        \]
    \end{theorem}

    \begin{strategy}
        The Intermediate Value Theorem is the theorem behind many existence
        arguments. To show that an equation has a solution, define a continuous
        function and show that its values cross the desired level.
    \end{strategy}

    \begin{theorem}[Extreme Value Theorem; Weierstrass Extreme Value Theorem]
        Let $f$ be continuous on $[a,b]$. Then $f$ attains both a maximum and
        a minimum on $[a,b]$.
    \end{theorem}

    \begin{remark}
        The closed interval condition is important. A continuous function on an
        open interval need not attain a maximum or a minimum.
    \end{remark}

    \begin{theorem}[Rolle's Theorem]
        Let $f$ be continuous on $[a,b]$ and differentiable on $(a,b)$. If
        $f(a)=f(b)$, then there exists $c\in(a,b)$ such that
        \[
            f'(c)=0.
        \]
    \end{theorem}

    \begin{theorem}[Mean Value Theorem; Lagrange Mean Value Theorem]
        Let $f$ be continuous on $[a,b]$ and differentiable on $(a,b)$. Then
        there exists $c\in(a,b)$ such that
        \[
            f'(c)
            =
            \frac{f(b)-f(a)}{b-a}.
        \]
    \end{theorem}

    \begin{idea}
        The Mean Value Theorem says that somewhere on the interval, the
        instantaneous rate of change equals the average rate of change.
    \end{idea}

    \begin{theorem}[Cauchy's Mean Value Theorem; Extended Mean Value Theorem]
        Let $f$ and $g$ be continuous on $[a,b]$ and differentiable on
        $(a,b)$. Then there exists $c\in(a,b)$ such that
        \[
            \bigl(f(b)-f(a)\bigr)g'(c)
            =
            \bigl(g(b)-g(a)\bigr)f'(c).
        \]
        If $g'(x)\neq 0$ for every $x\in(a,b)$, then $g(b)\neq g(a)$
        and the result may be written as
        \[
            \frac{f'(c)}{g'(c)}
            =
            \frac{f(b)-f(a)}{g(b)-g(a)}.
        \]
    \end{theorem}

    \begin{theorem}[Derivative of an Inverse Function]
        Let $f$ be one-to-one and differentiable near $x_0$, with
        $f'(x_0)\neq 0$. If $g=f^{-1}$ and $y_0=f(x_0)$, then
        \[
            g'(y_0)
            =
            \frac{1}{f'(x_0)}.
        \]
    \end{theorem}

    \begin{theorem}[Fundamental Theorem of Calculus]
        If $f$ is continuous on $[a,b]$ and
        \[
            F(x)=\int_a^x f(t)\dd t,
        \]
        then $F'(x)=f(x)$. Conversely, if $F'=f$ on $[a,b]$, then
        \[
            \int_a^b f(x)\dd x=F(b)-F(a).
        \]
    \end{theorem}

    \begin{theorem}[Taylor's Theorem]
        Suppose that $f$ is $n+1$ times differentiable near $a$. Then
        \[
            f(x)
            =
            f(a)+f'(a)(x-a)+\frac{f''(a)}{2!}(x-a)^2
            +\cdots
            +\frac{f^{(n)}(a)}{n!}(x-a)^n
            +R_n(x),
        \]
        where, in Lagrange's form of the remainder,
        \[
            R_n(x)
            =
            \frac{f^{(n+1)}(c)}{(n+1)!}(x-a)^{n+1}
        \]
        for some $c$ between $a$ and $x$.
    \end{theorem}

    \begin{corollary}[Taylor Error Bound]
        If
        \[
            |f^{(n+1)}(t)|\leq M
        \]
        for every $t$ between $a$ and $x$, then the Taylor remainder satisfies
        \[
            |R_n(x)|
            \leq
            \frac{M}{(n+1)!}|x-a|^{n+1}.
        \]
    \end{corollary}

    \begin{strategy}
        Taylor expansion is one of the fastest tools for limits, approximations,
        and hidden series evaluations. Expressions involving $e^x$, $\sin x$,
        $\cos x$, $\ln(1+x)$, or $\arctan x$ often invite Taylor expansion.
    \end{strategy}

    \begin{formula}[Maclaurin Series]
        The most important Maclaurin series are
        \[
            e^x
            =
            \sum_{n=0}^{\infty}\frac{x^n}{n!},
        \]
        \[
            \sin x
            =
            \sum_{n=0}^{\infty}(-1)^n\frac{x^{2n+1}}{(2n+1)!},
            \qquad
            \cos x
            =
            \sum_{n=0}^{\infty}(-1)^n\frac{x^{2n}}{(2n)!},
        \]
        \[
            \frac{1}{1-x}
            =
            \sum_{n=0}^{\infty}x^n
            \qquad (|x|<1),
        \]
        \[
            \ln(1+x)
            =
            \sum_{n=1}^{\infty}(-1)^{n+1}\frac{x^n}{n}
            \qquad (-1<x\leq 1),
        \]
        and
        \[
            \arctan x
            =
            \sum_{n=0}^{\infty}(-1)^n\frac{x^{2n+1}}{2n+1}
            \qquad (|x|\leq 1).
        \]
        Also,
        \[
            -\ln(1-x)
            =
            \sum_{n=1}^{\infty}\frac{x^n}{n}
            \qquad (-1\leq x<1),
        \]
        \[
            \sinh x
            =
            \sum_{n=0}^{\infty}\frac{x^{2n+1}}{(2n+1)!},
            \qquad
            \cosh x
            =
            \sum_{n=0}^{\infty}\frac{x^{2n}}{(2n)!},
        \]
        and
        \[
            \operatorname{artanh}x
            =
            \sum_{n=0}^{\infty}\frac{x^{2n+1}}{2n+1}
            \qquad (|x|<1).
        \]
    \end{formula}

    \begin{formula}[Binomial Series]
        For any real or complex number $\beta$,
        \[
            (1+x)^\beta
            =
            \sum_{n=0}^{\infty}\binom{\beta}{n}x^n
            \qquad (|x|<1),
        \]
        where
        \[
            \binom{\beta}{n}
            =
            \frac{\beta(\beta-1)\cdots(\beta-n+1)}{n!}.
        \]
    \end{formula}

    \begin{formula}[Chebyshev Polynomials]
        The Chebyshev polynomials of the first kind are characterized by
        \[
            T_n(\cos\theta)=\cos(n\theta).
        \]
        They satisfy
        \[
            T_0(x)=1,
            \qquad
            T_1(x)=x,
            \qquad
            T_{n+1}(x)=2xT_n(x)-T_{n-1}(x).
        \]
        Hence multiple-angle trigonometric expressions can be converted into
        polynomial calculations almost immediately; for example,
        \[
            \cos(3\theta)=4\cos^3\theta-3\cos\theta.
        \]
    \end{formula}

    \begin{formula}[Classical Series Values]
        The following standard series values are useful in computation:
        \[
            \sum_{n=1}^{\infty}\frac{1}{n^2}
            =
            \frac{\pi^2}{6},
            \qquad
            \sum_{n=1}^{\infty}\frac{1}{n^4}
            =
            \frac{\pi^4}{90},
        \]
        \[
            \sum_{n=1}^{\infty}\frac{(-1)^{n+1}}{n}
            =
            \ln 2,
            \qquad
            \sum_{n=0}^{\infty}\frac{(-1)^n}{2n+1}
            =
            \frac{\pi}{4}.
        \]
        More generally, the $p$-series
        \[
            \sum_{n=1}^{\infty}\frac{1}{n^p}
        \]
        converges if and only if $p>1$.
    \end{formula}

    \begin{formula}[Geometric-Series Differentiation]
        Starting from
        \[
            \frac{1}{1-x}
            =
            \sum_{n=0}^{\infty}x^n
            \qquad (|x|<1),
        \]
        term-by-term differentiation gives
        \[
            \frac{1}{(1-x)^2}
            =
            \sum_{n=1}^{\infty}nx^{n-1}
            \qquad (|x|<1),
        \]
        and hence
        \[
            \sum_{n=1}^{\infty}nx^n
            =
            \frac{x}{(1-x)^2}
            \qquad (|x|<1).
        \]
    \end{formula}

    \begin{formula}[Telescoping Series]
        If
        \[
            a_n=b_n-b_{n+1},
        \]
        then the partial sums satisfy
        \[
            \sum_{n=1}^{N}a_n
            =
            b_1-b_{N+1}.
        \]
        Hence, if $b_n\to L$, then
        \[
            \sum_{n=1}^{\infty}a_n=b_1-L.
        \]
    \end{formula}

    \begin{formula}[Abel's Summation Formula; Summation by Parts]
        Let
        \[
            A_k=\sum_{j=m}^{k}a_j.
        \]
        Then
        \[
            \sum_{k=m}^{n}a_kb_k
            =
            A_nb_n
            +
            \sum_{k=m}^{n-1}A_k(b_k-b_{k+1}).
        \]
        This is the discrete analogue of integration by parts and is useful when
        one factor has simple partial sums while the other changes slowly.
    \end{formula}

    \begin{formula}[Euler--Maclaurin Summation Formula]
        For a sufficiently smooth function $f$,
        \[
            \sum_{k=m}^{n}f(k)
            =
            \int_m^n f(x)\dd x
            +
            \frac{f(m)+f(n)}{2}
            +
            \sum_{r=1}^{p}
            \frac{B_{2r}}{(2r)!}
            \left(
                f^{(2r-1)}(n)-f^{(2r-1)}(m)
            \right)
            +
            R_p,
        \]
        where $B_{2r}$ are Bernoulli numbers. The first correction
        \[
            \sum_{k=m}^{n}f(k)
            \approx
            \int_m^n f(x)\dd x+\frac{f(m)+f(n)}{2}
        \]
        is often enough for a rapid asymptotic estimate.
    \end{formula}

    \begin{formula}[Asymptotic Growth Hierarchy]
        For fixed $p,q>0$ and $c>1$, as $n\to\infty$,
        \[
            (\ln n)^p=o(n^q),
            \qquad
            n^q=o(c^n),
            \qquad
            c^n=o(n!),
            \qquad
            n!=o(n^n).
        \]
    \end{formula}

    \begin{strategy}
        For fast series calculations, first try to reduce the expression to a
        geometric series, a differentiated geometric series, a Maclaurin series,
        or a known special value such as $\pi^2/6$, $\pi^4/90$, $\ln 2$, or
        $\pi/4$.
    \end{strategy}

    \begin{formula}[Integration by Parts]
        If $u$ and $v$ are differentiable functions, then
        \[
            \int u\dd v
            =
            uv-\int v\dd u.
        \]
        Equivalently,
        \[
            \int_a^b u(x)v'(x)\dd x
            =
            \left[u(x)v(x)\right]_a^b
            -
            \int_a^b u'(x)v(x)\dd x.
        \]
    \end{formula}

    \begin{theorem}[Leibniz Integral Rule; Differentiation under the Integral Sign]
        Under suitable continuity and domination assumptions,
        \[
            \diff{}{t}\int_{a(t)}^{b(t)} f(x,t)\dd x
            =
            f(b(t),t)b'(t)-f(a(t),t)a'(t)
            +
            \int_{a(t)}^{b(t)}\frac{\partial f}{\partial t}(x,t)\dd x.
        \]
        When the endpoints are constant, only the integral term remains.
    \end{theorem}

    \begin{tactic}
        Introducing a parameter and differentiating under the integral sign is
        often called Feynman's trick. It is especially effective when direct
        integration is difficult but the parameter derivative is elementary.
    \end{tactic}

    \begin{formula}[Frullani's Integral]
        If $f$ is continuous on $(0,\infty)$ and the endpoint limits exist under
        conditions ensuring convergence, then for $a,b>0$,
        \[
            \int_0^\infty \frac{f(ax)-f(bx)}{x}\dd x
            =
            \bigl(f(0^+)-f(\infty)\bigr)\ln\frac{b}{a}.
        \]
        A standard special case is
        \[
            \int_0^\infty \frac{e^{-ax}-e^{-bx}}{x}\dd x
            =
            \ln\frac{b}{a}.
        \]
    \end{formula}

    \begin{formula}[Arc Length Formula]
        The length of the graph $y=f(x)$ on $[a,b]$ is
        \[
            L
            =
            \int_a^b \sqrt{1+(f'(x))^2}\dd x.
        \]
        More generally, if a curve is parametrized by
        $\mathbf{r}(t)=(x(t),y(t),z(t))$, then
        \[
            L
            =
            \int_a^b \|\mathbf{r}'(t)\|\dd t.
        \]
    \end{formula}

    \begin{formula}[Vi\`ete's Infinite Product]
        For real $x\neq 0$, we have
        \[
            \frac{\sin x}{x}
            =
            \prod_{k=1}^{\infty}
            \cos\left(\frac{x}{2^k}\right).
        \]
        In particular,
        \[
            \frac{2}{\pi}
            =
            \cos\frac{\pi}{4}
            \cdot
            \cos\frac{\pi}{8}
            \cdot
            \cos\frac{\pi}{16}
            \cdots.
        \]
    \end{formula}

    \begin{idea}
        Vi\`ete's formula is a product formula for $\pi$. The key identity is
        the double-angle formula
        \[
            \sin(2x)=2\sin x\cos x.
        \]
    \end{idea}

    \begin{formula}[Euler's Formula]
        For every real number $x$,
        \[
            e^{ix}
            =
            \cos x+i\sin x.
        \]
        In particular,
        \[
            e^{i\pi}+1=0.
        \]
    \end{formula}

    \begin{remark}
        The identity $e^{i\pi}+1=0$ is called Euler's identity. It connects
        the constants $e$, $i$, $\pi$, $1$, and $0$ in a single formula.
    \end{remark}

    \begin{formula}[Gaussian Integral]
        The Gaussian integral is
        \[
            \int_{-\infty}^{\infty}e^{-x^2}\dd x
            =
            \sqrt{\pi}.
        \]
        Equivalently,
        \[
            \int_0^\infty e^{-x^2}\dd x
            =
            \frac{\sqrt{\pi}}{2}.
        \]
    \end{formula}

    \begin{remark}
        The Gaussian integral is one of the most famous integrals in calculus.
        It is closely connected to the normal distribution in probability.
    \end{remark}

    \begin{formula}[Fresnel Integrals]
        The Fresnel integrals satisfy
        \[
            \int_0^{\infty}\cos(x^2)\dd x
            =
            \int_0^{\infty}\sin(x^2)\dd x
            =
            \sqrt{\frac{\pi}{8}}.
        \]
    \end{formula}

    \begin{remark}
        Fresnel integrals appear in wave optics, diffraction, and the Cornu
        spiral. They are memorable examples of non-elementary definite
        integrals with clean values.
    \end{remark}

    \begin{formula}[Sine Integral; Dirichlet Integral]
        The sine integral is
        \[
            \operatorname{Si}(x)
            =
            \int_0^x \frac{\sin t}{t}\dd t.
        \]
        The classical Dirichlet integral is
        \[
            \int_0^{\infty}\frac{\sin x}{x}\dd x
            =
            \frac{\pi}{2}.
        \]
    \end{formula}

    \begin{formula}[Arctangent Integral]
        For $a>0$,
        \[
            \int_{-\infty}^{\infty}\frac{1}{x^2+a^2}\dd x
            =
            \frac{\pi}{a},
            \qquad
            \int_0^{\infty}\frac{1}{x^2+a^2}\dd x
            =
            \frac{\pi}{2a}.
        \]
        More generally,
        \[
            \int_{-\infty}^{\infty}\frac{c}{(x-b)^2+a^2}\dd x
            =
            \frac{c\pi}{a}.
        \]
    \end{formula}

    \begin{formula}[Gamma-Type Exponential Integrals]
        For $a>0$ and every nonnegative integer $n$,
        \[
            \int_0^{\infty}x^n e^{-ax}\dd x
            =
            \frac{n!}{a^{n+1}}.
        \]
        More generally, for $s>0$,
        \[
            \int_0^{\infty}x^{s-1}e^{-ax}\dd x
            =
            \frac{\Gamma(s)}{a^s}.
        \]
    \end{formula}

    \begin{formula}[Damped Sine and Cosine Integrals]
        For $a>0$,
        \[
            \int_0^{\infty}e^{-ax}\cos bx\dd x
            =
            \frac{a}{a^2+b^2},
            \qquad
            \int_0^{\infty}e^{-ax}\sin bx\dd x
            =
            \frac{b}{a^2+b^2}.
        \]
    \end{formula}

    \begin{strategy}
        Many fast integral problems reduce to this formula after completing a
        square or translating the variable. A shift such as $u=x-b$ does not
        change the infinite limits.
    \end{strategy}

    \begin{formula}[Symmetry Trick for Definite Integrals]
        If $f$ is integrable on $[a,b]$, then
        \[
            \int_a^b f(x)\dd x
            =
            \int_a^b f(a+b-x)\dd x.
        \]
        Hence, if
        \[
            f(x)+f(a+b-x)=C
        \]
        is constant, then
        \[
            \int_a^b f(x)\dd x
            =
            \frac{C(b-a)}{2}.
        \]
    \end{formula}

    \begin{strategy}
        This is one of the most useful Integration-Bee-style tricks. Typical
        signals are intervals such as $[0,\pi/2]$, expressions involving
        $\tan x$ and $\cot x$, or a denominator that becomes complementary after
        the substitution $x\mapsto a+b-x$.
    \end{strategy}

    \begin{formula}[Parity and Periodicity Shortcuts for Integrals]
        If $f$ is integrable on $[-a,a]$, then
        \[
            f(-x)=-f(x)
            \quad\Longrightarrow\quad
            \int_{-a}^{a}f(x)\dd x=0,
        \]
        and
        \[
            f(-x)=f(x)
            \quad\Longrightarrow\quad
            \int_{-a}^{a}f(x)\dd x
            =
            2\int_{0}^{a}f(x)\dd x.
        \]
        If $f$ has period $T$, then for every $c$,
        \[
            \int_c^{c+T}f(x)\dd x
            =
            \int_0^T f(x)\dd x.
        \]
    \end{formula}

    \begin{formula}[Reciprocal-Substitution Identity]
        If the improper integrals converge, the substitution $x=1/t$ gives
        \[
            \int_0^\infty f(x)\dd x
            =
            \int_0^\infty \frac{f(1/x)}{x^2}\dd x.
        \]
        Therefore,
        \[
            \int_0^\infty f(x)\dd x
            =
            \frac12
            \int_0^\infty
            \left(f(x)+\frac{f(1/x)}{x^2}\right)\dd x.
        \]
        This often reveals a cancellation or symmetry hidden in an integral on
        $(0,\infty)$.
    \end{formula}

    \begin{formula}[Weierstrass Substitution; Tangent Half-Angle Substitution]
        The substitution
        \[
            t=\tan\frac{x}{2}
        \]
        converts rational expressions in $\sin x$ and $\cos x$ into rational
        functions of $t$:
        \[
            \sin x=\frac{2t}{1+t^2},
            \qquad
            \cos x=\frac{1-t^2}{1+t^2},
            \qquad
            \dd x=\frac{2\dd t}{1+t^2}.
        \]
    \end{formula}

    \begin{algorithm}[Fast Integral Recognition Checklist]
        Before expanding or integrating directly, check in this order:
        \begin{enumerate}
            \item parity, periodicity, and interval symmetry;
            \item the substitutions $x\mapsto a-x$, $x\mapsto 1/x$, or
                  $x\mapsto a/x$;
            \item a hidden derivative for substitution or integration by parts;
            \item a parameter that permits differentiation under the integral sign;
            \item Beta--Gamma, Gaussian, Wallis, or a standard trigonometric
                  integral;
            \item the tangent half-angle substitution for rational trigonometric
                  expressions.
        \end{enumerate}
    \end{algorithm}

    \begin{formula}[Wallis Product]
        The Wallis product is
        \[
            \frac{\pi}{2}
            =
            \prod_{n=1}^{\infty}
            \frac{(2n)(2n)}{(2n-1)(2n+1)}
            =
            \frac{2}{1}\cdot\frac{2}{3}\cdot
            \frac{4}{3}\cdot\frac{4}{5}\cdot
            \frac{6}{5}\cdot\frac{6}{7}\cdots.
        \]
    \end{formula}

    \begin{definition}[Gamma Function]
        For $x>0$, the Gamma function is defined by
        \[
            \Gamma(x)
            =
            \int_0^\infty t^{x-1}e^{-t}\dd t.
        \]
    \end{definition}

    \begin{proposition}[Gamma Function and Factorials]
        The Gamma function satisfies
        \[
            \Gamma(x+1)=x\Gamma(x).
        \]
        In particular, for every positive integer $n$,
        \[
            \Gamma(n)=(n-1)!.
        \]
        Also,
        \[
            \Gamma\left(\frac12\right)=\sqrt{\pi}.
        \]
    \end{proposition}

    \begin{definition}[Beta Function; Euler Beta Integral]
        For $x,y>0$, the Beta function is defined by
        \[
            B(x,y)
            =
            \int_0^1 t^{x-1}(1-t)^{y-1}\dd t.
        \]
    \end{definition}

    \begin{formula}[Beta--Gamma Identity]
        The Beta and Gamma functions are related by
        \[
            B(x,y)
            =
            \frac{\Gamma(x)\Gamma(y)}{\Gamma(x+y)}.
        \]
        In particular, for nonnegative integers $m,n$,
        \[
            \int_0^1 x^m(1-x)^n\dd x
            =
            B(m+1,n+1)
            =
            \frac{m!n!}{(m+n+1)!}
            =
            \frac{1}{(m+n+1)\binom{m+n}{m}}.
        \]
    \end{formula}

    \begin{formula}[Stirling's Formula]
        As $n\to\infty$,
        \[
            n!
            \sim
            \sqrt{2\pi n}\left(\frac{n}{e}\right)^n.
        \]
    \end{formula}

    \begin{remark}
        Stirling's formula is the standard approximation for large factorials.
        It often appears in counting, probability, and asymptotic estimation.
    \end{remark}

    \begin{formula}[Trapezoidal Rule]
        The trapezoidal approximation is
        \[
            \int_a^b f(x)\dd x
            \approx
            \frac{b-a}{2}\bigl(f(a)+f(b)\bigr).
        \]
        It has degree of precision $1$.
    \end{formula}

    \begin{formula}[Simpson's Rule]
        Simpson's rule is
        \[
            \int_a^b f(x)\dd x
            \approx
            \frac{b-a}{6}
            \left(
                f(a)+4f\left(\frac{a+b}{2}\right)+f(b)
            \right).
        \]
        It has degree of precision $3$.
    \end{formula}

    \begin{remark}
        Simpson's rule uses quadratic interpolation, but the resulting formula
        is exact for every polynomial of degree at most $3$.
    \end{remark}

    \begin{theorem}[$p$-Series Test]
        The series
        \[
            \sum_{n=1}^{\infty}\frac{1}{n^p}
        \]
        converges if and only if $p>1$.
    \end{theorem}

    \begin{theorem}[Direct Comparison Test]
        Let $0\leq a_n\leq b_n$ for all sufficiently large $n$. If
        $\sum b_n$ converges, then $\sum a_n$ converges. If $\sum a_n$
        diverges, then $\sum b_n$ diverges.
    \end{theorem}

    \begin{proposition}[Absolute Convergence Implies Convergence]
        If
        \[
            \sum_{n=1}^{\infty}|a_n|
        \]
        converges, then $\sum_{n=1}^{\infty}a_n$ converges.
    \end{proposition}

    \begin{theorem}[Integral Test]
        Suppose $f$ is positive, continuous, and decreasing on $[1,\infty)$,
        and let $a_n=f(n)$. Then
        \[
            \sum_{n=1}^{\infty}a_n
        \]
        and
        \[
            \int_1^{\infty}f(x)\dd x
        \]
        either both converge or both diverge.
    \end{theorem}

    \begin{theorem}[Limit Comparison Test]
        Let $a_n,b_n>0$. If
        \[
            \lim_{n\to\infty}\frac{a_n}{b_n}=c
        \]
        for some finite number $c>0$, then
        \[
            \sum_{n=1}^{\infty}a_n
            \qquad\text{and}\qquad
            \sum_{n=1}^{\infty}b_n
        \]
        converge or diverge together.
    \end{theorem}

    \begin{theorem}[Ratio Test; d'Alembert's Ratio Test]
        Let $a_n>0$ and suppose
        \[
            L=\lim_{n\to\infty}\frac{a_{n+1}}{a_n}
        \]
        exists. If $L<1$, then $\sum a_n$ converges. If $L>1$, then
        $\sum a_n$ diverges. If $L=1$, the test is inconclusive.
    \end{theorem}

    \begin{theorem}[Root Test; Cauchy's Root Test]
        Let $a_n\geq 0$ and suppose
        \[
            L=\lim_{n\to\infty}\sqrt[n]{a_n}
        \]
        exists. If $L<1$, then $\sum a_n$ converges. If $L>1$, then
        $\sum a_n$ diverges. If $L=1$, the test is inconclusive.
    \end{theorem}

    \begin{theorem}[Cauchy--Hadamard Theorem]
        For a power series
        \[
            \sum_{n=0}^{\infty}a_n(x-a)^n,
        \]
        the radius of convergence $R$ is given by
        \[
            \frac{1}{R}
            =
            \limsup_{n\to\infty}\sqrt[n]{|a_n|}.
        \]
        Equivalently, the series converges absolutely for $|x-a|<R$ and
        diverges for $|x-a|>R$.
    \end{theorem}

    \begin{remark}
        This is the standard theorem behind the phrase ``radius of convergence''
        for a power series. It is also the source of the name
        Cauchy--Hadamard.
    \end{remark}

    \begin{theorem}[Termwise Differentiation and Integration of Power Series]
        Suppose
        \[
            f(x)=\sum_{n=0}^{\infty}a_n(x-a)^n
        \]
        has radius of convergence $R>0$. For $|x-a|<R$,
        \[
            f'(x)
            =
            \sum_{n=1}^{\infty}n a_n(x-a)^{n-1},
        \]
        and
        \[
            \int_a^x f(t)\dd t
            =
            \sum_{n=0}^{\infty}\frac{a_n}{n+1}(x-a)^{n+1}.
        \]
        The differentiated and integrated series have the same radius of
        convergence $R$. Convergence at the endpoints must be checked
        separately.
    \end{theorem}

    \begin{theorem}[Alternating Series Test; Leibniz Test]
        If
        \[
            a_1\geq a_2\geq a_3\geq\cdots\geq 0
        \]
        and $a_n\to 0$, then the alternating series
        \[
            \sum_{n=1}^{\infty}(-1)^{n+1}a_n
        \]
        converges.
    \end{theorem}

    \begin{formula}[Alternating Series Error Estimate]
        Under the hypotheses of the Alternating Series Test, if
        \[
            S=\sum_{n=1}^{\infty}(-1)^{n+1}a_n,
            \qquad
            S_N=\sum_{n=1}^{N}(-1)^{n+1}a_n,
        \]
        then
        \[
            |S-S_N|\leq a_{N+1}.
        \]
    \end{formula}

    \begin{strategy}
        Factorials and powers often suggest the Ratio Test, $n$th powers suggest the Root Test, and
        alternating signs with decreasing terms suggest the Alternating Series
        Test.
    \end{strategy}

    \begin{theorem}[Divergence Test; Term Test]
        If the series
        \[
            \sum_{n=1}^{\infty}a_n
        \]
        converges, then
        \[
            \lim_{n\to\infty}a_n=0.
        \]
        Equivalently, if $\lim_{n\to\infty}a_n$ does not exist or is not $0$,
        then $\sum_{n=1}^{\infty}a_n$ diverges.
    \end{theorem}

    \begin{strategy}
        Before trying a sophisticated convergence test, first check whether the
        terms go to $0$. If the terms do not tend to $0$, the series diverges
        immediately.
    \end{strategy}

    \begin{algorithm}[Fast Series Checklist]
        For a convergence or summation problem:
        \begin{enumerate}
            \item Check the term limit first.
            \item Look for telescoping, a geometric series, or a known power
                  series.
            \item For positive terms, compare only the dominant asymptotic
                  factors.
            \item Use the Ratio Test for factorials and exponentials, and the
                  Root Test for expressions raised to the $n$th power.
            \item For alternating terms, check monotonicity and use the first
                  omitted term to control the error.
        \end{enumerate}
    \end{algorithm}

    \section*{Multivariable Calculus}

    \begin{definition}[Gradient, Divergence, Curl, and Laplacian]
        For a scalar-valued function $f(x,y,z)$ and a vector field
        $\mathbf{F}=(P,Q,R)$, we write
        \[
            \nabla f
            =
            \left(\frac{\partial f}{\partial x},
                  \frac{\partial f}{\partial y},
                  \frac{\partial f}{\partial z}\right),
        \]
        \[
            \nabla\cdot\mathbf{F}
            =
            \frac{\partial P}{\partial x}
            +
            \frac{\partial Q}{\partial y}
            +
            \frac{\partial R}{\partial z},
        \]
        \[
            \nabla\times\mathbf{F}
            =
            \begin{vmatrix}
                \mathbf{i} & \mathbf{j} & \mathbf{k} \\
                \partial/\partial x & \partial/\partial y & \partial/\partial z \\
                P & Q & R
            \end{vmatrix},
        \]
        and
        \[
            \Delta f
            =
            \nabla\cdot\nabla f
            =
            f_{xx}+f_{yy}+f_{zz}.
        \]
    \end{definition}

    \begin{idea}
        The gradient points in the direction of steepest increase. Divergence
        measures source or sink behavior. Curl measures local rotation. The
        Laplacian measures how a value differs from its surrounding average.
    \end{idea}

    \begin{theorem}[Clairaut's Theorem; Mixed Partial Derivative Theorem]
        If $f$ has continuous second partial derivatives near a point, then the
        mixed partial derivatives agree at that point:
        \[
            \frac{\partial^2 f}{\partial x\partial y}
            =
            \frac{\partial^2 f}{\partial y\partial x}.
        \]
    \end{theorem}

    \begin{definition}[Differentiability; Total Derivative; Fr\'echet Derivative]
        Let $\mathbf{F}:U\subseteq\R^n\to\R^m$, where $U$ is open, and
        let $\mathbf{a}\in U$. The map $\mathbf{F}$ is differentiable at
        $\mathbf{a}$ if there exists a linear map
        $D\mathbf{F}(\mathbf{a}):\R^n\to\R^m$ such that
        \[
            \mathbf{F}(\mathbf{a}+\mathbf{h})
            =
            \mathbf{F}(\mathbf{a})
            +D\mathbf{F}(\mathbf{a})\mathbf{h}
            +\mathbf{r}(\mathbf{h}),
        \]
        where
        \[
            \lim_{\mathbf{h}\to\mathbf{0}}
            \frac{\|\mathbf{r}(\mathbf{h})\|}{\|\mathbf{h}\|}
            =0.
        \]
        Equivalently,
        \[
            \begin{aligned}
                &\forall\eps>0,
                \quad
                \exists\delta>0,
                \quad
                \forall\mathbf{h}\in\R^n,\\
                &\mathbf{a}+\mathbf{h}\in U
                \text{ and }
                0<\|\mathbf{h}\|<\delta
                \implies\\
                &\qquad
                \frac{
                    \|\mathbf{F}(\mathbf{a}+\mathbf{h})
                    -\mathbf{F}(\mathbf{a})
                    -D\mathbf{F}(\mathbf{a})\mathbf{h}\|
                }{\|\mathbf{h}\|}
                <\eps.
            \end{aligned}
        \]
        The linear map $D\mathbf{F}(\mathbf{a})$ is the total derivative;
        in standard coordinates, its matrix is the Jacobian matrix.
    \end{definition}

    \begin{theorem}[Multivariable Chain Rule]
        Let $\mathbf{F}:\R^n\to\R^m$ and $\mathbf{G}:\R^m\to\R^k$ be
        differentiable. Then
        \[
            D(\mathbf{G}\circ\mathbf{F})(\mathbf{x})
            =
            D\mathbf{G}(\mathbf{F}(\mathbf{x}))D\mathbf{F}(\mathbf{x}).
        \]
    \end{theorem}

    \begin{theorem}[Fubini's Theorem]
        Under suitable continuity or integrability assumptions, an iterated
        integral over a plane region may be computed in either order. For
        example, if
        \[
            R=\{(x,y)\mid a\leq x\leq b,\ g_1(x)\leq y\leq g_2(x)\},
        \]
        then
        \[
            \iint_R f(x,y)\dd A
            =
            \int_a^b\int_{g_1(x)}^{g_2(x)}f(x,y)\dd y\dd x.
        \]
    \end{theorem}

    \begin{definition}[Jacobian]
        For a map
        \[
            \mathbf{F}(u_1,\dots,u_n)
            =
            (x_1(u_1,\dots,u_n),\dots,x_n(u_1,\dots,u_n)),
        \]
        the Jacobian matrix is
        \[
            D\mathbf{F}(\mathbf{u})
            =
            \left(\frac{\partial x_i}{\partial u_j}\right)_{i,j=1}^{n}.
        \]
        Its determinant
        \[
            \Jac(\mathbf{F})
            =
            \det(D\mathbf{F})
        \]
        is called the Jacobian determinant.
    \end{definition}

    \begin{theorem}[Change of Variables Formula]
        Under suitable regularity and one-to-one assumptions,
        \[
            \int_{\mathbf{F}(D)} f(\mathbf{x})\dd\mathbf{x}
            =
            \int_D f(\mathbf{F}(\mathbf{u}))
            \left|\det D\mathbf{F}(\mathbf{u})\right|\dd\mathbf{u}.
        \]
    \end{theorem}

    \begin{remark}
        The Jacobian determinant is the local area or volume scaling factor of
        a change of variables.
    \end{remark}

    \begin{formula}[Polar Coordinates]
        For polar coordinates $(r,\theta)$,
        \[
            x=r\cos\theta,
            \qquad
            y=r\sin\theta,
            \qquad
            r\geq0,
            \qquad
            0\leq\theta<2\pi.
        \]
        The area element is
        \[
            \dd A=r\dd r\dd\theta.
        \]
        The area enclosed by a polar curve $r=r(\theta)$ is
        \[
            A
            =
            \frac12\int_{\alpha}^{\beta}r(\theta)^2\dd\theta.
        \]
        Its arc length is
        \[
            L
            =
            \int_{\alpha}^{\beta}
            \sqrt{r(\theta)^2+\left(r'(\theta)\right)^2}\dd\theta.
        \]
    \end{formula}

    \begin{formula}[Cylindrical and Spherical Coordinates]
        Cylindrical coordinates $(r,\theta,z)$ satisfy
        \[
            x=r\cos\theta,
            \qquad
            y=r\sin\theta,
            \qquad
            z=z,
        \]
        with volume element
        \[
            \dd V=r\dd r\dd\theta\dd z.
        \]
        Spherical coordinates $(\rho,\theta,\phi)$ satisfy
        \[
            x=\rho\sin\phi\cos\theta,
            \qquad
            y=\rho\sin\phi\sin\theta,
            \qquad
            z=\rho\cos\phi,
        \]
        where
        \[
            \rho\geq0,
            \qquad
            0\leq\theta<2\pi,
            \qquad
            0\leq\phi\leq\pi.
        \]
        The volume element is
        \[
            \dd V
            =
            \rho^2\sin\phi\dd\rho\dd\theta\dd\phi.
        \]
        Here $\theta$ is the azimuthal angle measured from the positive
        $x$-axis in the $xy$-plane, and $\phi$ is the polar angle measured from
        the positive $z$-axis.
    \end{formula}

    \begin{formula}[Differential Operators in Standard Coordinates]
        Let
        \[
            \mathbf{F}
            =F_r\mathbf{e}_r+F_\theta\mathbf{e}_\theta+F_z\mathbf{e}_z
        \]
        in cylindrical coordinates. Then
        \[
            \nabla f
            =
            \frac{\partial f}{\partial r}\mathbf{e}_r
            +\frac{1}{r}\frac{\partial f}{\partial\theta}\mathbf{e}_\theta
            +\frac{\partial f}{\partial z}\mathbf{e}_z,
        \]
        \[
            \nabla\cdot\mathbf{F}
            =
            \frac{1}{r}\frac{\partial}{\partial r}(rF_r)
            +\frac{1}{r}\frac{\partial F_\theta}{\partial\theta}
            +\frac{\partial F_z}{\partial z},
        \]
        and
        \[
            \Delta f
            =
            \frac{1}{r}\frac{\partial}{\partial r}
            \left(r\frac{\partial f}{\partial r}\right)
            +\frac{1}{r^2}\frac{\partial^2f}{\partial\theta^2}
            +\frac{\partial^2f}{\partial z^2}.
        \]

        In spherical coordinates $(\rho,\theta,\phi)$, where $\theta$ is the
        azimuthal angle and $\phi$ is the polar angle, write
        \[
            \mathbf{F}
            =F_\rho\mathbf{e}_\rho
            +F_\theta\mathbf{e}_\theta
            +F_\phi\mathbf{e}_\phi.
        \]
        Then
        \[
            \nabla f
            =
            \frac{\partial f}{\partial\rho}\mathbf{e}_\rho
            +\frac{1}{\rho\sin\phi}
             \frac{\partial f}{\partial\theta}\mathbf{e}_\theta
            +\frac{1}{\rho}
             \frac{\partial f}{\partial\phi}\mathbf{e}_\phi,
        \]
        \[
            \nabla\cdot\mathbf{F}
            =
            \frac{1}{\rho^2}\frac{\partial}{\partial\rho}(\rho^2F_\rho)
            +\frac{1}{\rho\sin\phi}\frac{\partial F_\theta}{\partial\theta}
            +\frac{1}{\rho\sin\phi}
             \frac{\partial}{\partial\phi}(\sin\phi\,F_\phi),
        \]
        and
        \[
            \Delta f
            =
            \frac{1}{\rho^2}\frac{\partial}{\partial\rho}
            \left(\rho^2\frac{\partial f}{\partial\rho}\right)
            +\frac{1}{\rho^2\sin^2\phi}
             \frac{\partial^2f}{\partial\theta^2}
            +\frac{1}{\rho^2\sin\phi}\frac{\partial}{\partial\phi}
            \left(\sin\phi\frac{\partial f}{\partial\phi}\right).
        \]
    \end{formula}

    \begin{theorem}[Inverse Function Theorem]
        Let $\mathbf{F}:\R^n\to\R^n$ be continuously differentiable near
        $\mathbf{a}$. If
        \[
            \det D\mathbf{F}(\mathbf{a})\neq 0,
        \]
        then $\mathbf{F}$ has a continuously differentiable local inverse near
        $\mathbf{a}$.
    \end{theorem}

    \begin{theorem}[Implicit Function Theorem]
        Let $F(x,y)$ be continuously differentiable near $(a,b)$, with
        \[
            F(a,b)=0,
            \qquad
            \frac{\partial F}{\partial y}(a,b)\neq 0.
        \]
        Then, near $(a,b)$, the equation $F(x,y)=0$ determines $y$ as a
        differentiable function of $x$, and
        \[
            \diff{y}{x}
            =
            -\frac{F_x}{F_y}.
        \]
    \end{theorem}

    \begin{formula}[Second Derivative Test and Hessian]
        At a critical point $(a,b)$ of a twice differentiable function
        $f(x,y)$, let
        \[
            \Delta_H
            =
            f_{xx}(a,b)f_{yy}(a,b)-f_{xy}(a,b)^2.
        \]
        If $\Delta_H>0$ and $f_{xx}(a,b)>0$, the point is a local minimum. If
        $\Delta_H>0$ and $f_{xx}(a,b)<0$, it is a local maximum. If
        $\Delta_H<0$, it is a saddle point. If $\Delta_H=0$, the test is
        inconclusive.
    \end{formula}

    \begin{theorem}[Lagrange Multipliers]
        Suppose that $f$ has a local maximum or minimum subject to the constraint
        \[
            g(x_1,\dots,x_n)=0
        \]
        at a point where $\nabla g\neq\mathbf{0}$. Then there exists a scalar
        $\lambda$ such that
        \[
            \nabla f=\lambda\nabla g.
        \]
    \end{theorem}

    \begin{idea}
        At a constrained optimum, the level surface of $f$ is tangent to the
        constraint surface. Thus their normal vectors are parallel.
    \end{idea}

    \begin{theorem}[Euler's Homogeneous Function Theorem]
        Suppose $f:\R^n\to\R$ is differentiable and homogeneous of degree
        $k$, meaning
        \[
            f(t\mathbf{x})=t^k f(\mathbf{x})
        \]
        for $t>0$. Then
        \[
            \mathbf{x}\cdot\nabla f(\mathbf{x})=k f(\mathbf{x}).
        \]
        In coordinates,
        \[
            \sum_{i=1}^{n}x_i\frac{\partial f}{\partial x_i}=kf.
        \]
    \end{theorem}

    \begin{formula}[Parametrized Curves]
        Let $\mathbf{r}(t)$ be a twice differentiable curve. Its velocity,
        speed, and acceleration are
        \[
            \mathbf{v}(t)=\mathbf{r}'(t),
            \qquad
            \|\mathbf{v}(t)\|=\|\mathbf{r}'(t)\|,
            \qquad
            \mathbf{a}(t)=\mathbf{r}''(t).
        \]
        If $\mathbf{r}'(t_0)\neq\mathbf{0}$, the tangent line at $t_0$ is
        \[
            \mathbf{r}(t_0)+s\mathbf{r}'(t_0),
            \qquad
            s\in\R.
        \]
        The speed is constant if and only if
        \[
            \mathbf{r}'(t)\cdot\mathbf{r}''(t)=0.
        \]
    \end{formula}

    \begin{definition}[Scalar Line Integral]
        If a curve $C$ is parametrized by $\mathbf{r}(t)$ for $a\leq t\leq b$,
        then
        \[
            \int_C f\dd s
            =
            \int_a^b f(\mathbf{r}(t))\|\mathbf{r}'(t)\|\dd t.
        \]
        The value is unchanged by any regular reparametrization.
    \end{definition}

    \begin{definition}[Line Integral]
        If a curve $C$ is parametrized by $\mathbf{r}(t)$ for $a\leq t\leq b$,
        then the line integral of a vector field $\mathbf{F}$ along $C$ is
        \[
            \int_C \mathbf{F}\cdot\dd\mathbf{r}
            =
            \int_a^b \mathbf{F}(\mathbf{r}(t))\cdot\mathbf{r}'(t)\dd t.
        \]
    \end{definition}

    \begin{theorem}[Fundamental Theorem for Line Integrals; Gradient Theorem]
        If $\mathbf{F}=\nabla f$ and $C$ is a smooth curve from $\mathbf{a}$ to
        $\mathbf{b}$, then
        \[
            \int_C \mathbf{F}\cdot\dd\mathbf{r}
            =
            f(\mathbf{b})-f(\mathbf{a}).
        \]
        Hence the integral is path independent.
    \end{theorem}

    \begin{proposition}[Conservative-Field Test]
        On a simply connected open region, a continuously differentiable vector
        field is conservative if and only if its curl is zero.
    \end{proposition}

    \begin{theorem}[Green's Theorem]
        Let $C$ be a positively oriented, simple closed curve in the plane, and
        let $D$ be the region enclosed by $C$. Then
        \[
            \oint_C P\dd x+Q\dd y
            =
            \iint_D
            \left(
                \frac{\partial Q}{\partial x}
                -
                \frac{\partial P}{\partial y}
            \right)
            \dd A.
        \]
    \end{theorem}

    \begin{idea}
        Green's Theorem converts circulation along a boundary curve into curl
        over the region inside. It is the two-dimensional ancestor of Stokes'
        Theorem.
    \end{idea}

    \begin{theorem}[Stokes' Theorem]
        Let $S$ be an oriented smooth surface with positively oriented boundary
        curve $\partial S$. Then
        \[
            \oint_{\partial S}\mathbf{F}\cdot\dd\mathbf{r}
            =
            \iint_S (\nabla\times\mathbf{F})\cdot\mathbf{n}\dd S.
        \]
    \end{theorem}

    \begin{idea}
        Stokes' Theorem says that circulation around the boundary equals total
        curl through the surface. It turns a boundary integral into an interior
        integral.
    \end{idea}

    \begin{theorem}[Divergence Theorem; Gauss's Theorem]
        Let $E$ be a solid region in $\R^3$ with outward-oriented boundary
        surface $\partial E$. Then
        \[
            \iint_{\partial E}\mathbf{F}\cdot\mathbf{n}\dd S
            =
            \iiint_E \nabla\cdot\mathbf{F}\dd V.
        \]
    \end{theorem}

    \begin{idea}
        The Divergence Theorem says that outward flux through the boundary
        equals total source strength inside the region.
    \end{idea}

    \begin{remark}
        Green's Theorem, Stokes' Theorem, and the Divergence Theorem share the
        same philosophy: a boundary integral is equal to an interior integral.
        This is the geometric heart of vector calculus.
    \end{remark}

    \section*{Real Analysis}

    \begin{definition}[Metric Space and Open Ball]
        A metric space is a pair $(X,d)$ in which $d:X\times X\to\R$ satisfies,
        for all $x,y,z\in X$,
        \[
            d(x,y)\geq 0,
            \qquad
            d(x,y)=0\iff x=y,
        \]
        \[
            d(x,y)=d(y,x),
            \qquad
            d(x,z)\leq d(x,y)+d(y,z).
        \]
        The open ball of radius $r>0$ centered at $x$ is
        \[
            B_r(x)=\{y\in X\mid d(x,y)<r\}.
        \]
    \end{definition}

    \begin{definition}[Open and Closed Sets in a Metric Space]
        A subset $U\subseteq X$ is open if
        \[
            \forall x\in U,
            \quad
            \exists r>0,
            \qquad
            B_r(x)\subseteq U.
        \]
        A subset $F\subseteq X$ is closed if $X\setminus F$ is open.
    \end{definition}

    \begin{definition}[Limit of a Map Between Metric Spaces]
        Let $(X,d_X)$ and $(Y,d_Y)$ be metric spaces, let $D\subseteq X$, let
        $a$ be an accumulation point of $D$, and let $f:D\to Y$. The statement
        \[
            \lim_{x\to a}f(x)=L
        \]
        means
        \[
            \forall\eps>0,
            \quad
            \exists\delta>0,
            \quad
            \forall x\in D,
            \qquad
            0<d_X(x,a)<\delta
            \implies
            d_Y\bigl(f(x),L\bigr)<\eps.
        \]
    \end{definition}

    \begin{definition}[Convergence of a Sequence]
        A sequence $(x_n)$ in a metric space $(X,d)$ converges to $x\in X$,
        written $x_n\to x$, if
        \[
            \forall\eps>0,
            \quad
            \exists N\in\N,
            \quad
            \forall n\geq N,
            \qquad
            d(x_n,x)<\eps.
        \]
    \end{definition}

    \begin{definition}[Cauchy Sequence and Complete Metric Space]
        A sequence $(x_n)$ in $(X,d)$ is Cauchy if
        \[
            \forall\eps>0,
            \quad
            \exists N\in\N,
            \quad
            \forall m,n\geq N,
            \qquad
            d(x_m,x_n)<\eps.
        \]
        The metric space is complete if every Cauchy sequence in $X$ converges
        to a point of $X$.
    \end{definition}

    \begin{definition}[Continuity in Metric Spaces]
        Let $(X,d_X)$ and $(Y,d_Y)$ be metric spaces. A map $f:X\to Y$ is
        continuous at $x_0\in X$ if
        \[
            \forall\eps>0,
            \quad
            \exists\delta>0,
            \quad
            \forall x\in X,
            \qquad
            d_X(x,x_0)<\delta
            \implies
            d_Y\bigl(f(x),f(x_0)\bigr)<\eps.
        \]
        It is continuous on $X$ if it is continuous at every $x_0\in X$.
    \end{definition}

    \begin{theorem}[Sequential Criterion for Continuity]
        A map $f:X\to Y$ between metric spaces is continuous at $x\in X$ if
        and only if
        \[
            x_n\to x
            \implies
            f(x_n)\to f(x)
        \]
        for every sequence $(x_n)$ in $X$.
    \end{theorem}

    \begin{definition}[Uniform Continuity]
        A map $f:(X,d_X)\to(Y,d_Y)$ is uniformly continuous if
        \[
            \forall\eps>0,
            \quad
            \exists\delta>0,
            \quad
            \forall x,y\in X,
            \qquad
            d_X(x,y)<\delta
            \implies
            d_Y\bigl(f(x),f(y)\bigr)<\eps.
        \]
        Unlike ordinary continuity, the same $\delta$ must work for every pair
        $x,y\in X$.
    \end{definition}

    \begin{definition}[Pointwise and Uniform Convergence]
        Let $f_n:E\to Y$ and $f:E\to Y$, where $(Y,d)$ is a metric space. The
        sequence $(f_n)$ converges pointwise to $f$, written
        $f_n\ptwto f$, if
        \[
            \forall x\in E,
            \quad
            \forall\eps>0,
            \quad
            \exists N=N(x,\eps)\in\N,
            \quad
            \forall n\geq N,
            \qquad
            d\bigl(f_n(x),f(x)\bigr)<\eps.
        \]
        It converges uniformly to $f$, written $f_n\rightrightarrows f$, if
        \[
            \forall\eps>0,
            \quad
            \exists N=N(\eps)\in\N,
            \quad
            \forall n\geq N,
            \quad
            \forall x\in E,
            \qquad
            d\bigl(f_n(x),f(x)\bigr)<\eps.
        \]
        For real- or complex-valued functions, this is equivalent to
        \[
            \|f_n-f\|_\infty
            =
            \sup_{x\in E}|f_n(x)-f(x)|
            \longrightarrow 0.
        \]
    \end{definition}

    \begin{definition}[Equicontinuity and Uniform Boundedness]
        A family $\mathcal{F}$ of maps from $(X,d_X)$ to $(Y,d_Y)$ is
        equicontinuous at $x_0\in X$ if
        \[
            \forall\eps>0,
            \quad
            \exists\delta>0,
            \quad
            \forall f\in\mathcal{F},
            \quad
            \forall x\in X,
            \qquad
            d_X(x,x_0)<\delta
            \implies
            d_Y\bigl(f(x),f(x_0)\bigr)<\eps.
        \]
        A sequence of real-valued functions $(f_n)$ is uniformly bounded if
        \[
            \exists M>0,
            \quad
            \forall n\in\N,
            \quad
            \forall x\in X,
            \qquad
            |f_n(x)|\leq M.
        \]
    \end{definition}

    \begin{theorem}[Least Upper Bound Property; Completeness Axiom of $\R$]
        Every nonempty subset of $\R$ that is bounded above has a least upper
        bound in $\R$. Equivalently, every nonempty subset bounded below has a
        greatest lower bound.
    \end{theorem}

    \begin{definition}[Limit Superior and Limit Inferior]
        For a real sequence $(a_n)$, define
        \[
            \limsup_{n\to\infty} a_n
            =
            \lim_{n\to\infty}\sup_{k\geq n} a_k,
            \qquad
            \liminf_{n\to\infty} a_n
            =
            \lim_{n\to\infty}\inf_{k\geq n} a_k,
        \]
        whenever these extended real limits are allowed.
    \end{definition}

    \begin{theorem}[Monotone Convergence Theorem for Sequences]
        Every bounded monotone real sequence converges. More precisely, every
        increasing sequence bounded above converges to its supremum, and every
        decreasing sequence bounded below converges to its infimum.
    \end{theorem}

    \begin{theorem}[Bolzano--Weierstrass Theorem]
        Every bounded sequence in $\R^n$ has a convergent subsequence.
    \end{theorem}

    \begin{theorem}[Heine--Borel Theorem]
        A subset of $\R^n$ is compact if and only if it is closed and bounded.
    \end{theorem}

    \begin{theorem}[Heine--Cantor Theorem; Uniform Continuity on Compact Sets]
        If $K\subseteq\R^n$ is compact and $f:K\to\R^m$ is continuous, then $f$
        is uniformly continuous on $K$.
    \end{theorem}

    \begin{strategy}
        In analysis questions, compactness usually means that one may extract a
        convergent subsequence or that a continuous function attains extrema.
        In $\R^n$, the fast recognition phrase is ``closed and bounded.''
    \end{strategy}

    \begin{theorem}[Cauchy Criterion for Convergence]
        In a complete metric space, a sequence converges if and only if it is
        Cauchy. In particular, for $(a_n)$ in $\R$ or $\C$,
        \[
            (a_n)\text{ converges}
            \iff
            \forall\eps>0,
            \ \exists N\in\N,
            \ \forall m,n\geq N,
            \ |a_m-a_n|<\eps.
        \]
    \end{theorem}

    \begin{theorem}[Uniform Cauchy Criterion]
        Let $(Y,d)$ be complete and let $f_n:E\to Y$. Then $(f_n)$ converges
        uniformly on $E$ if and only if
        \[
            \forall\eps>0,
            \quad
            \exists N\in\N,
            \quad
            \forall m,n\geq N,
            \quad
            \forall x\in E,
            \qquad
            d\bigl(f_m(x),f_n(x)\bigr)<\eps.
        \]
    \end{theorem}

    \begin{theorem}[Weierstrass M-Test]
        Let $(f_n)$ be a sequence of functions on a set $E$. If
        \[
            |f_n(x)|\leq M_n
            \quad\text{for all }x\in E
        \]
        and
        \[
            \sum_{n=1}^{\infty}M_n
        \]
        converges, then the series
        \[
            \sum_{n=1}^{\infty} f_n(x)
        \]
        converges uniformly on $E$.
    \end{theorem}

    \begin{theorem}[Uniform Limit Theorem]
        If $(f_n)$ is a sequence of continuous functions on a set $E$ and
        $f_n\rightrightarrows f$ on $E$, then $f$ is continuous on $E$.
    \end{theorem}

    \begin{theorem}[Term-by-Term Differentiation]
        Suppose $(f_n)$ is a sequence of continuously differentiable functions
        on a closed bounded interval $[a,b]$, and suppose $(f_n(x_0))$
        converges for some $x_0\in[a,b]$. If $f_n'\rightrightarrows g$ on
        $[a,b]$, then $f_n\rightrightarrows f$ for a differentiable function
        $f$, and
        \[
            f'=g=\lim_{n\to\infty} f_n'.
        \]
    \end{theorem}

    \begin{theorem}[Arzelà--Ascoli Theorem]
        Every uniformly bounded and equicontinuous sequence of real-valued
        continuous functions on a compact metric space has a uniformly
        convergent subsequence.
    \end{theorem}

    \begin{theorem}[Stone--Weierstrass Theorem]
        A subalgebra of $C(K,\R)$ that contains the constants and separates
        points is dense in $C(K,\R)$ under the supremum norm, provided $K$ is
        compact.
    \end{theorem}

    \begin{theorem}[Banach Fixed-Point Theorem; Contraction Mapping Theorem]
        Let $(X,d)$ be a nonempty complete metric space, and suppose that
        $T:X\to X$ is a contraction: there exists $q\in[0,1)$ such that
        \[
            \forall x,y\in X,
            \qquad
            d\bigl(T(x),T(y)\bigr)\leq qd(x,y).
        \]
        Then $T$ has a unique fixed point $x^\ast\in X$, and the iteration
        $x_{n+1}=T(x_n)$ converges to $x^\ast$ for every initial point $x_0\in X$.
    \end{theorem}

    \begin{theorem}[Baire Category Theorem]
        A nonempty complete metric space cannot be written as a countable union
        of nowhere dense closed sets. Equivalently, a countable intersection of
        open dense subsets of a complete metric space is dense.
    \end{theorem}

    \begin{theorem}[Fatou's Lemma]
        If $f_n\geq 0$ are measurable, then
        \[
            \int \liminf_{n\to\infty}f_n\dd\mu
            \leq
            \liminf_{n\to\infty}\int f_n\dd\mu.
        \]
    \end{theorem}

    \begin{theorem}[Lebesgue Monotone Convergence Theorem; Beppo Levi Theorem]
        If $0\leq f_1\leq f_2\leq\cdots$ and $f_n\aeto f$,
        then
        \[
            \lim_{n\to\infty}\int f_n\dd\mu
            =
            \int f\dd\mu.
        \]
    \end{theorem}

    \begin{theorem}[Lebesgue Dominated Convergence Theorem]
        If $f_n\aeto f$ and there exists an integrable function
        $g$ such that $|f_n|\leq g$ for all $n$, then
        \[
            \lim_{n\to\infty}\int f_n\dd\mu
            =
            \int f\dd\mu.
        \]
    \end{theorem}

    \section*{Complex Analysis}

    \begin{definition}[Complex Differentiability; Holomorphic and Entire Functions]
        Let $U\subseteq\C$ be open and let $f:U\to\C$. The complex derivative of
        $f$ at $z_0\in U$ is
        \[
            f'(z_0)
            =
            \lim_{z\to z_0}
            \frac{f(z)-f(z_0)}{z-z_0},
        \]
        provided the limit exists independently of the direction of approach.
        The function is holomorphic on $U$ if it is complex differentiable at
        every point of $U$, and entire if it is holomorphic on all of $\C$.
    \end{definition}

    \begin{theorem}[Cauchy--Riemann Equations]
        Write
        \[
            f(z)=u(x,y)+iv(x,y),
            \qquad
            z=x+iy.
        \]
        If $f$ is complex differentiable at $z_0=(x_0,y_0)$ and the first partial
        derivatives exist there, then
        \[
            \frac{\partial u}{\partial x}(x_0,y_0)
            =
            \frac{\partial v}{\partial y}(x_0,y_0),
            \qquad
            \frac{\partial u}{\partial y}(x_0,y_0)
            =
            -\frac{\partial v}{\partial x}(x_0,y_0).
        \]
        Conversely, if these partial derivatives are continuous near $z_0$ and
        satisfy the equations there, then $f$ is complex differentiable at
        $z_0$.
    \end{theorem}

    \begin{strategy}
        To test whether a given $f(z)=u(x,y)+iv(x,y)$ can be holomorphic, compute
        the four first partial derivatives and check the Cauchy--Riemann equations.
        Failure at one point immediately rules out holomorphicity there.
    \end{strategy}

    \begin{proposition}[Holomorphic Functions Have Harmonic Components]
        If $f=u+iv$ is holomorphic and $u,v$ have continuous second partial
        derivatives, then
        \[
            \Delta u=0,
            \qquad
            \Delta v=0.
        \]
        Thus the real and imaginary parts of a holomorphic function are harmonic
        conjugates locally.
    \end{proposition}

    \begin{theorem}[Cauchy--Goursat Integral Theorem]
        If $f$ is holomorphic on a simply connected domain $U\subseteq\C$, then
        for every closed piecewise smooth curve $\gamma$ in $U$,
        \[
            \oint_{\gamma}f(z)\dd z=0.
        \]
        Equivalently, contour integrals of $f$ in $U$ are path-independent.
    \end{theorem}

    \begin{theorem}[Cauchy Integral Formula]
        Let $f$ be holomorphic on an open set containing a positively oriented
        simple closed contour $\gamma$ and its interior. If $a$ lies inside
        $\gamma$, then
        \[
            f(a)
            =
            \frac{1}{2\pi i}
            \oint_{\gamma}\frac{f(z)}{z-a}\dd z.
        \]
        More generally, for every integer $n\geq 0$,
        \[
            f^{(n)}(a)
            =
            \frac{n!}{2\pi i}
            \oint_{\gamma}
            \frac{f(z)}{(z-a)^{n+1}}\dd z.
        \]
    \end{theorem}

    \begin{strategy}
        When a contour integral contains $(z-a)^{-m}$ and the remaining factor
        is holomorphic, compare it directly with Cauchy's integral formula before
        attempting a parametrization of the contour.
    \end{strategy}

    \begin{theorem}[Maximum Modulus Principle]
        A nonconstant holomorphic function on a connected open set cannot attain
        a maximum of $|f|$ at an interior point. Consequently, if $f$ is
        holomorphic on a bounded domain and continuous on its closure, then the
        maximum of $|f|$ occurs on the boundary.
    \end{theorem}

    \begin{theorem}[Liouville's Theorem]
        Every bounded entire function is constant.
    \end{theorem}

    \begin{remark}
        Liouville's Theorem gives a short proof of the Fundamental Theorem of
        Algebra: if a nonconstant complex polynomial had no zero, then its
        reciprocal would be a bounded entire function.
    \end{remark}

    \begin{theorem}[Taylor Series for Holomorphic Functions]
        If $f$ is holomorphic on the open disk $|z-a|<R$, then
        \[
            f(z)
            =
            \sum_{n=0}^{\infty}
            \frac{f^{(n)}(a)}{n!}(z-a)^n,
            \qquad
            |z-a|<R.
        \]
        Thus a holomorphic function is analytic: complex differentiability implies
        local representability by a convergent power series.
    \end{theorem}

    \begin{theorem}[Laurent Series]
        If $f$ is holomorphic on an annulus
        \[
            r<|z-a|<R,
        \]
        then it has a unique Laurent expansion
        \[
            f(z)
            =
            \sum_{n=-\infty}^{\infty}c_n(z-a)^n
        \]
        converging throughout the annulus. The coefficient $c_{-1}$ is the residue
        of $f$ at $a$.
    \end{theorem}

    \begin{definition}[Isolated Singularities and Residues]
        Let $a$ be an isolated singularity of $f$. In the Laurent expansion about
        $a$, the residue is
        \[
            \Res(f;a)=c_{-1}.
        \]
        The singularity is removable if every negative-power coefficient is zero,
        a pole if only finitely many negative-power coefficients are nonzero, and
        essential if infinitely many are nonzero.
    \end{definition}

    \begin{formula}[Residues at Poles]
        If $a$ is a simple pole of $f$, then
        \[
            \Res(f;a)
            =
            \lim_{z\to a}(z-a)f(z).
        \]
        More generally, if $a$ is a pole of order $m$, then
        \[
            \Res(f;a)
            =
            \frac{1}{(m-1)!}
            \lim_{z\to a}
            \frac{\dd^{m-1}}{\dd z^{m-1}}
            \left((z-a)^m f(z)\right).
        \]
        If $f(z)=g(z)/h(z)$, where $g$ and $h$ are holomorphic,
        $h(a)=0$, and $h'(a)\neq 0$, then
        \[
            \Res(f;a)=\frac{g(a)}{h'(a)}.
        \]
    \end{formula}

    \begin{theorem}[Residue Theorem]
        Let $f$ be holomorphic inside and on a positively oriented simple closed
        contour $\gamma$, except for isolated singularities
        $a_1,\dots,a_m$ inside $\gamma$. Then
        \[
            \oint_{\gamma}f(z)\dd z
            =
            2\pi i
            \sum_{k=1}^{m}\Res(f;a_k).
        \]
    \end{theorem}

    \begin{strategy}[Real Integrals by Residues]
        For a rational integral over $\R$, extend the integrand to a complex
        rational function, choose a semicircular contour, sum the residues in the
        appropriate half-plane, and verify that the contribution from the large
        arc vanishes. Trigonometric integrals over $[0,2\pi]$ often become contour
        integrals under the substitution $z=e^{i\theta}$ and
        \[
            \dd\theta=\frac{\dd z}{iz}.
        \]
    \end{strategy}

    \begin{theorem}[Rouch\'e's Theorem]
        Let $f$ and $g$ be holomorphic inside and on a simple closed contour
        $\gamma$. If
        \[
            |g(z)|<|f(z)|
            \qquad
            \text{for every }z\in\gamma,
        \]
        then $f$ and $f+g$ have the same number of zeros inside $\gamma$, counted
        with multiplicity.
    \end{theorem}

    \begin{definition}[Conformal Map]
        A holomorphic function $f$ is conformal at $z_0$ if $f'(z_0)\neq 0$.
        At such a point, $f$ preserves oriented angles and infinitesimal shapes,
        although it may change scale.
    \end{definition}

    \begin{definition}[M\"obius Transformation; Linear Fractional Transformation]
        A M\"obius transformation is a map of the extended complex plane of the
        form
        \[
            T(z)=\frac{az+b}{cz+d},
            \qquad
            ad-bc\neq 0.
        \]
        It is conformal wherever its derivative is finite and nonzero, and it maps
        generalized circles---circles and lines---to generalized circles.
    \end{definition}

    \begin{formula}[Standard Conformal Maps]
        The following transformations are useful to recognize quickly:
        \[
            z\mapsto az+b
            \quad\text{(rotation, dilation, and translation)},
        \]
        \[
            z\mapsto \frac{1}{z}
            \quad\text{(inversion)},
            \qquad
            z\mapsto z^n
            \quad\text{(angle multiplication)},
        \]
        and
        \[
            z\mapsto \frac{z-i}{z+i},
        \]
        which maps the upper half-plane conformally onto the unit disk.
    \end{formula}

    \chapter{Linear Algebra}

    \relatedcourses{(MATH203) Applied Linear Algebra}

    \relatedbooks{Linear Algebra and Its Applications (Gilbert Strang)}

    \begin{definition}[Linear Combination, Span, and Linear Independence]
        A vector of the form
        \[
            c_1\mathbf{v}_1+\cdots+c_k\mathbf{v}_k
        \]
        is a linear combination of $\mathbf{v}_1,\dots,\mathbf{v}_k$. Their
        span is
        \[
            \Span(\mathbf{v}_1,\dots,\mathbf{v}_k)
            =
            \left\{
                c_1\mathbf{v}_1+\cdots+c_k\mathbf{v}_k
                :c_1,\dots,c_k\in F
            \right\}.
        \]
        The vectors are linearly independent if
        \[
            c_1\mathbf{v}_1+\cdots+c_k\mathbf{v}_k=\mathbf{0}
        \]
        implies $c_1=\cdots=c_k=0$.
    \end{definition}

    \begin{definition}[Linear Map; Linear Transformation]
        Let $V$ and $W$ be vector spaces over the same field $F$. A map
        $T:V\to W$ is linear if
        \[
            \forall\alpha,\beta\in F,
            \quad
            \forall\mathbf{u},\mathbf{v}\in V,
            \qquad
            T(\alpha\mathbf{u}+\beta\mathbf{v})
            =
            \alpha T(\mathbf{u})+\beta T(\mathbf{v}).
        \]
        Equivalently, $T$ preserves vector addition and scalar
        multiplication. Its kernel and image are
        \[
            \ker T=\{\mathbf{v}\in V\mid T(\mathbf{v})=\mathbf{0}\},
            \qquad
            \image T=\{T(\mathbf{v})\mid\mathbf{v}\in V\}.
        \]
    \end{definition}

    \begin{fact}[Translation Is Affine, Not Linear]
        For a fixed vector $\mathbf{v}\in\R^n$, the translation
        \[
            T_{\mathbf{v}}(\mathbf{x})
            =
            \mathbf{x}+\mathbf{v}
        \]
        is bijective, with inverse $T_{-\mathbf{v}}$. It is a linear map if and
        only if $\mathbf{v}=\mathbf{0}$.
    \end{fact}

    \begin{formula}[Inner Product and Adjoint Transfer]
        For real column vectors,
        \[
            \langle\mathbf{u},\mathbf{v}\rangle
            =
            \mathbf{u}^{\mathsf T}\mathbf{v}.
        \]
        For a real matrix $\mathbf{A}$ of compatible size,
        \[
            \langle\mathbf{x},\mathbf{A}\mathbf{y}\rangle
            =
            \langle\mathbf{A}^{\mathsf T}\mathbf{x},\mathbf{y}\rangle.
        \]
        Hence, if $\mathbf{A}$ is symmetric,
        \[
            \langle\mathbf{x},\mathbf{A}\mathbf{y}\rangle
            =
            \langle\mathbf{A}\mathbf{x},\mathbf{y}\rangle.
        \]
        Over $\C$, use
        \[
            \langle\mathbf{u},\mathbf{v}\rangle
            =
            \mathbf{u}^{\mathsf H}\mathbf{v}
        \]
        and replace $\mathbf{A}^{\mathsf T}$ by
        $\mathbf{A}^{\mathsf H}$.
    \end{formula}

    \begin{theorem}[Rank--Nullity Theorem]
        Let $T:V\to W$ be a linear map between finite-dimensional vector spaces.
        Then
        \[
            \dim V
            =
            \dim\ker T+\dim\image T.
        \]
        For a matrix $\mathbf{A}\in F^{m\times n}$, this becomes
        \[
            n=\rk(\mathbf{A})+\dim\Null(\mathbf{A}).
        \]
    \end{theorem}

    \begin{strategy}
        Rank--Nullity connects variables, pivots, and free variables. It is the
        linear algebra version of accounting: nothing disappears; dimensions
        are redistributed.
    \end{strategy}

    \begin{theorem}[Row Rank Equals Column Rank]
        For every matrix $\mathbf{A}$,
        \[
            \dim\Row(\mathbf{A})
            =
            \dim\Col(\mathbf{A}).
        \]
        This common dimension is called the rank of $\mathbf{A}$ and is denoted
        by $\rk(\mathbf{A})$.
    \end{theorem}

    \begin{theorem}[Invertible Matrix Theorem]
        For a square matrix $\mathbf{A}\in F^{n\times n}$, the following
        statements are equivalent:
        \[
            \mathbf{A}\text{ is invertible},
            \qquad
            \det(\mathbf{A})\neq 0,
            \qquad
            \rk(\mathbf{A})=n,
        \]
        \[
            \Null(\mathbf{A})=\{\mathbf{0}\},
        \]
        and
        \[
            \mathbf{A}\mathbf{x}=\mathbf{b}
        \]
        has a unique solution for every $\mathbf{b}\in F^n$.
    \end{theorem}

    \begin{remark}
        This theorem is a dictionary. It translates between determinants, rank,
        null spaces, and solvability of linear systems.
    \end{remark}

    \begin{proposition}[One-Sided Inverse for Square Matrices]
        If $\mathbf{A},\mathbf{B}\in F^{n\times n}$ satisfy
        \[
            \mathbf{A}\mathbf{B}=\mathbf{I}_n,
        \]
        then both matrices are invertible and
        \[
            \mathbf{B}\mathbf{A}=\mathbf{I}_n.
        \]
        Thus, for square matrices, a right inverse is automatically a left
        inverse, and conversely.
    \end{proposition}

    \begin{formula}[Inverse of a $2\times 2$ Matrix]
        If
        \[
            \mathbf{A}
            =
            \begin{pmatrix}
                a & b \\
                c & d
            \end{pmatrix}
        \]
        and $ad-bc\neq 0$, then
        \[
            \mathbf{A}^{-1}
            =
            \frac{1}{ad-bc}
            \begin{pmatrix}
                d & -b \\
                -c & a
            \end{pmatrix}.
        \]
    \end{formula}

    \begin{theorem}[Cramer's Rule]
        Let $\mathbf{A}\in F^{n\times n}$ be invertible. The unique solution of
        $\mathbf{A}\mathbf{x}=\mathbf{b}$ is given by
        \[
            x_i
            =
            \frac{\det(\mathbf{A}_i)}{\det(\mathbf{A})},
        \]
        where $\mathbf{A}_i$ is obtained from $\mathbf{A}$ by replacing its
        $i$th column with $\mathbf{b}$.
    \end{theorem}

    \begin{formula}[Adjugate Formula for the Inverse]
        For every square matrix $\mathbf{A}$,
        \[
            \mathbf{A}\adj(\mathbf{A})
            =
            \adj(\mathbf{A})\mathbf{A}
            =
            \det(\mathbf{A})\mathbf{I}_n.
        \]
        Hence, if $\det(\mathbf{A})\neq 0$, then
        \[
            \mathbf{A}^{-1}
            =
            \frac{1}{\det(\mathbf{A})}\adj(\mathbf{A}).
        \]
    \end{formula}

    \begin{definition}[Determinant]
        The determinant is a scalar associated to a square matrix
        $\mathbf{A}\in F^{n\times n}$, denoted by $\det(\mathbf{A})$.
        Geometrically, for real matrices, $|\det(\mathbf{A})|$ is the volume
        scaling factor of the linear transformation
        $\mathbf{x}\mapsto\mathbf{A}\mathbf{x}$.
    \end{definition}

    \begin{formula}[Determinant as Volume]
        If $\mathbf{v}_1,\dots,\mathbf{v}_n\in\R^n$, then the $n$-dimensional
        volume of the parallelepiped spanned by these vectors is
        \[
            \left|
            \det
            \begin{pmatrix}
                \mathbf{v}_1 & \mathbf{v}_2 & \cdots & \mathbf{v}_n
            \end{pmatrix}
            \right|.
        \]
        In $\R^3$, this is the absolute value of the scalar triple product
        \[
            |(\mathbf{a}\times\mathbf{b})\cdot\mathbf{c}|.
        \]
    \end{formula}

    \begin{formula}[Scalar and Vector Triple Products]
        For $\mathbf{a},\mathbf{b},\mathbf{c}\in\R^3$,
        \[
            \mathbf{a}\cdot(\mathbf{b}\times\mathbf{c})
            =
            \mathbf{b}\cdot(\mathbf{c}\times\mathbf{a})
            =
            \mathbf{c}\cdot(\mathbf{a}\times\mathbf{b})
            =
            \det(\mathbf{a},\mathbf{b},\mathbf{c}).
        \]
        The vector triple-product identities are
        \[
            \mathbf{a}\times(\mathbf{b}\times\mathbf{c})
            =
            (\mathbf{a}\cdot\mathbf{c})\mathbf{b}
            -
            (\mathbf{a}\cdot\mathbf{b})\mathbf{c},
        \]
        \[
            (\mathbf{a}\times\mathbf{b})\times\mathbf{c}
            =
            (\mathbf{a}\cdot\mathbf{c})\mathbf{b}
            -
            (\mathbf{b}\cdot\mathbf{c})\mathbf{a}.
        \]
    \end{formula}

    \begin{formula}[Cross-Product Determinant Identity]
        For $\mathbf{a},\mathbf{b},\mathbf{c}\in\R^3$,
        \[
            \det(
                \mathbf{a}\times\mathbf{b},
                \mathbf{b}\times\mathbf{c},
                \mathbf{c}\times\mathbf{a}
            )
            =
            \det(\mathbf{a},\mathbf{b},\mathbf{c})^2.
        \]
        Therefore, if $\mathbf{a},\mathbf{b},\mathbf{c}$ are linearly
        independent, then
        \[
            \mathbf{a}\times\mathbf{b},
            \qquad
            \mathbf{b}\times\mathbf{c},
            \qquad
            \mathbf{c}\times\mathbf{a}
        \]
        are also linearly independent.
    \end{formula}

    \begin{proposition}[Basic Properties of Determinants]
        For square matrices $\mathbf{A},\mathbf{B}\in F^{n\times n}$, we have
        \[
            \det(\mathbf{A}\mathbf{B})
            =
            \det(\mathbf{A})\det(\mathbf{B}),
        \]
        \[
            \det(\mathbf{A}^{\mathsf T})=\det(\mathbf{A}),
        \]
        and, if $\mathbf{A}$ is invertible,
        \[
            \det(\mathbf{A}^{-1})
            =
            \frac{1}{\det(\mathbf{A})}.
        \]
    \end{proposition}

    \begin{formula}[Determinant of a Triangular Matrix]
        If $\mathbf{A}$ is upper triangular or lower triangular with diagonal
        entries $a_{11},a_{22},\dots,a_{nn}$, then
        \[
            \det(\mathbf{A})
            =
            a_{11}a_{22}\cdots a_{nn}.
        \]
    \end{formula}

    \begin{formula}[Determinant of a Block Triangular Matrix]
        If the diagonal blocks are square, then
        \[
            \det
            \begin{pmatrix}
                \mathbf{A} & \mathbf{B} \\
                \mathbf{0} & \mathbf{D}
            \end{pmatrix}
            =
            \det(\mathbf{A})\det(\mathbf{D}),
        \]
        and the same formula holds for a lower block triangular matrix.
    \end{formula}

    \begin{proposition}[Multilinearity and Alternation of the Determinant]
        The determinant is linear in each row separately and also linear in each
        column separately. If two rows or two columns are interchanged, the sign
        of the determinant changes. In particular, if the rows or columns of
        $\mathbf{A}$ are linearly dependent, then
        \[
            \det(\mathbf{A})=0.
        \]
    \end{proposition}

    \begin{strategy}
        In one-minute determinant questions, first look for dependent rows or
        columns, row or column sums, repeated patterns, triangular form, and
        easy row operations. Direct expansion should usually be the last resort,
        except for $2\times 2$ and $3\times 3$ matrices.
    \end{strategy}

    \begin{proposition}[Characterization of the Determinant]
        Let
        \[
            f:(F^n)^n\to F
        \]
        be alternating and multilinear. Then
        \[
            f(\mathbf{v}_1,\dots,\mathbf{v}_n)
            =
            f(\mathbf{e}_1,\dots,\mathbf{e}_n)
            \det
            \begin{pmatrix}
                \mathbf{v}_1 & \cdots & \mathbf{v}_n
            \end{pmatrix}.
        \]
        In particular, the determinant is the unique alternating multilinear
        function satisfying
        \[
            \det(\mathbf{I}_n)=1.
        \]
    \end{proposition}

    \begin{formula}[Laplace Expansion; Cofactor Expansion]
        Let $\mathbf{A}=(a_{ij})\in F^{n\times n}$, and let $C_{ij}$ be the
        cofactor of $a_{ij}$. Expanding along row $i$ gives
        \[
            \det(\mathbf{A})
            =
            \sum_{j=1}^{n}a_{ij}C_{ij}.
        \]
        Expanding along column $j$ gives
        \[
            \det(\mathbf{A})
            =
            \sum_{i=1}^{n}a_{ij}C_{ij}.
        \]
    \end{formula}

    \begin{formula}[Sarrus' Rule]
        For a $3\times 3$ matrix,
        \[
            \det
            \begin{pmatrix}
                a & b & c \\
                d & e & f \\
                g & h & i
            \end{pmatrix}
            =
            aei+bfg+cdh-ceg-bdi-afh.
        \]
    \end{formula}

    \begin{formula}[Determinant and Elementary Row Operations]
        The determinant changes under elementary row operations as follows.
        \begin{itemize}
            \item Swapping two rows multiplies the determinant by $-1$.
            \item Multiplying one row by $c$ multiplies the determinant by $c$.
            \item Adding a multiple of one row to another row does not change
                  the determinant.
        \end{itemize}
        The same rules hold for columns.
    \end{formula}

    \begin{formula}[Matrix Determinant Lemma]
        If $\mathbf{A}\in F^{n\times n}$ is invertible and
        $\mathbf{u},\mathbf{v}\in F^n$, then
        \[
            \det(\mathbf{A}+\mathbf{u}\mathbf{v}^{\mathsf T})
            =
            \det(\mathbf{A})
            \left(1+\mathbf{v}^{\mathsf T}\mathbf{A}^{-1}\mathbf{u}\right).
        \]
        In particular,
        \[
            \det(\mathbf{I}_n+\mathbf{u}\mathbf{v}^{\mathsf T})
            =
            1+\mathbf{v}^{\mathsf T}\mathbf{u}.
        \]
    \end{formula}

    \begin{theorem}[Sylvester's Determinant Theorem]
        If $\mathbf{A}\in F^{m\times n}$ and
        $\mathbf{B}\in F^{n\times m}$, then
        \[
            \det(\mathbf{I}_m+\mathbf{A}\mathbf{B})
            =
            \det(\mathbf{I}_n+\mathbf{B}\mathbf{A}).
        \]
        This allows a determinant involving a large product to be replaced by a
        determinant of the smaller dimension.
    \end{theorem}

    \begin{formula}[Sherman--Morrison Formula]
        If $\mathbf{A}$ is invertible and
        $1+\mathbf{v}^{\mathsf T}\mathbf{A}^{-1}\mathbf{u}\neq0$, then
        \[
            (\mathbf{A}+\mathbf{u}\mathbf{v}^{\mathsf T})^{-1}
            =
            \mathbf{A}^{-1}
            -
            \frac{
                \mathbf{A}^{-1}\mathbf{u}\mathbf{v}^{\mathsf T}\mathbf{A}^{-1}
            }{
                1+\mathbf{v}^{\mathsf T}\mathbf{A}^{-1}\mathbf{u}
            }.
        \]
        It is the inverse counterpart of the Matrix Determinant Lemma.
    \end{formula}

    \begin{formula}[Schur Complement Determinant Formula]
        If $\mathbf{A}$ is invertible, then
        \[
            \det
            \begin{pmatrix}
                \mathbf{A} & \mathbf{B} \\
                \mathbf{C} & \mathbf{D}
            \end{pmatrix}
            =
            \det(\mathbf{A})
            \det(\mathbf{D}-\mathbf{C}\mathbf{A}^{-1}\mathbf{B}).
        \]
        The matrix $\mathbf{D}-\mathbf{C}\mathbf{A}^{-1}\mathbf{B}$ is the
        Schur complement of $\mathbf{A}$.
    \end{formula}

    \begin{formula}[Vandermonde Determinant]
        For scalars $x_1,x_2,\dots,x_n$,
        \[
            \det
            \begin{pmatrix}
                1 & x_1 & x_1^2 & \cdots & x_1^{n-1} \\
                1 & x_2 & x_2^2 & \cdots & x_2^{n-1} \\
                \vdots & \vdots & \vdots & \ddots & \vdots \\
                1 & x_n & x_n^2 & \cdots & x_n^{n-1}
            \end{pmatrix}
            =
            \prod_{1\leq i<j\leq n}(x_j-x_i).
        \]
    \end{formula}

    \begin{fact}[Pascal Matrices]
        Define the lower triangular Pascal matrix by
        \[
            \mathbf{P}_n=(p_{ij})_{0\leq i,j\leq n},
            \qquad
            p_{ij}
            =
            \begin{cases}
                \binom{i}{j}, & j\leq i, \\
                0, & j>i.
            \end{cases}
        \]
        Since all diagonal entries are equal to $1$, we have
        \[
            \det(\mathbf{P}_n)=1.
        \]
        Define the symmetric Pascal matrix by
        \[
            \mathbf{S}_n=(s_{ij})_{0\leq i,j\leq n},
            \qquad
            s_{ij}=\binom{i+j}{i}.
        \]
        This matrix also has determinant $1$.
    \end{fact}

    \begin{remark}
        Pascal matrices have combinatorial entries while their determinants admit
        simple closed forms.
    \end{remark}

    \begin{formula}[Cauchy Determinant]
        Let
        \[
            c_{ij}=\frac{1}{x_i+y_j},
        \]
        where every denominator is nonzero and the $x_i$ and $y_j$ are pairwise
        distinct within their respective families. Then
        \[
            \det(c_{ij})_{1\leq i,j\leq n}
            =
            \frac{
                \displaystyle
                \prod_{1\leq i<j\leq n}(x_j-x_i)(y_j-y_i)
            }{
                \displaystyle
                \prod_{i=1}^{n}\prod_{j=1}^{n}(x_i+y_j)
            }.
        \]
        The Hilbert matrix is a famous special case of a Cauchy matrix.
    \end{formula}


    \begin{formula}[Hilbert Matrix Determinant]
        The $n\times n$ Hilbert matrix is
        \[
            \mathbf{H}_n
            =
            \left(\frac{1}{i+j-1}\right)_{1\leq i,j\leq n}.
        \]
        As a special Cauchy matrix, it satisfies
        \[
            \det(\mathbf{H}_n)
            =
            \frac{\displaystyle\prod_{k=1}^{n-1}(k!)^4}
                 {\displaystyle\prod_{k=1}^{2n-1}k!}.
        \]
    \end{formula}

    \begin{definition}[Hadamard Matrix]
        A Hadamard matrix of order $n$ is an $n\times n$ matrix
        $\mathbf{H}$ whose entries are all $+1$ or $-1$ and whose rows are
        mutually orthogonal. Equivalently,
        \[
            \mathbf{H}\mathbf{H}^{\mathsf T}=n\mathbf{I}_n.
        \]
    \end{definition}

    \begin{formula}[Hadamard Determinant]
        If $\mathbf{H}$ is a Hadamard matrix of order $n$, then
        \[
            |\det(\mathbf{H})|
            =
            n^{n/2}.
        \]
    \end{formula}

    \begin{formula}[Permutation Matrices]
        A permutation matrix $\mathbf{P}_\sigma$ has exactly one entry equal
        to $1$ in each row and each column. It satisfies
        \[
            \mathbf{P}_\sigma^{-1}
            =
            \mathbf{P}_\sigma^{\mathsf T},
            \qquad
            \det(\mathbf{P}_\sigma)=\sgn(\sigma).
        \]
        Its powers are determined by the cycle structure of $\sigma$.
    \end{formula}

    \begin{formula}[Idempotent, Involutory, and Nilpotent Matrices]
        If $\mathbf{P}^2=\mathbf{P}$, then $\mathbf{P}$ is idempotent and its
        eigenvalues belong to $\{0,1\}$. If $\mathbf{A}^2=\mathbf{I}_n$, then
        $\mathbf{A}$ is involutory and its eigenvalues belong to
        $\{-1,1\}$. If $\mathbf{N}^k=\mathbf{0}$ for some $k$, then
        $\mathbf{N}$ is nilpotent and every eigenvalue is $0$.
    \end{formula}

    \begin{formula}[Eigenvalues of a Circulant Matrix]
        Let $\mathbf{C}$ be the circulant matrix whose first row is
        \[
            (c_0,c_1,\dots,c_{n-1}).
        \]
        With $\omega_n=e^{-2\pi i/n}$, its eigenvalues are
        \[
            \lambda_j
            =
            \sum_{k=0}^{n-1}c_k\omega_n^{jk},
            \qquad
            j=0,1,\dots,n-1.
        \]
        Hence $\det(\mathbf{C})=\prod_{j=0}^{n-1}\lambda_j$.
    \end{formula}


    \begin{formula}[Discrete Fourier Matrix]
        Let $\omega_n=e^{2\pi i/n}$. The normalized discrete Fourier matrix is
        \[
            \mathbf{F}_n
            =
            \frac{1}{\sqrt{n}}
            \left(\omega_n^{-jk}\right)_{0\leq j,k\leq n-1}.
        \]
        It is unitary:
        \[
            \mathbf{F}_n^{\mathsf H}\mathbf{F}_n=\mathbf{I}_n.
        \]
        Every circulant matrix is diagonalized by $\mathbf{F}_n$.
    \end{formula}

    \begin{formula}[Symmetric Tridiagonal Toeplitz Matrix]
        For
        \[
            \mathbf{T}_n
            =
            \begin{pmatrix}
                a & b & 0 & \cdots & 0 \\
                b & a & b & \ddots & \vdots \\
                0 & b & a & \ddots & 0 \\
                \vdots & \ddots & \ddots & \ddots & b \\
                0 & \cdots & 0 & b & a
            \end{pmatrix},
        \]
        the eigenvalues are
        \[
            a+2b\cos\frac{k\pi}{n+1},
            \qquad
            k=1,2,\dots,n.
        \]
    \end{formula}


    \begin{formula}[Continuant Recurrence for Tridiagonal Determinants]
        Let
        \[
            \mathbf{T}_n
            =
            \begin{pmatrix}
                a_1 & b_1 & 0 & \cdots & 0 \\
                c_1 & a_2 & b_2 & \ddots & \vdots \\
                0 & c_2 & a_3 & \ddots & 0 \\
                \vdots & \ddots & \ddots & \ddots & b_{n-1} \\
                0 & \cdots & 0 & c_{n-1} & a_n
            \end{pmatrix}.
        \]
        If $D_n=\det(\mathbf{T}_n)$, then
        \[
            D_0=1,
            \qquad
            D_1=a_1,
            \qquad
            D_n=a_nD_{n-1}-b_{n-1}c_{n-1}D_{n-2}.
        \]
    \end{formula}

    \begin{formula}[Householder Reflection]
        For a nonzero vector $\mathbf{v}$,
        \[
            \mathbf{H}
            =
            \mathbf{I}_n
            -2\frac{\mathbf{v}\mathbf{v}^{\mathsf T}}
                     {\mathbf{v}^{\mathsf T}\mathbf{v}}
        \]
        is symmetric, orthogonal, and involutory:
        \[
            \mathbf{H}^{\mathsf T}=\mathbf{H},
            \qquad
            \mathbf{H}^{\mathsf T}\mathbf{H}=\mathbf{I}_n,
            \qquad
            \mathbf{H}^2=\mathbf{I}_n.
        \]
        It reflects vectors across the hyperplane perpendicular to
        $\mathbf{v}$.
    \end{formula}

    \begin{formula}[Rotation and Reflection Matrices]
        Rotation by angle $\theta$ in the plane is represented by
        \[
            \begin{pmatrix}
                \cos\theta & -\sin\theta \\
                \sin\theta & \cos\theta
            \end{pmatrix}.
        \]
        Reflection across the line making angle $\theta$ with the positive
        $x$-axis is represented by
        \[
            \begin{pmatrix}
                \cos 2\theta & \sin 2\theta \\
                \sin 2\theta & -\cos 2\theta
            \end{pmatrix}.
        \]
    \end{formula}

    \begin{formula}[Eigenvalues of a Three-Dimensional Rotation Matrix]
        A rotation of $\R^3$ by angle $\theta$ around an axis through the
        origin has eigenvalues
        \[
            1,
            \qquad
            e^{i\theta},
            \qquad
            e^{-i\theta}.
        \]
        The eigenvalue $1$ corresponds to the rotation axis. Therefore, if
        $\theta\not\equiv 0,\pi\pmod{2\pi}$, the rotation matrix has exactly one
        real eigenvalue.
    \end{formula}

    \begin{proposition}[Trace, Determinant, and Eigenvalues]
        Let $\mathbf{A}\in\C^{n\times n}$ have eigenvalues
        $\lambda_1,\dots,\lambda_n$, counted with algebraic multiplicity. Then
        \[
            \tr(\mathbf{A})
            =
            \lambda_1+\cdots+\lambda_n
        \]
        and
        \[
            \det(\mathbf{A})
            =
            \lambda_1\lambda_2\cdots\lambda_n.
        \]
    \end{proposition}

    \begin{proposition}[Cyclic Property of the Trace]
        Whenever the products are defined and square,
        \[
            \tr(\mathbf{A}\mathbf{B})
            =
            \tr(\mathbf{B}\mathbf{A}).
        \]
        More generally,
        \[
            \tr(\mathbf{A}\mathbf{B}\mathbf{C})
            =
            \tr(\mathbf{B}\mathbf{C}\mathbf{A})
            =
            \tr(\mathbf{C}\mathbf{A}\mathbf{B}).
        \]
        The factors may be cyclically permuted, but they may not in general be
        rearranged arbitrarily.
    \end{proposition}

    \begin{definition}[Eigenvalue and Characteristic Polynomial]
        Let $\mathbf{A}\in F^{n\times n}$. A scalar $\lambda\in F$ is an
        eigenvalue of $\mathbf{A}$ if there exists a nonzero vector
        $\mathbf{v}\in F^n$ such that
        \[
            \mathbf{A}\mathbf{v}=\lambda\mathbf{v}.
        \]
        The characteristic polynomial of $\mathbf{A}$ is
        \[
            p_{\mathbf{A}}(\lambda)
            =
            \det(\lambda\mathbf{I}_n-\mathbf{A}).
        \]
    \end{definition}

    \begin{proposition}[Eigenvalues and the Characteristic Polynomial]
        A scalar $\lambda$ is an eigenvalue of $\mathbf{A}$ if and only if
        \[
            \det(\lambda\mathbf{I}_n-\mathbf{A})=0.
        \]
        Equivalently, the eigenvalues of $\mathbf{A}$ are the roots of the
        characteristic polynomial $p_{\mathbf{A}}(\lambda)$.
    \end{proposition}

    \begin{formula}[All-Ones Matrix]
        Let $\mathbf{J}_n$ be the $n\times n$ matrix whose entries are all $1$.
        Then
        \[
            \mathbf{J}_n\ones=n\ones,
            \qquad
            \mathbf{J}_n\mathbf{v}=\mathbf{0}
            \quad\text{if}\quad
            v_1+\cdots+v_n=0.
        \]
        Hence the eigenvalues of $\mathbf{J}_n$ are $n$ with multiplicity $1$
        and $0$ with multiplicity $n-1$. The matrix $\mathbf{J}_n-\mathbf{I}_n$
        has eigenvalues $n-1$ with multiplicity $1$ and $-1$ with multiplicity
        $n-1$.
    \end{formula}

    \begin{formula}[Constant Diagonal and Off-Diagonal Matrix]
        Let $\mathbf{A}\in F^{n\times n}$ have diagonal entries $a$ and
        off-diagonal entries $b$. Then
        \[
            \mathbf{A}=(a-b)\mathbf{I}_n+b\mathbf{J}_n.
        \]
        Counting algebraic multiplicity, its eigenvalues consist of
        \[
            a+(n-1)b
        \]
        once and
        \[
            a-b
        \]
        $n-1$ times. If these two values coincide, they together give the single
        eigenvalue of multiplicity $n$.
    \end{formula}

    \begin{formula}[Rank-One Outer Product]
        If $\mathbf{v},\mathbf{w}\in F^n$ are nonzero column vectors, then
        \[
            \mathbf{v}\mathbf{w}^{\mathsf T}
        \]
        has
        \[
            \rk(\mathbf{v}\mathbf{w}^{\mathsf T})=1.
        \]
        Its trace is
        \[
            \tr(\mathbf{v}\mathbf{w}^{\mathsf T})
            =
            \mathbf{w}^{\mathsf T}\mathbf{v},
        \]
        so its nonzero eigenvalue is $\mathbf{w}^{\mathsf T}\mathbf{v}$, and
        all remaining eigenvalues are $0$.
    \end{formula}

    \begin{theorem}[Cayley--Hamilton Theorem]
        Every square matrix satisfies its own characteristic polynomial. That
        is, if
        \[
            p_{\mathbf{A}}(\lambda)
            =
            \det(\lambda\mathbf{I}_n-\mathbf{A}),
        \]
        then
        \[
            p_{\mathbf{A}}(\mathbf{A})=\mathbf{0}.
        \]
    \end{theorem}

    \begin{remark}
        Cayley--Hamilton is a famous named theorem. It is especially useful for
        reducing high powers of a matrix to lower powers.
    \end{remark}


    \begin{formula}[Companion Matrix and Linear Recurrences]
        For
        \[
            p(x)=x^m-c_1x^{m-1}-\cdots-c_m,
        \]
        the companion matrix
        \[
            \mathbf{C}
            =
            \begin{pmatrix}
                c_1 & c_2 & \cdots & c_{m-1} & c_m \\
                1 & 0 & \cdots & 0 & 0 \\
                0 & 1 & \ddots & 0 & 0 \\
                \vdots & \ddots & \ddots & \vdots & \vdots \\
                0 & 0 & \cdots & 1 & 0
            \end{pmatrix}
        \]
        has characteristic polynomial $p$. The recurrence
        \[
            a_n=c_1a_{n-1}+\cdots+c_ma_{n-m}
        \]
        is encoded by multiplication by $\mathbf{C}$, so fast matrix
        exponentiation computes distant terms efficiently.
    \end{formula}

    \begin{formula}[Fibonacci $Q$-Matrix]
        For
        \[
            \mathbf{Q}
            =
            \begin{pmatrix}
                1 & 1 \\
                1 & 0
            \end{pmatrix},
        \]
        we have
        \[
            \mathbf{Q}^n
            =
            \begin{pmatrix}
                F_{n+1} & F_n \\
                F_n & F_{n-1}
            \end{pmatrix}
            \qquad(n\geq1).
        \]
        Taking determinants gives Cassini's identity
        \[
            F_{n+1}F_{n-1}-F_n^2=(-1)^n.
        \]
    \end{formula}

    \begin{formula}[Matrix Geometric Series]
        Let $\mathbf{A}$ be an $n\times n$ matrix. If
        $\mathbf{I}_n-\mathbf{A}$ is invertible, then
        \[
            \sum_{k=0}^{N-1}\mathbf{A}^k
            =
            (\mathbf{I}_n-\mathbf{A}^N)(\mathbf{I}_n-\mathbf{A})^{-1}.
        \]
        If the spectral radius of $\mathbf{A}$ is less than $1$, then
        \[
            \sum_{k=0}^{\infty}\mathbf{A}^k
            =
            (\mathbf{I}_n-\mathbf{A})^{-1}.
        \]
    \end{formula}

    \begin{strategy}[Fast Matrix Powers]
        To compute a large power $\mathbf{A}^N$, first look for one of the
        following: diagonalization, a short minimal polynomial, idempotence,
        involution, nilpotence, rank one, a circulant structure, or a decomposition
        into $a\mathbf{I}_n+b\mathbf{J}_n$. Cayley--Hamilton then reduces every
        sufficiently high power to a linear combination of lower powers.
    \end{strategy}

    \begin{theorem}[Diagonalization Criterion]
        A matrix $\mathbf{A}\in F^{n\times n}$ is diagonalizable over $F$ if
        and only if $F^n$ has a basis consisting of eigenvectors of
        $\mathbf{A}$.
    \end{theorem}

    \begin{corollary}[Distinct Eigenvalues Imply Diagonalizability]
        If an $n\times n$ matrix $\mathbf{A}$ has $n$ distinct eigenvalues in
        $F$, then $\mathbf{A}$ is diagonalizable over $F$.
    \end{corollary}

    \begin{proposition}[Polynomials of a Diagonalizable Matrix]
        Let $\mathbf{A}\in F^{n\times n}$ be diagonalizable over $F$, and suppose
        \[
            \mathbf{A}
            =
            \mathbf{S}\mathbf{D}\mathbf{S}^{-1},
            \qquad
            \mathbf{D}
            =
            \diag(\lambda_1,\dots,\lambda_n).
        \]
        For a polynomial $p(x)\in F[x]$, we have
        \[
            p(\mathbf{A})
            =
            \mathbf{S}p(\mathbf{D})\mathbf{S}^{-1},
        \]
        where
        \[
            p(\mathbf{D})
            =
            \diag\bigl(p(\lambda_1),\dots,p(\lambda_n)\bigr).
        \]
        In particular,
        \[
            \det(p(\mathbf{A}))
            =
            p(\lambda_1)p(\lambda_2)\cdots p(\lambda_n).
        \]
    \end{proposition}

    \begin{idea}
        Similar matrices represent the same linear transformation in different
        bases. Therefore, applying a polynomial to a diagonalizable matrix is
        the same as applying that polynomial to each diagonal entry after
        diagonalization.
    \end{idea}

    \begin{strategy}
        Distinct eigenvalues are an easy sufficient condition for
        diagonalizability. They are not necessary: a matrix can be
        diagonalizable even when some eigenvalues are repeated.
    \end{strategy}


    \begin{theorem}[Gershgorin Circle Theorem]
        Let $\mathbf{A}=(a_{ij})\in\C^{n\times n}$. Every eigenvalue of
        $\mathbf{A}$ lies in at least one disk
        \[
            \left\{z\in\C\mid
            |z-a_{ii}|\leq\sum_{j\neq i}|a_{ij}|
            \right\}.
        \]
    \end{theorem}

    \begin{theorem}[Perron--Frobenius Theorem, Positive-Matrix Form]
        If every entry of a real square matrix $\mathbf{A}$ is positive, then
        $\rho(\mathbf{A})$ is a positive simple eigenvalue with a positive
        eigenvector, and every other eigenvalue $\lambda$ satisfies
        \[
            |\lambda|<\rho(\mathbf{A}).
        \]
    \end{theorem}

    \begin{theorem}[Spectral Theorem for Real Symmetric Matrices]
        If $\mathbf{A}\in\R^{n\times n}$ is symmetric, then all eigenvalues of
        $\mathbf{A}$ are real, and there exists an orthogonal matrix
        $\mathbf{Q}$ and a diagonal matrix $\mathbf{D}$ such that
        \[
            \mathbf{A}
            =
            \mathbf{Q}\mathbf{D}\mathbf{Q}^{\mathsf T}.
        \]
        Equivalently, $\R^n$ has an orthonormal basis consisting of eigenvectors
        of $\mathbf{A}$.
    \end{theorem}

    \begin{remark}
        A real symmetric matrix has real eigenvalues and admits orthogonal
        diagonalization.
    \end{remark}

    \begin{definition}[Conjugate Transpose; Hermitian Transpose]
        For a complex matrix $\mathbf{A}=(a_{ij})$, its conjugate transpose, or
        Hermitian transpose, is
        \[
            \mathbf{A}^{\mathsf H}
            =
            \overline{\mathbf{A}}^{\mathsf T}
            =
            (\overline{a_{ji}}).
        \]
        A complex square matrix is Hermitian if
        \[
            \mathbf{A}^{\mathsf H}=\mathbf{A},
        \]
        and unitary if
        \[
            \mathbf{A}^{\mathsf H}\mathbf{A}
            =
            \mathbf{A}\mathbf{A}^{\mathsf H}
            =
            \mathbf{I}_n.
        \]
    \end{definition}

    \begin{proposition}[Basic Rules for the Hermitian Transpose]
        For complex matrices of compatible sizes,
        \[
            (\mathbf{A}+\mathbf{B})^{\mathsf H}
            =
            \mathbf{A}^{\mathsf H}+\mathbf{B}^{\mathsf H},
            \qquad
            (c\mathbf{A})^{\mathsf H}
            =
            \overline{c}\,\mathbf{A}^{\mathsf H},
        \]
        \[
            (\mathbf{A}\mathbf{B})^{\mathsf H}
            =
            \mathbf{B}^{\mathsf H}\mathbf{A}^{\mathsf H},
            \qquad
            (\mathbf{A}^{\mathsf H})^{\mathsf H}
            =
            \mathbf{A}.
        \]
        If $\mathbf{A}$ is real, then $\mathbf{A}^{\mathsf H}=\mathbf{A}^{\mathsf T}$.
    \end{proposition}

    \begin{proposition}[Hermitian Matrix Quick Facts]
        If $\mathbf{A}$ is Hermitian, then
        \[
            \mathbf{x}^{\mathsf H}\mathbf{A}\mathbf{x}\in\R
        \]
        for every $\mathbf{x}\in\C^n$. All eigenvalues of $\mathbf{A}$ are real,
        and eigenvectors corresponding to distinct eigenvalues are orthogonal.
    \end{proposition}

    \begin{theorem}[Spectral Theorem for Hermitian Matrices]
        If $\mathbf{A}\in\C^{n\times n}$ is Hermitian, then there exists a
        unitary matrix $\mathbf{U}$ and a real diagonal matrix $\mathbf{D}$ such
        that
        \[
            \mathbf{A}
            =
            \mathbf{U}\mathbf{D}\mathbf{U}^{\mathsf H}.
        \]
        Equivalently, $\C^n$ has an orthonormal basis consisting of eigenvectors
        of $\mathbf{A}$.
    \end{theorem}

    \begin{remark}
        The real analogue of a Hermitian matrix is a real symmetric matrix.
        The real analogue of a unitary matrix is an orthogonal matrix.
    \end{remark}

    \begin{definition}[Markov Matrix; Stochastic Matrix]
        A column-stochastic Markov matrix is a matrix whose entries are
        nonnegative and whose columns each sum to $1$.
        For such a matrix $\mathbf{A}$,
        \[
            \begin{pmatrix}1&1&\cdots&1\end{pmatrix}\mathbf{A}
            =
            \begin{pmatrix}1&1&\cdots&1\end{pmatrix}.
        \]
        Hence $1$ is an eigenvalue of $\mathbf{A}$.
    \end{definition}

    \begin{remark}
        Markov matrices are linear algebra objects behind finite-state Markov
        chains. A stationary distribution is an eigenvector corresponding to the
        eigenvalue $1$.
    \end{remark}

    \begin{theorem}[Cauchy--Schwarz Inequality]
        For vectors $\mathbf{u},\mathbf{v}\in\R^n$, we have
        \[
            |\langle\mathbf{u},\mathbf{v}\rangle|
            \leq
            \|\mathbf{u}\|\,\|\mathbf{v}\|.
        \]
        Equivalently, for real numbers $a_1,\dots,a_n$ and
        $b_1,\dots,b_n$,
        \[
            \left(\sum_{i=1}^{n}a_i b_i\right)^2
            \leq
            \left(\sum_{i=1}^{n}a_i^2\right)
            \left(\sum_{i=1}^{n}b_i^2\right).
        \]
    \end{theorem}

    \begin{strategy}
        The expression $\sum a_i b_i$ is a strong signal for Cauchy--Schwarz.
        The vector form is often the fastest way to recognize the inequality.
    \end{strategy}

    \begin{corollary}[Triangle Inequality]
        For vectors $\mathbf{u},\mathbf{v}\in\R^n$, we have
        \[
            \|\mathbf{u}+\mathbf{v}\|
            \leq
            \|\mathbf{u}\|+\|\mathbf{v}\|.
        \]
    \end{corollary}

    \begin{proposition}[Properties of Orthogonal Matrices]
        A real square matrix $\mathbf{Q}$ is orthogonal if
        \[
            \mathbf{Q}^{\mathsf T}\mathbf{Q}=\mathbf{I}_n.
        \]
        If $\mathbf{Q}$ is orthogonal, then
        \[
            \|\mathbf{Q}\mathbf{v}\|=\|\mathbf{v}\|,
            \qquad
            \langle\mathbf{Q}\mathbf{u},\mathbf{Q}\mathbf{v}\rangle
            =
            \langle\mathbf{u},\mathbf{v}\rangle,
            \qquad
            \det(\mathbf{Q})=\pm 1.
        \]
    \end{proposition}

    \begin{algorithm}[Gaussian Elimination]
        Gaussian elimination uses elementary row operations to transform a
        matrix into an upper triangular or row echelon form. It is the standard
        computational method for solving linear systems, finding ranks, and
        computing inverses.
    \end{algorithm}

    \begin{theorem}[LU Decomposition]
        When Gaussian elimination can be performed without row exchanges, a
        square matrix $\mathbf{A}$ can be written as
        \[
            \mathbf{A}=\mathbf{L}\mathbf{U},
        \]
        where $\mathbf{L}$ is lower triangular and $\mathbf{U}$ is upper
        triangular.
    \end{theorem}

    \begin{theorem}[Gram--Schmidt Process; QR Decomposition]
        From any linearly independent list
        \[
            \mathbf{v}_1,\dots,\mathbf{v}_k
        \]
        in an inner product space, one can construct an orthonormal list
        \[
            \mathbf{q}_1,\dots,\mathbf{q}_k
        \]
        with the same span. In matrix form, this leads to the QR decomposition
        \[
            \mathbf{A}=\mathbf{Q}\mathbf{R},
        \]
        where the columns of $\mathbf{Q}$ are orthonormal and $\mathbf{R}$ is
        upper triangular.
    \end{theorem}

    \begin{remark}
        Gram--Schmidt is the standard procedure for turning an arbitrary basis
        into an orthonormal basis.
    \end{remark}

    \begin{formula}[Orthogonal Projection Matrices]
        For a nonzero vector $\mathbf{v}\in\R^n$,
        \[
            \proj_{\mathbf{v}}\mathbf{x}
            =
            \frac{\mathbf{v}^{\mathsf T}\mathbf{x}}
                 {\mathbf{v}^{\mathsf T}\mathbf{v}}\mathbf{v},
        \]
        and the projection matrix onto $\Span(\mathbf{v})$ is
        \[
            \mathbf{P}_{\mathbf{v}}
            =
            \frac{\mathbf{v}\mathbf{v}^{\mathsf T}}
                 {\mathbf{v}^{\mathsf T}\mathbf{v}}.
        \]
        If $\|\mathbf{v}\|=1$, then
        \[
            \mathbf{P}_{\mathbf{v}}
            =
            \mathbf{v}\mathbf{v}^{\mathsf T},
            \qquad
            \mathbf{I}_n-\mathbf{v}\mathbf{v}^{\mathsf T}
        \]
        projects onto the hyperplane $\mathbf{v}^{\perp}$.

        If the columns of $\mathbf{Q}\in\R^{n\times k}$ are orthonormal, then
        the orthogonal projection matrix onto $\Col(\mathbf{Q})$ is
        \[
            \mathbf{P}
            =
            \mathbf{Q}\mathbf{Q}^{\mathsf T}.
        \]
        More generally, if the columns of $\mathbf{A}\in\R^{n\times k}$ are
        linearly independent, then
        \[
            \mathbf{P}
            =
            \mathbf{A}(\mathbf{A}^{\mathsf T}\mathbf{A})^{-1}
            \mathbf{A}^{\mathsf T}.
        \]
        Over $\C$, replace every transpose by the Hermitian transpose.
    \end{formula}

    \begin{proposition}[Projection Matrix Criterion]
        A matrix $\mathbf{P}$ is a projection matrix if
        \[
            \mathbf{P}^2=\mathbf{P}.
        \]
        Over $\R$, it is an orthogonal projection matrix if, in addition,
        \[
            \mathbf{P}^{\mathsf T}=\mathbf{P}.
        \]
        Over $\C$, the corresponding condition is
        \[
            \mathbf{P}^{\mathsf H}=\mathbf{P}.
        \]
    \end{proposition}

    \begin{definition}[Positive Definite Matrix]
        A real symmetric matrix $\mathbf{A}$ is positive definite if
        \[
            \mathbf{x}^{\mathsf T}\mathbf{A}\mathbf{x}>0
        \]
        for every nonzero vector $\mathbf{x}$.
    \end{definition}

    \begin{theorem}[Sylvester's Criterion]
        A real symmetric matrix is positive definite if and only if all of its
        leading principal minors are positive.
    \end{theorem}

    \begin{definition}[Singular Values]
        Let $\mathbf{A}\in\R^{m\times n}$. The singular values of
        $\mathbf{A}$ are the nonnegative square roots of the eigenvalues of
        \[
            \mathbf{A}^{\mathsf T}\mathbf{A}.
        \]
    \end{definition}

    \begin{theorem}[Singular Value Decomposition; SVD]
        Every real matrix $\mathbf{A}\in\R^{m\times n}$ can be written as
        \[
            \mathbf{A}
            =
            \mathbf{U}\boldsymbol{\Sigma}\mathbf{V}^{\mathsf T},
        \]
        where $\mathbf{U}$ and $\mathbf{V}$ are orthogonal matrices and
        $\boldsymbol{\Sigma}$ is a rectangular diagonal matrix whose diagonal
        entries are the singular values of $\mathbf{A}$.
    \end{theorem}

    \begin{remark}
        Singular Value Decomposition applies to every real matrix, not only to
        square or diagonalizable matrices. This is why SVD is one of the most
        useful decompositions in applied linear algebra.
    \end{remark}

    \begin{theorem}[Jordan Canonical Form]
        Over the complex numbers, every square matrix is similar to a Jordan
        matrix. That is, for every $\mathbf{A}\in\C^{n\times n}$, there exists
        an invertible matrix $\mathbf{P}$ such that
        \[
            \mathbf{P}^{-1}\mathbf{A}\mathbf{P}
        \]
        is in Jordan canonical form.
    \end{theorem}

    \begin{remark}
        Jordan canonical form describes the structure that remains when a matrix
        cannot be diagonalized.
    \end{remark}

    \chapter{Differential Equations}

    \relatedcourses{(MATH200) Differential Equations}

    \relatedbooks{Elementary Differential Equations and Boundary Value Problems
        (William E. Boyce, Richard C. DiPrima, Douglas B. Meade)}

    \begin{theorem}[Picard--Lindel\"of Existence and Uniqueness Theorem]
        Consider the initial value problem
        \[
            y'=f(t,y),
            \qquad
            y(t_0)=y_0.
        \]
        If $f$ and $\partial f/\partial y$ are continuous near
        $(t_0,y_0)$, then there exists a unique solution near $t=t_0$.
    \end{theorem}

    \begin{remark}
        The decisive conclusion is existence and uniqueness.
        It tells us when an initial condition determines exactly one solution.
    \end{remark}

    \begin{theorem}[Superposition Principle]
        If $y_1,\dots,y_k$ are solutions of a homogeneous linear differential
        equation, then any linear combination
        \[
            c_1y_1+\cdots+c_ky_k
        \]
        is also a solution.
    \end{theorem}

    \begin{remark}
        The Superposition Principle is the differential-equation version of
        linear algebra. It is one of the main reasons linear equations are much
        easier to handle than nonlinear equations.
    \end{remark}

    \begin{proposition}[General Solution of a Nonhomogeneous Linear Equation]
        For a nonhomogeneous linear differential equation, the general solution
        has the form
        \[
            y=y_h+y_p,
        \]
        where $y_h$ is the general solution of the associated homogeneous
        equation and $y_p$ is one particular solution of the nonhomogeneous
        equation.
    \end{proposition}

    \begin{formula}[Separable First-Order Equation]
        A first-order equation of the form
        \[
            \frac{\dd y}{\dd x}=g(x)h(y)
        \]
        is separable. When $h(y)\neq 0$, rewrite it as
        \[
            \frac{1}{h(y)}\dd y=g(x)\dd x
        \]
        and integrate both sides.
    \end{formula}

    \begin{formula}[First-Order Linear Equation; Integrating Factor]
        The equation
        \[
            y'+p(t)y=q(t)
        \]
        has integrating factor
        \[
            \mu(t)=e^{\int p(t)\dd t}.
        \]
        Then
        \[
            (\mu y)'=\mu q,
        \]
        so
        \[
            y(t)=\frac{1}{\mu(t)}
            \left(
                \int \mu(t)q(t)\dd t+C
            \right).
        \]
    \end{formula}

    \begin{formula}[Exact Differential Equation]
        An equation
        \[
            M(x,y)\dd x+N(x,y)\dd y=0
        \]
        is exact if
        \[
            \frac{\partial M}{\partial y}
            =
            \frac{\partial N}{\partial x}.
        \]
        Then there is a potential function $F(x,y)$ such that
        \[
            F_x=M,
            \qquad
            F_y=N,
        \]
        and the solution is given implicitly by
        \[
            F(x,y)=C.
        \]
    \end{formula}

    \begin{formula}[Bernoulli Equation]
        The Bernoulli equation
        \[
            y'+p(t)y=q(t)y^n
            \qquad(n\neq 0,1)
        \]
        becomes a first-order linear equation after the substitution
        \[
            u=y^{1-n}.
        \]
    \end{formula}

    \begin{tactic}
        For a first-order ODE, first check whether it is separable, linear,
        exact, or Bernoulli. These four patterns cover many quick undergraduate
        differential-equation questions.
    \end{tactic}

    \begin{formula}[Second-Order Constant-Coefficient Equation]
        Consider the homogeneous equation
        \[
            ay''+by'+cy=0,
            \qquad
            a\neq 0.
        \]
        Its characteristic equation is
        \[
            ar^2+br+c=0.
        \]
        If the roots are distinct real numbers $r_1,r_2$, then
        \[
            y_h=c_1e^{r_1t}+c_2e^{r_2t}.
        \]
        If the equation has a repeated real root $r$, then
        \[
            y_h=(c_1+c_2t)e^{rt}.
        \]
        If the roots are $\alpha\pm i\beta$ with $\beta\neq 0$, then
        \[
            y_h=e^{\alpha t}
            \left(
                c_1\cos \beta t+c_2\sin \beta t
            \right).
        \]
    \end{formula}

    \begin{remark}
        The characteristic equation converts a constant-coefficient linear ODE
        into an algebraic equation. This is the key trick.
    \end{remark}

    \begin{definition}[Wronskian]
        For two differentiable functions $y_1$ and $y_2$, their Wronskian is
        \[
            W(y_1,y_2)(t)
            =
            \begin{vmatrix}
                y_1(t) & y_2(t) \\
                y_1'(t) & y_2'(t)
            \end{vmatrix}
            =
            y_1(t)y_2'(t)-y_1'(t)y_2(t).
        \]
    \end{definition}

    \begin{proposition}[Wronskian and Linear Independence]
        If
        \[
            W(y_1,y_2)(t_0)\neq 0
        \]
        for some $t_0$, then $y_1$ and $y_2$ are linearly independent.
    \end{proposition}

    \begin{formula}[Abel's Identity]
        For
        \[
            y''+p(t)y'+q(t)y=0,
        \]
        the Wronskian of two solutions satisfies
        \[
            W(t)=Ce^{-\int p(t)\dd t}.
        \]
        In particular, the Wronskian is either always zero or never zero on an
        interval where $p$ is continuous.
    \end{formula}

    \begin{formula}[Euler--Cauchy Equation]
        The equation
        \[
            ax^2y''+bxy'+cy=0
        \]
        suggests the trial solution $y=x^r$. This gives the indicial equation
        \[
            ar(r-1)+br+c=0.
        \]
    \end{formula}

    \begin{algorithm}[Undetermined Coefficients]
        For a constant-coefficient equation with forcing term made from
        polynomials, exponentials, sines, or cosines, guess a particular solution
        of the same type. If the guess overlaps with $y_h$, multiply the guess
        by the smallest power of $t$ that removes the overlap.
    \end{algorithm}

    \begin{formula}[Variation of Parameters]
        For
        \[
            y''+p(t)y'+q(t)y=g(t)
        \]
        with fundamental solutions $y_1,y_2$ of the homogeneous equation, a
        particular solution is
        \[
            y_p
            =
            -y_1\int\frac{y_2g}{W}\dd t
            +
            y_2\int\frac{y_1g}{W}\dd t,
        \]
        where
        \[
            W=y_1y_2'-y_1'y_2.
        \]
    \end{formula}

    \begin{formula}[Exponential Growth and Decay]
        The equation
        \[
            y'=ky
        \]
        has solution
        \[
            y(t)=Ce^{kt}.
        \]
        If $k>0$, this models exponential growth. If $k<0$, this models
        exponential decay.
    \end{formula}

    \begin{formula}[Logistic Equation]
        The logistic equation is
        \[
            y'=ky\left(1-\frac{y}{K}\right),
        \]
        where $k>0$ is the growth rate and $K>0$ is the carrying capacity.
        Its equilibrium solutions are
        \[
            y=0
            \qquad\text{and}\qquad
            y=K.
        \]
    \end{formula}

    \begin{remark}
        The logistic equation is the standard model for population growth with
        limited resources. The parameter $K$ represents the limiting population.
    \end{remark}

    \begin{formula}[Simple Harmonic Oscillator]
        The equation
        \[
            x''+\omega^2x=0
        \]
        has general solution
        \[
            x(t)=c_1\cos \omega t+c_2\sin \omega t.
        \]
        Equivalently,
        \[
            x(t)=A\cos(\omega t-\phi)
        \]
        for suitable constants $A$ and $\phi$.
    \end{formula}

    \begin{remark}
        This is the equation behind ideal mass--spring motion and many small
        oscillation models in physics.
    \end{remark}

    \begin{formula}[Damped Harmonic Oscillator]
        The equation
        \[
            mx''+cx'+kx=0
        \]
        models a damped harmonic oscillator. Its characteristic equation is
        \[
            mr^2+cr+k=0.
        \]
        The motion is underdamped, critically damped, or overdamped according
        as
        \[
            c^2-4mk<0,
            \qquad
            c^2-4mk=0,
            \qquad
            c^2-4mk>0.
        \]
    \end{formula}

    \begin{definition}[Laplace Transform]
        The Laplace transform of a function $f(t)$ is
        \[
            \mathcal{L}\{f(t)\}(s)
            =
            \int_0^\infty e^{-st}f(t)\dd t.
        \]
    \end{definition}

    \begin{formula}[Laplace Transform Table]
        Let $F(s)=\mathcal{L}\{f(t)\}(s)$. The basic transforms are

        {\small
        \renewcommand{\arraystretch}{1.18}
        \begin{longtable}{@{}L{0.38\textwidth}L{0.42\textwidth}L{0.12\textwidth}@{}}
            \toprule
            $f(t)$ & $\mathcal{L}\{f(t)\}(s)$ & Conditions \\
            \midrule
            \endfirsthead

            \toprule
            $f(t)$ & $\mathcal{L}\{f(t)\}(s)$ & Conditions \\
            \midrule
            \endhead

            \bottomrule
            \endlastfoot

            $1$ & $\displaystyle\frac{1}{s}$ & $s>0$ \\
            $t^n$ & $\displaystyle\frac{n!}{s^{n+1}}$ & $n\in\Nzero$ \\
            $e^{at}$ & $\displaystyle\frac{1}{s-a}$ & $s>a$ \\
            $\sin at$ & $\displaystyle\frac{a}{s^2+a^2}$ & $s>0$ \\
            $\cos at$ & $\displaystyle\frac{s}{s^2+a^2}$ & $s>0$ \\
            $\sinh at$ & $\displaystyle\frac{a}{s^2-a^2}$ & $s>|a|$ \\
            $\cosh at$ & $\displaystyle\frac{s}{s^2-a^2}$ & $s>|a|$ \\
            $f'(t)$ & $sF(s)-f(0)$ & \\
            $f''(t)$ & $s^2F(s)-sf(0)-f'(0)$ & \\
            $t f(t)$ & $-F'(s)$ & \\
            $\displaystyle\int_0^t f(\tau)\dd\tau$ & $\displaystyle\frac{F(s)}{s}$ & \\
            $u(t-a)f(t-a)$ & $e^{-as}F(s)$ & $a>0$ \\
        \end{longtable}
        }
    \end{formula}

    \begin{proposition}[Laplace Shift and Convolution]
        If $F(s)=\mathcal{L}\{f(t)\}(s)$, then
        \[
            \mathcal{L}\{e^{at}f(t)\}(s)=F(s-a).
        \]
        If
        \[
            (f*g)(t)=\int_0^t f(\tau)g(t-\tau)\dd\tau,
        \]
        then
        \[
            \mathcal{L}\{f*g\}=\mathcal{L}\{f\}\mathcal{L}\{g\}.
        \]
    \end{proposition}

    \begin{formula}[Fourier Series]
        For a piecewise smooth function $f$ on $[-L,L]$,
        \[
            f(x)
            \sim
            \frac{a_0}{2}
            +\sum_{n=1}^{\infty}
            \left(
                a_n\cos\frac{n\pi x}{L}
                +b_n\sin\frac{n\pi x}{L}
            \right),
        \]
        where
        \[
            a_n
            =
            \frac{1}{L}\int_{-L}^{L}
            f(x)\cos\frac{n\pi x}{L}\dd x,
            \qquad
            b_n
            =
            \frac{1}{L}\int_{-L}^{L}
            f(x)\sin\frac{n\pi x}{L}\dd x.
        \]
        If $f$ is even, then $b_n=0$; if $f$ is odd, then $a_n=0$.
        At a jump discontinuity, the series converges to the midpoint of the
        left and right limits.
    \end{formula}

    \begin{formula}[Classical Fourier Series]
        For $-\pi<x<\pi$,
        \[
            x
            =
            2\sum_{n=1}^{\infty}
            \frac{(-1)^{n+1}}{n}\sin(nx),
        \]
        \[
            x^2
            =
            \frac{\pi^2}{3}
            +4\sum_{n=1}^{\infty}
            \frac{(-1)^n}{n^2}\cos(nx),
        \]
        and, for $0<x<\pi$,
        \[
            1
            =
            \frac{4}{\pi}
            \sum_{k=0}^{\infty}
            \frac{\sin((2k+1)x)}{2k+1}.
        \]
    \end{formula}

    \begin{definition}[Fourier Transform]
        For an integrable function $f:\R\to\C$, define
        \[
            \widehat{f}(\omega)
            =
            \int_{-\infty}^{\infty}f(t)e^{-i\omega t}\dd t.
        \]
        Under suitable hypotheses, the inversion formula is
        \[
            f(t)
            =
            \frac{1}{2\pi}
            \int_{-\infty}^{\infty}\widehat{f}(\omega)e^{i\omega t}\dd\omega.
        \]
    \end{definition}

    \begin{formula}[Fourier Transform Table]
        With the convention above, the standard identities are

        {\small
        \renewcommand{\arraystretch}{1.18}
        \begin{longtable}{@{}L{0.42\textwidth}L{0.46\textwidth}@{}}
            \toprule
            Function or operation & Fourier transform \\
            \midrule
            \endfirsthead

            \toprule
            Function or operation & Fourier transform \\
            \midrule
            \endhead

            \bottomrule
            \endlastfoot

            $f'(t)$ & $i\omega\widehat{f}(\omega)$ \\
            $f(t-a)$ & $e^{-i\omega a}\widehat{f}(\omega)$ \\
            $e^{iat}f(t)$ & $\widehat{f}(\omega-a)$ \\
            $f(at)$, $a\neq0$ & $\displaystyle\frac{1}{|a|}\widehat{f}\left(\frac{\omega}{a}\right)$ \\
            $t f(t)$ & $\displaystyle i\frac{\dd}{\dd\omega}\widehat{f}(\omega)$ \\
            $(f*g)(t)$ & $\widehat{f}(\omega)\widehat{g}(\omega)$ \\
            $e^{-at^2}$, $a>0$ & $\displaystyle\sqrt{\frac{\pi}{a}}e^{-\omega^2/(4a)}$ \\
        \end{longtable}
        }
    \end{formula}

    \begin{formula}[Linear Systems and Matrix Exponential]
        The homogeneous linear system
        \[
            \mathbf{x}'=\mathbf{A}\mathbf{x}
        \]
        has solution
        \[
            \mathbf{x}(t)=e^{t\mathbf{A}}\mathbf{x}(0).
        \]
        If $\mathbf{A}$ is diagonalizable as
        $\mathbf{A}=\mathbf{S}\mathbf{D}\mathbf{S}^{-1}$, then
        \[
            e^{t\mathbf{A}}
            =
            \mathbf{S}e^{t\mathbf{D}}\mathbf{S}^{-1},
            \qquad
            e^{t\mathbf{D}}
            =
            \diag(e^{\lambda_1t},\dots,e^{\lambda_nt}).
        \]
    \end{formula}

    \begin{strategy}
        For a constant-coefficient linear system, the eigenvalues of
        $\mathbf{A}$ control the qualitative behavior. Negative real parts
        indicate decay, positive real parts indicate growth, and nonzero
        imaginary parts indicate rotation or oscillation.
    \end{strategy}

    \begin{remark}
        The Laplace transform turns many differential equations into algebraic
        equations. This is its main computational advantage.
    \end{remark}

    \begin{definition}[Heat Equation]
        The heat equation is the partial differential equation
        \[
            u_t=\alpha u_{xx},
            \qquad
            \alpha>0.
        \]
        It models diffusion, especially the flow of heat in a one-dimensional
        medium.
    \end{definition}

    \begin{definition}[Wave Equation]
        The one-dimensional wave equation is
        \[
            u_{tt}=c^2u_{xx}.
        \]
        It models waves traveling with speed $c$.
    \end{definition}

    \begin{theorem}[d'Alembert's Formula]
        The solution of the wave equation
        \[
            u_{tt}=c^2u_{xx}
        \]
        on the real line with initial conditions
        \[
            u(x,0)=f(x),
            \qquad
            u_t(x,0)=g(x),
        \]
        is
        \[
            u(x,t)
            =
            \frac{f(x-ct)+f(x+ct)}{2}
            +
            \frac{1}{2c}
            \int_{x-ct}^{x+ct}g(s)\dd s.
        \]
    \end{theorem}

    \begin{remark}
        d'Alembert's formula shows that waves move along the characteristic
        lines $x-ct=\text{constant}$ and $x+ct=\text{constant}$.
    \end{remark}

    \begin{definition}[Laplace Equation]
        Laplace's equation is
        \[
            \Delta u=0.
        \]
        In two variables, this means
        \[
            u_{xx}+u_{yy}=0.
        \]
        A function satisfying Laplace's equation is called harmonic.
    \end{definition}

    \begin{definition}[Poisson Equation]
        Poisson's equation is
        \[
            \Delta u=f.
        \]
        It is an inhomogeneous version of Laplace's equation.
    \end{definition}

    \begin{remark}
        The heat equation, wave equation, and Laplace equation are three fundamental
        partial differential equations with distinct physical interpretations.
    \end{remark}

    \chapter{Combinatorics and Discrete Mathematics}

    \relatedcourses{(MATH261) Discrete Mathematics}

    \relatedbooks{Introductory Combinatorics (Richard A. Brualdi)}

    \begin{formula}[Binomial Theorem; Newton's Binomial Formula]
        For every nonnegative integer $n$,
        \[
            (x+y)^n
            =
            \sum_{k=0}^{n}\binom{n}{k}x^{n-k}y^k.
        \]
        In particular,
        \[
            (1+x)^n
            =
            \sum_{k=0}^{n}\binom{n}{k}x^k.
        \]
    \end{formula}

    \begin{remark}
        The coefficients $\binom{n}{k}$ are the entries of Pascal's triangle.
        This is one of the most recognizable bridges between algebra and
        counting.
    \end{remark}


    \begin{strategy}[Bijection, Double Counting, and Complementary Counting]
        Before expanding a complicated sum, try to count the same objects in a
        second way. Bijections, double counting, and counting the complement
        are among the fastest and most reusable contest methods.
    \end{strategy}

    \begin{formula}[Binomial Inversion]
        The relations
        \[
            b_n=\sum_{k=0}^{n}\binom{n}{k}a_k
        \]
        and
        \[
            a_n=\sum_{k=0}^{n}(-1)^{n-k}\binom{n}{k}b_k
        \]
        are equivalent.
    \end{formula}

    \begin{formula}[Pascal's Identity and Vandermonde's Identity]
        Pascal's identity is
        \[
            \binom{n}{k}
            =
            \binom{n-1}{k-1}+\binom{n-1}{k}.
        \]
        Vandermonde's identity is
        \[
            \sum_{k=0}^{r}\binom{m}{k}\binom{n}{r-k}
            =
            \binom{m+n}{r}.
        \]
    \end{formula}

    \begin{formula}[Multinomial Theorem]
        For nonnegative integers $k_1,\dots,k_m$ with
        $k_1+\cdots+k_m=n$,
        \[
            (x_1+\cdots+x_m)^n
            =
            \sum_{k_1+\cdots+k_m=n}
            \frac{n!}{k_1!\cdots k_m!}
            x_1^{k_1}\cdots x_m^{k_m}.
        \]
        The coefficient
        \[
            \frac{n!}{k_1!\cdots k_m!}
        \]
        is called a multinomial coefficient.
    \end{formula}

    \begin{formula}[Stars and Bars]
        The number of nonnegative integer solutions of
        \[
            x_1+x_2+\cdots+x_n=k
        \]
        is
        \[
            \binom{n+k-1}{k}
            =
            \multichoose{n}{k}.
        \]
        Equivalently, the expansion of $(x_1+x_2+\cdots+x_n)^k$ has
        \[
            \binom{n+k-1}{k}
        \]
        distinct monomial terms.
    \end{formula}

    \begin{formula}[Lattice Path Counting]
        The number of shortest lattice paths from $(0,0)$ to $(m,n)$ using only
        unit steps to the right and upward is
        \[
            \binom{m+n}{m}
            =
            \binom{m+n}{n}.
        \]
    \end{formula}

    \begin{theorem}[Reflection Principle and Ballot-Path Formula]
        The reflection principle converts a path that first crosses a boundary
        into an unrestricted reflected path. In particular, for integers
        $a\geq b\geq0$, the number of paths from $(0,0)$ to $(a,b)$ using right
        and up steps and never entering the region $y>x$ is
        \[
            \binom{a+b}{b}-\binom{a+b}{b-1}
            =
            \frac{a-b+1}{a+1}\binom{a+b}{b}.
        \]
        When $a=b=n$, this becomes the Catalan number
        \[
            \frac{1}{n+1}\binom{2n}{n}.
        \]
    \end{theorem}

    \begin{strategy}[Tail Switching for Intersecting Paths]
        When two directed lattice paths meet, exchange their tails after the
        first intersection. This produces a bijection with a pair of paths whose
        endpoints have been crossed. The method is a standard way to count or
        cancel intersecting path pairs.
    \end{strategy}

    \begin{formula}[Ordinary Generating Function Table]
        For a sequence $(a_n)_{n\geq 0}$, its ordinary generating function is
        \[
            A(x)=\sum_{n=0}^{\infty}a_nx^n.
        \]
        The basic examples are

        {\small
        \renewcommand{\arraystretch}{1.18}
        \begin{longtable}{@{}L{0.38\textwidth}L{0.50\textwidth}@{}}
            \toprule
            Sequence $a_n$ & Ordinary generating function $A(x)$ \\
            \midrule
            \endfirsthead

            \toprule
            Sequence $a_n$ & Ordinary generating function $A(x)$ \\
            \midrule
            \endhead

            \bottomrule
            \endlastfoot

            $1$ & $\displaystyle\frac{1}{1-x}$ \\
            $n$ & $\displaystyle\frac{x}{(1-x)^2}$ \\
            $n^2$ & $\displaystyle\frac{x(1+x)}{(1-x)^3}$ \\
            $\displaystyle\binom{n+r}{r}$ & $\displaystyle\frac{1}{(1-x)^{r+1}}$ \\
            $F_n$, where $F_0=0$ and $F_1=1$ & $\displaystyle\frac{x}{1-x-x^2}$ \\
            $\displaystyle\frac{1}{n+1}\binom{2n}{n}$ &
            $\displaystyle\frac{1-\sqrt{1-4x}}{2x}$ \\
        \end{longtable}
        }
    \end{formula}

    \begin{tactic}
        Generating functions are useful when a recurrence or counting problem has a
        repeated structure; a known generating function can replace a long
        derivation.
    \end{tactic}


    \begin{formula}[Stirling Numbers and Bell Numbers]
        The Stirling number of the second kind
        \[
            \stirlingsecond{n}{k}
        \]
        counts partitions of an $n$-element set into $k$ nonempty blocks and
        satisfies
        \[
            \stirlingsecond{n}{k}
            =
            k\stirlingsecond{n-1}{k}
            +
            \stirlingsecond{n-1}{k-1}.
        \]
        The Bell number
        \[
            B_n=\sum_{k=0}^{n}\stirlingsecond{n}{k}
        \]
        counts all set partitions of an $n$-element set.
    \end{formula}

    \begin{formula}[Fibonacci Numbers; Binet's Formula]
        If $F_0=0$, $F_1=1$, and $F_{n+1}=F_n+F_{n-1}$, then
        \[
            F_n
            =
            \frac{\phi^n-\psi^n}{\sqrt{5}},
            \qquad
            \phi=\frac{1+\sqrt{5}}{2},
            \qquad
            \psi=\frac{1-\sqrt{5}}{2}.
        \]
        In particular,
        \[
            \lim_{n\to\infty}\frac{F_{n+1}}{F_n}=\phi.
        \]
    \end{formula}

    \begin{formula}[Roots-of-Unity Filter]
        Let $\omega_m=e^{2\pi i/m}$. Then
        \[
            \frac{1}{m}\sum_{j=0}^{m-1}\omega_m^{j(n-r)}
            =
            \begin{cases}
                1, & n\equiv r\pmod m, \\
                0, & n\not\equiv r\pmod m.
            \end{cases}
        \]
        Consequently, if $A(x)=\sum_{n\geq0}a_nx^n$, then the terms with
        indices congruent to $r$ modulo $m$ are extracted by
        \[
            \sum_{n\equiv r\, (\mathrm{mod}\,m)}a_nx^n
            =
            \frac1m\sum_{j=0}^{m-1}\omega_m^{-rj}A(\omega_m^j x).
        \]
    \end{formula}

    \begin{formula}[Tower of Hanoi]
        Let $T_n$ be the minimum number of moves needed to move $n$ disks in the
        Tower of Hanoi puzzle. Then
        \[
            T_n=2T_{n-1}+1,
            \qquad
            T_1=1,
        \]
        and hence
        \[
            T_n=2^n-1.
        \]
    \end{formula}

    \begin{remark}
        The recurrence comes from moving the top $n-1$ disks away, moving the
        largest disk once, and then moving the $n-1$ disks back on top of it.
    \end{remark}


    \begin{strategy}[State Compression and Transfer Matrices]
        When a counting process depends only on finitely many recent states,
        place the transition counts in a matrix $\mathbf{T}$. Counts of length
        $n$ then have the form
        \[
            \mathbf{u}^{\mathsf T}\mathbf{T}^n\mathbf{v}.
        \]
        This is effective for forbidden substrings, tilings, walks, and finite
        automata.
    \end{strategy}

    \begin{theorem}[Pigeonhole Principle; Dirichlet's Box Principle]
        If $n+1$ objects are placed into $n$ boxes, then at least one box
        contains at least two objects.
        More generally, if $N$ objects are placed into $k$ boxes, then at least
        one box contains at least
        \[
            \left\lceil\frac{N}{k}\right\rceil
        \]
        objects.
    \end{theorem}

    \begin{strategy}
        The Pigeonhole Principle is often hidden in problems asking us to prove
        that two objects must share a property. The hard part is usually
        choosing the right ``boxes.''
    \end{strategy}

    \begin{theorem}[Consecutive-Sum Divisibility Lemma]
        For any integers $a_1,\dots,a_n$, there exists a nonempty consecutive
        block whose sum is divisible by $n$. Define prefix sums
        \[
            S_0=0,
            \qquad
            S_k=\sum_{i=1}^{k}a_i.
        \]
        If some $S_k\equiv0\pmod n$, use the first $k$ terms. Otherwise, two
        of $S_1,\dots,S_n$ have the same residue modulo $n$, and their difference
        is the required block sum.
    \end{theorem}

    \begin{strategy}[Coloring, Parity, and Bipartite Obstructions]
        Colorings convert geometric or movement constraints into parity data.
        On a bipartite graph, every path alternates colors. Thus a Hamiltonian
        path can exist only if the two color classes differ in size by at most
        $1$, and a Hamiltonian cycle requires the two classes to have equal size.
        Checkerboard colorings are especially effective for tiling and movement
        problems.
    \end{strategy}

    \begin{theorem}[Inclusion--Exclusion Principle; Sieve Formula]
        For finite sets $A_1,\dots,A_n$,
        \[
            \left|\bigcup_{i=1}^{n}A_i\right|
            =
            \sum_i |A_i|
            -
            \sum_{i<j}|A_i\cap A_j|
            +
            \sum_{i<j<k}|A_i\cap A_j\cap A_k|
            -\cdots
            +(-1)^{n+1}|A_1\cap\cdots\cap A_n|.
        \]
    \end{theorem}

    \begin{remark}
        Inclusion--Exclusion is the standard correction mechanism for
        overcounting. It is also historically connected to sieve methods.
    \end{remark}

    \begin{formula}[Derangement Formula]
        Let $D_n$ be the number of derangements of $n$ objects, that is,
        permutations with no fixed points. Then
        \[
            D_n
            =
            n!\sum_{k=0}^{n}\frac{(-1)^k}{k!}.
        \]
        Moreover,
        \[
            D_n
            =
            \left\lfloor \frac{n!}{e}+\frac12\right\rfloor.
        \]
    \end{formula}

    \begin{remark}
        The approximation $D_n\approx n!/e$ is very memorable. In probability
        language, a random permutation has no fixed point with probability
        approximately $1/e$.
    \end{remark}

    \begin{formula}[Catalan Numbers]
        The $n$th Catalan number is
        \[
            C_n
            =
            \frac{1}{n+1}\binom{2n}{n}.
        \]
        It begins as
        \[
            1,\ 1,\ 2,\ 5,\ 14,\ 42,\ 132,\dots.
        \]
    \end{formula}

    \begin{remark}
        Catalan numbers count many different objects, including correct
        parenthesizations, Dyck paths, triangulations of a polygon, and rooted
        full binary trees. The number $42$ is a classic giveaway.
    \end{remark}

    \begin{theorem}[Cayley's Formula]
        The number of labeled trees on $n$ vertices is
        \[
            n^{n-2}.
        \]
    \end{theorem}

    \begin{remark}
        Cayley's formula gives a compact and surprising enumeration of labeled
        trees.
    \end{remark}

    \begin{theorem}[Burnside's Lemma; Cauchy--Frobenius--Burnside Lemma]
        Suppose a finite group $G$ acts on a finite set $X$. The number of
        orbits is
        \[
            \frac{1}{|G|}
            \sum_{g\in G}|\Fix(g)|,
        \]
        where
        \[
            \Fix(g)
            =
            \{x\in X\mid g\cdot x=x\}.
        \]
    \end{theorem}

    \begin{idea}
        Burnside's Lemma says that the number of essentially different objects
        is the average number of objects fixed by a symmetry.
    \end{idea}

    \begin{lemma}[Handshaking Lemma]
        For a finite undirected graph $G=(V,E)$,
        \[
            \sum_{v\in V}\deg(v)=2|E|.
        \]
        In particular, the number of vertices of odd degree is even.
    \end{lemma}

    \begin{remark}
        The name comes from the idea that every handshake contributes $2$ to
        the total count of hands shaken.
    \end{remark}

    \begin{theorem}[Euler's Formula for Planar Graphs]
        If a connected planar graph has $V$ vertices, $E$ edges, and $F$ faces,
        then
        \[
            V-E+F=2.
        \]
    \end{theorem}

    \begin{remark}
        This is the graph-theoretic version of Euler's polyhedron formula.
        It is one of the most famous formulas involving planar graphs.
    \end{remark}

    \begin{formula}[Regions Cut by Hyperplanes in General Position]
        The maximum number of regions into which $n$ hyperplanes in general
        position divide $d$-dimensional space is
        \[
            R_d(n)=\sum_{k=0}^{d}\binom{n}{k}.
        \]
        In particular, $n$ lines divide the plane into at most
        \[
            1+n+\binom{n}{2}
        \]
        regions, and $n$ planes divide $\R^3$ into at most
        \[
            1+n+\binom{n}{2}+\binom{n}{3}
        \]
        regions.
    \end{formula}

    \begin{formula}[Chords and Diagonals in Convex Position]
        If $n$ points lie on a circle and no three interior chords are concurrent,
        drawing every chord divides the disk into
        \[
            1+\binom{n}{2}+\binom{n}{4}
        \]
        regions. A convex $n$-gon has at most
        \[
            \binom{n}{4}
        \]
        interior intersection points of diagonals.
    \end{formula}

    \begin{formula}[Triangles in a Complete Diagonal Drawing]
        For $n$ vertices in convex position with no three diagonals concurrent,
        the number of triangles whose sides lie on sides or diagonals of the
        complete drawing is
        \[
            \binom{n}{3}
            +4\binom{n}{4}
            +5\binom{n}{5}
            +\binom{n}{6}.
        \]
    \end{formula}

    \begin{formula}[Triangulation with Interior Vertices]
        Suppose a polygonal disk is triangulated using $V$ vertices in total,
        of which $b$ lie on the boundary. Then the number of edges and triangular
        faces are
        \[
            E=3V-b-3,
            \qquad
            F_{\triangle}=2V-b-2.
        \]
        These values follow from Euler's formula and double counting the
        triangle-edge incidences.
    \end{formula}

    \begin{strategy}[Double Counting with Euler's Formula]
        In a planar subdivision, count edge-face incidences first and then use
        $V-E+F=2$. Interior edges are counted twice and boundary edges once.
        This is often faster and safer than trying to enumerate faces directly.
    \end{strategy}

    \begin{theorem}[Eulerian Trail Criterion]
        A connected finite graph has an Eulerian circuit if and only if every
        vertex has even degree. It has an Eulerian trail but not an Eulerian
        circuit if and only if exactly two vertices have odd degree.
    \end{theorem}

    \begin{remark}
        This theorem is historically connected to Euler's solution of the
        Seven Bridges of K\"onigsberg problem, which marks the beginning of
        graph theory and topology. The historical city K\"onigsberg is now
        Kaliningrad, Russia.
    \end{remark}

    \begin{theorem}[Hall's Marriage Theorem]
        Let $G=(X\cup Y,E)$ be a bipartite graph. There is a matching that
        covers every vertex of $X$ if and only if, for every subset
        $S\subseteq X$,
        \[
            |N(S)|\geq |S|,
        \]
        where $N(S)$ is the set of neighbors of vertices in $S$.
    \end{theorem}

    \begin{strategy}
        Hall's Marriage Theorem is the standard named theorem for proving that
        a perfect assignment or matching exists.
    \end{strategy}


    \begin{theorem}[K\H{o}nig's Theorem]
        In every finite bipartite graph, the size of a maximum matching equals
        the size of a minimum vertex cover.
    \end{theorem}

    \begin{theorem}[Max-Flow Min-Cut Theorem]
        In a finite capacitated network, the maximum value of a source--sink
        flow equals the minimum capacity of a source--sink cut.
    \end{theorem}

    \begin{theorem}[Menger's Theorem, Vertex Form]
        For nonadjacent vertices $u$ and $v$, the maximum number of pairwise
        internally vertex-disjoint $u$--$v$ paths equals the minimum number of
        vertices whose deletion separates $u$ from $v$.
    \end{theorem}

    \begin{theorem}[Kirchhoff's Matrix--Tree Theorem]
        Let
        \[
            \mathbf{L}=\mathbf{D}-\mathbf{A}
        \]
        be the Laplacian matrix of a finite graph. The number of spanning trees
        equals any cofactor of $\mathbf{L}$ obtained by deleting the same row
        and column.
    \end{theorem}

    \begin{theorem}[Sperner's Theorem]
        The largest family of subsets of $[n]$ in which no member contains
        another has size
        \[
            \binom{n}{\lfloor n/2\rfloor}.
        \]
    \end{theorem}

    \begin{theorem}[Dilworth's Theorem]
        In a finite partially ordered set, the maximum size of an antichain
        equals the minimum number of chains covering the poset.
    \end{theorem}

    \begin{theorem}[Tur\'an's Theorem]
        Among all $n$-vertex graphs containing no $K_{r+1}$, the maximum number
        of edges is attained by the complete $r$-partite graph whose part sizes
        differ by at most $1$.
    \end{theorem}

    \begin{theorem}[Dirac's and Ore's Hamiltonian Criteria]
        Let $G$ be a simple graph on $n\geq3$ vertices. If every vertex has
        degree at least $n/2$, then $G$ is Hamiltonian. More generally, it is
        enough that
        \[
            \deg(u)+\deg(v)\geq n
        \]
        for every pair of nonadjacent vertices $u,v$.
    \end{theorem}

    \begin{theorem}[Ramsey's Theorem]
        For any positive integers $r$ and $s$, there exists an integer
        $R(r,s)$ such that every red-blue coloring of the edges of a sufficiently
        large complete graph contains either a red $K_r$ or a blue $K_s$.
    \end{theorem}

    \begin{formula}[The Party Problem]
        The classical Ramsey number
        \[
            R(3,3)=6
        \]
        says that among any six people, there are either three mutual
        acquaintances or three mutual strangers.
    \end{formula}

    \begin{remark}
        Ramsey theory is the mathematics of unavoidable structure. Its slogan
        is that complete disorder is impossible in a large enough system.
    \end{remark}

    \begin{theorem}[Erd\H{o}s--Szekeres Theorem]
        Any sequence of
        \[
            (r-1)(s-1)+1
        \]
        distinct real numbers contains either an increasing subsequence of
        length $r$ or a decreasing subsequence of length $s$.
    \end{theorem}

    \begin{remark}
        This is a beautiful pigeonhole-style result. It is often remembered as
        the monotone subsequence theorem.
    \end{remark}

    \chapter{Probability and Statistics}

    \relatedcourses{(MATH230) Probability \& Statistics}

    \relatedbooks{Introduction to Probability
        (Dimitri P. Bertsekas, John N. Tsitsiklis)}

    \begin{theorem}[Law of Total Probability]
        If $B_1,\dots,B_n$ form a partition of the sample space and
        $\Prob(B_i)>0$ for all $i$, then
        \[
            \Prob(A)
            =
            \sum_{i=1}^{n}\Prob(A\mid B_i)\Prob(B_i).
        \]
    \end{theorem}


    \begin{theorem}[Law of Total Expectation; Tower Property]
        For a discrete random variable $Y$,
        \[
            \Exp[X]
            =
            \sum_y\Exp[X\mid Y=y]\Prob(Y=y)
            =
            \Exp\!\left[\Exp[X\mid Y]\right].
        \]
    \end{theorem}

    \begin{formula}[Law of Total Variance]
        \[
            \Var(X)
            =
            \Exp\!\left[\Var(X\mid Y)\right]
            +
            \Var\!\left(\Exp[X\mid Y]\right).
        \]
    \end{formula}

    \begin{theorem}[Bayes' Theorem]
        If $B_1,\dots,B_n$ form a partition of the sample space and
        $\Prob(A)>0$, then
        \[
            \Prob(B_j\mid A)
            =
            \frac{\Prob(A\mid B_j)\Prob(B_j)}
                 {\sum_{i=1}^{n}\Prob(A\mid B_i)\Prob(B_i)}.
        \]
        In the two-event form,
        \[
            \Prob(B\mid A)
            =
            \frac{\Prob(A\mid B)\Prob(B)}{\Prob(A)}.
        \]
    \end{theorem}

    \begin{remark}
        Bayes' Theorem is the standard formula for reversing conditional
        probability. It updates a prior probability after observing evidence.
    \end{remark}

    \begin{definition}[Independence]
        Events $A$ and $B$ are independent if
        \[
            \Prob(A\cap B)=\Prob(A)\Prob(B).
        \]
        Equivalently, if $\Prob(B)>0$, then
        \[
            \Prob(A\mid B)=\Prob(A).
        \]
        Random variables $X$ and $Y$ are independent if events determined by
        $X$ and events determined by $Y$ are independent.
    \end{definition}

    \begin{tactic}
        ``Disjoint'' and ``independent'' are different. If two nonempty events
        are disjoint, then knowing that one occurred makes the other impossible;
        they are usually not independent.
    \end{tactic}

    \begin{formula}[Complement Rule and Inclusion--Exclusion for Probability]
        The complement rule is
        \[
            \Prob(A^{\mathsf c})=1-\Prob(A).
        \]
        For two events,
        \[
            \Prob(A\cup B)
            =
            \Prob(A)+\Prob(B)-\Prob(A\cap B).
        \]
    \end{formula}

    \begin{proposition}[Linearity of Expectation]
        For random variables $X_1,\dots,X_n$,
        \[
            \Exp[X_1+\cdots+X_n]
            =
            \Exp[X_1]+\cdots+\Exp[X_n].
        \]
        This holds even when the random variables are not independent.
    \end{proposition}

    \begin{strategy}
        Linearity of expectation is often the fastest way to count an expected
        number of objects. Introduce indicator random variables and add their
        expectations.
    \end{strategy}

    \begin{proposition}[Expectation of a Product of Independent Random Variables]
        If $X_1,X_2,\dots,X_n$ are independent random variables, then
        \[
            \Exp[X_1X_2\cdots X_n]
            =
            \Exp[X_1]\Exp[X_2]\cdots\Exp[X_n].
        \]
        In particular, if $X_1,\dots,X_n$ are independent and identically
        distributed, then
        \[
            \Exp[X_1X_2\cdots X_n]
            =
            (\Exp[X_1])^n.
        \]
    \end{proposition}

    \begin{formula}[Common Probability Distributions]
        The most frequently used distributions are summarized below.

        {\small
        \renewcommand{\arraystretch}{1.18}
        \begin{longtable}{@{}L{0.19\textwidth}L{0.35\textwidth}L{0.17\textwidth}L{0.17\textwidth}@{}}
            \toprule
            Distribution & Mass or density function & Mean & Variance \\
            \midrule
            \endfirsthead

            \toprule
            Distribution & Mass or density function & Mean & Variance \\
            \midrule
            \endhead

            \bottomrule
            \endlastfoot

            Bernoulli$(p)$ &
            $\Prob(X=x)=p^x(1-p)^{1-x}$, $x\in\{0,1\}$ &
            $p$ & $p(1-p)$ \\
            Binomial$(n,p)$ &
            $\displaystyle\Prob(X=k)=\binom{n}{k}p^k(1-p)^{n-k}$ &
            $np$ & $np(1-p)$ \\
            Geometric$(p)$ &
            $\Prob(X=k)=(1-p)^{k-1}p$, $k\geq1$ &
            $\displaystyle\frac1p$ & $\displaystyle\frac{1-p}{p^2}$ \\
            Poisson$(\lambda)$ &
            $\displaystyle\Prob(X=k)=e^{-\lambda}\frac{\lambda^k}{k!}$ &
            $\lambda$ & $\lambda$ \\
            Uniform$(a,b)$ &
            $\displaystyle f(x)=\frac{1}{b-a}$, $a\leq x\leq b$ &
            $\displaystyle\frac{a+b}{2}$ & $\displaystyle\frac{(b-a)^2}{12}$ \\
            Exponential$(\lambda)$ &
            $f(x)=\lambda e^{-\lambda x}$, $x\geq0$ &
            $\displaystyle\frac1\lambda$ & $\displaystyle\frac1{\lambda^2}$ \\
            Normal$(\mu,\sigma^2)$ &
            $\displaystyle f(x)=\frac{1}{\sigma\sqrt{2\pi}}
            e^{-(x-\mu)^2/(2\sigma^2)}$ &
            $\mu$ & $\sigma^2$ \\
        \end{longtable}
        }
    \end{formula}

    \begin{formula}[Geometric Distribution and Waiting Time]
        If $T$ is the number of independent trials needed to obtain the first
        success, and each trial has success probability $p$, then
        \[
            \Prob(T=k)=(1-p)^{k-1}p,
            \qquad
            k=1,2,3,\dots,
        \]
        and
        \[
            \Exp[T]=\frac{1}{p}.
        \]
    \end{formula}

    \begin{strategy}
        Waiting-until questions usually start with a geometric distribution. If
        a running sum is also involved, combine the expected number of trials
        with the expected contribution per trial, or condition on the first
        step.
    \end{strategy}

    \begin{formula}[Variance and Covariance Identities]
        For a random variable $X$,
        \[
            \Var(X)
            =
            \Exp[X^2]-(\Exp[X])^2.
        \]
        For random variables $X$ and $Y$,
        \[
            \Var(X+Y)
            =
            \Var(X)+\Var(Y)+2\Cov(X,Y).
        \]
        If $X$ and $Y$ are independent, then
        \[
            \Var(X+Y)=\Var(X)+\Var(Y).
        \]
    \end{formula}

    \begin{formula}[Binomial Distribution]
        If $X\sim\Bin(n,p)$, then
        \[
            \Prob(X=k)
            =
            \binom{n}{k}p^k(1-p)^{n-k},
            \qquad
            k=0,1,\dots,n.
        \]
        Its mean and variance are
        \[
            \Exp[X]=np,
            \qquad
            \Var(X)=np(1-p).
        \]
    \end{formula}

    \begin{formula}[Hypergeometric Distribution]
        Suppose a population has $N$ objects, of which $K$ are successes. If
        $n$ objects are chosen without replacement and $X$ is the number of
        successes chosen, then
        \[
            \Prob(X=k)
            =
            \frac{\binom{K}{k}\binom{N-K}{n-k}}{\binom{N}{n}}.
        \]
        Its mean is
        \[
            \Exp[X]=n\frac{K}{N}.
        \]
    \end{formula}

    \begin{formula}[Poisson Distribution]
        If $X\sim\Poi(\lambda)$, then
        \[
            \Prob(X=k)
            =
            e^{-\lambda}\frac{\lambda^k}{k!},
            \qquad
            k=0,1,2,\dots.
        \]
        Its mean and variance are
        \[
            \Exp[X]=\lambda,
            \qquad
            \Var(X)=\lambda.
        \]
    \end{formula}

    \begin{theorem}[Poisson Limit Theorem; Law of Small Numbers]
        If $n\to\infty$, $p\to 0$, and $np\to\lambda$, then
        \[
            \binom{n}{k}p^k(1-p)^{n-k}
            \to
            e^{-\lambda}\frac{\lambda^k}{k!}
        \]
        for each fixed $k\geq 0$.
    \end{theorem}

    \begin{remark}
        This theorem explains why the Poisson distribution models rare events:
        many trials, small success probability, and finite expected count.
    \end{remark}

    \begin{formula}[Normal Distribution; Gaussian Distribution]
        A random variable $X\sim\Normal(\mu,\sigma^2)$ has density
        \[
            f(x)
            =
            \frac{1}{\sigma\sqrt{2\pi}}
            \exp\left(
                -\frac{(x-\mu)^2}{2\sigma^2}
            \right).
        \]
        The standard normal distribution is $\Normal(0,1)$.
    \end{formula}

    \begin{formula}[Empirical Rule; 68--95--99.7 Rule]
        For a normal distribution, approximately
        \[
            68\%
        \]
        of the mass lies within one standard deviation of the mean,
        \[
            95\%
        \]
        lies within two standard deviations, and
        \[
            99.7\%
        \]
        lies within three standard deviations.
    \end{formula}

    \begin{definition}[Modes of Convergence for Random Variables]
        Let $X_n$ and $X$ be random variables on the same probability space.
        The principal modes of convergence used in probability are:
        \begin{enumerate}
            \item[(i)] $X_n\asto X$ if
            \[
                \Prob\bigl(\{\omega\mid X_n(\omega)\to X(\omega)\}\bigr)=1;
            \]
            \item[(ii)] $X_n\probto X$ if
            \[
                \forall\eps>0,
                \qquad
                \Prob\bigl(|X_n-X|>\eps\bigr)\to 0;
            \]
            \item[(iii)] $X_n\Lpto{p}X$ for $p\geq 1$ if
            \[
                \Exp\bigl[|X_n-X|^p\bigr]\to 0;
            \]
            \item[(iv)] $X_n\distto X$ if
            \[
                F_{X_n}(x)\to F_X(x)
            \]
            at every continuity point $x$ of the distribution function $F_X$.
        \end{enumerate}
        Almost-sure convergence implies convergence in probability, and
        convergence in probability implies convergence in distribution.
    \end{definition}


    \begin{theorem}[Borel--Cantelli Lemmas]
        If
        \[
            \sum_{n=1}^{\infty}\Prob(A_n)<\infty,
        \]
        then $\Prob(A_n\text{ infinitely often})=0$. If the $A_n$ are
        independent and the series diverges, then
        \[
            \Prob(A_n\text{ infinitely often})=1.
        \]
    \end{theorem}

    \begin{theorem}[Weak Law of Large Numbers]
        Let $X_1,X_2,\dots$ be independent identically distributed random
        variables with mean $\mu$. Then, for every $\eps>0$,
        \[
            \Prob\left(
                \left|
                    \frac{X_1+\cdots+X_n}{n}-\mu
                \right|
                >
                \eps
            \right)
            \to 0
        \]
        as $n\to\infty$.
    \end{theorem}

    \begin{idea}
        The Law of Large Numbers says that the sample average stabilizes near
        the true mean when the number of trials becomes large.
    \end{idea}

    \begin{theorem}[Central Limit Theorem]
        Let $X_1,X_2,\dots$ be independent identically distributed random
        variables with mean $\mu$ and variance $\sigma^2>0$. Then
        \[
            \frac{X_1+\cdots+X_n-n\mu}{\sigma\sqrt n}
            \distto
            Z,
            \qquad
            Z\sim\Normal(0,1).
        \]
        Thus the standardized sum converges in distribution to the standard
        normal distribution.
    \end{theorem}

    \begin{remark}
        The Central Limit Theorem explains why the normal distribution appears
        so often in nature and statistics.
    \end{remark}

    \begin{theorem}[Markov's Inequality]
        If $X$ is a nonnegative random variable and $a>0$, then
        \[
            \Prob(X\geq a)
            \leq
            \frac{\Exp[X]}{a}.
        \]
    \end{theorem}

    \begin{theorem}[Chebyshev's Inequality]
        If $X$ has mean $\mu$ and variance $\sigma^2$, then for every $k>0$,
        \[
            \Prob(|X-\mu|\geq k\sigma)
            \leq
            \frac{1}{k^2}.
        \]
        Equivalently, for every $\eps>0$,
        \[
            \Prob(|X-\mu|\geq \eps)
            \leq
            \frac{\Var(X)}{\eps^2}.
        \]
    \end{theorem}

    \begin{remark}
        Markov's inequality uses only the mean, while Chebyshev's inequality
        uses the variance. Both are probability bounds that work without
        assuming a normal distribution.
    \end{remark}


    \begin{theorem}[Chernoff Bound for a Binomial Random Variable]
        If $X\sim\Bin(n,p)$ and $\mu=np$, then for $\delta>0$,
        \[
            \Prob\bigl(X\geq(1+\delta)\mu\bigr)
            \leq
            \left(\frac{e^\delta}{(1+\delta)^{1+\delta}}\right)^\mu.
        \]
        For $0<\delta<1$,
        \[
            \Prob\bigl(|X-\mu|\geq\delta\mu\bigr)
            \leq
            2e^{-\delta^2\mu/3}.
        \]
    \end{theorem}

    \begin{theorem}[Jensen's Inequality]
        If $\varphi$ is convex, then
        \[
            \varphi(\Exp[X])
            \leq
            \Exp[\varphi(X)].
        \]
        If $\varphi$ is concave, the inequality is reversed.
    \end{theorem}

    \begin{idea}
        Jensen's inequality says that, for a convex function, applying the
        function after averaging gives a value no larger than averaging after
        applying the function.
    \end{idea}

    \begin{formula}[Exact Birthday Problem]
        In a group of $n$ people, assuming $365$ equally likely birthdays and
        ignoring leap years, the probability that at least two people share a
        birthday is
        \[
            1-
            \frac{365\cdot 364\cdot 363\cdots(365-n+1)}{365^n}
            =
            1-
            \prod_{k=0}^{n-1}\frac{365-k}{365}.
        \]
        Equivalently, the probability that all birthdays are distinct is
        \[
            \prod_{k=0}^{n-1}\frac{365-k}{365}.
        \]
    \end{formula}

    \begin{formula}[Birthday Problem Approximation]
        In a group of $n$ people, assuming $365$ equally likely birthdays, the
        probability that at least two people share a birthday is approximately
        \[
            1-\exp\left(-\frac{n(n-1)}{730}\right).
        \]
        For $n=23$, this probability is already greater than $1/2$.
    \end{formula}

    \begin{remark}
        The birthday problem is famous because the answer is surprisingly
        small: only $23$ people are needed for a better-than-even chance of a
        shared birthday.
    \end{remark}

    \begin{formula}[Expected Number of Runs in Coin Tosses]
        In $n$ tosses of a fair coin, a run is a maximal consecutive block of
        identical outcomes. If $R_n$ is the number of runs, then
        \[
            \Exp[R_n]
            =
            1+\frac{n-1}{2}
            =
            \frac{n+1}{2}.
        \]
    \end{formula}

    \begin{idea}
        The first toss starts one run. For each of the remaining $n-1$ positions,
        a new run begins exactly when the current toss differs from the previous
        toss, which has probability $1/2$.
    \end{idea}

    \begin{formula}[Monty Hall Problem]
        In the standard three-door Monty Hall problem, the probability of
        winning by staying with the original choice is
        \[
            \frac{1}{3},
        \]
        while the probability of winning by switching is
        \[
            \frac{2}{3}.
        \]
    \end{formula}

    \begin{remark}
        The host's action is not random information: the host always reveals a
        losing door. This is why the original $2/3$ probability of having chosen
        incorrectly transfers to the remaining unopened door.
    \end{remark}

    \begin{formula}[Coupon Collector]
        If there are $n$ equally likely coupon types, the expected number of
        trials needed to see all $n$ types is
        \[
            nH_n
            =
            n\left(1+\frac12+\cdots+\frac1n\right).
        \]
    \end{formula}

    \begin{strategy}
        For a fair die, the expected number of rolls needed to see all six faces
        is $6H_6$. If a question asks about the order in which the first
        appearances occur, ignore repeated rolls and look only at the random
        ordering of the six faces.
    \end{strategy}

    \begin{formula}[Gambler's Ruin]
        Suppose a gambler starts with $i$ dollars and stops when reaching
        either $0$ dollars or $N$ dollars. If the probability of winning each
        round is $p$ and $q=1-p$, then the probability of reaching $N$ before
        $0$ is
        \[
            \frac{1-(q/p)^i}{1-(q/p)^N}
            \qquad
            (p\neq q).
        \]
        In the fair case $p=q=1/2$, the probability is
        \[
            \frac{i}{N}.
        \]
    \end{formula}

    \begin{remark}
        Gambler's ruin is a classic random walk result. It is a good example
        where the fair case has a very simple answer.
    \end{remark}

    \chapter{Geometry}

    \relatedcourses{(MATH428) Modern Geometry,
        (MATH570/CSED508) Discrete and Computational Geometry}

    \relatedbooks{Geometry Revisited (H. S. M. Coxeter, Samuel L. Greitzer),
        Euclidean Geometry in Mathematical Olympiads (Evan Chen)}

    \section*{Synthetic Geometry}

    \begin{formula}[Basic Plane and Solid Geometry Table]
        The basic measurement formulas are

        {\small
        \renewcommand{\arraystretch}{1.18}
        \begin{longtable}{@{}L{0.27\textwidth}L{0.28\textwidth}L{0.35\textwidth}@{}}
            \toprule
            Figure & Perimeter or surface area & Area or volume \\
            \midrule
            \endfirsthead

            \toprule
            Figure & Perimeter or surface area & Area or volume \\
            \midrule
            \endhead

            \bottomrule
            \endlastfoot

            Triangle & $a+b+c$ & $\displaystyle\frac12 bh$ \\
            Parallelogram & $2(a+b)$ & $bh$ \\
            Trapezoid & $a+b+c+d$ & $\displaystyle\frac12(a+c)h$ \\
            Regular $n$-gon & $na$ & $\displaystyle\frac12(na)r=\frac{na^2}{4\tan(\pi/n)}$ \\
            Circle & $2\pi r$ & $\pi r^2$ \\
            Circular sector & $r\theta$ & $\displaystyle\frac12r^2\theta$ \\
            Rectangular box & $2(ab+bc+ca)$ & $abc$ \\
            Right circular cylinder & $2\pi r(r+h)$ & $\pi r^2h$ \\
            Right circular cone & $\pi r(r+\ell)$ & $\displaystyle\frac13\pi r^2h$ \\
            Sphere & $4\pi r^2$ & $\displaystyle\frac43\pi r^3$ \\
        \end{longtable}
        }
        In the regular-polygon row, $r$ is the apothem; in the cone row,
        $\ell=\sqrt{r^2+h^2}$ is the slant height. Sector angles are measured in
        radians.
    \end{formula}

    \begin{theorem}[Triangle Angle Sum and Exterior Angle Theorem]
        For every triangle $\triangle ABC$,
        \[
            \angle A+\angle B+\angle C=\pi.
        \]
        An exterior angle equals the sum of the two nonadjacent interior angles.
    \end{theorem}

    \begin{proposition}[Triangle Congruence and Similarity Criteria]
        Two triangles are congruent under any of the criteria
        \[
            \mathrm{SSS},\qquad \mathrm{SAS},\qquad \mathrm{ASA},
            \qquad \mathrm{AAS},\qquad \mathrm{RHS}.
        \]
        They are similar under any of the criteria
        \[
            \mathrm{AA},\qquad \mathrm{SAS},\qquad \mathrm{SSS},
        \]
        where the latter two use proportional corresponding side lengths.
    \end{proposition}

    \begin{theorem}[Midpoint Theorem and Centroid Ratio]
        The segment joining the midpoints of two sides of a triangle is parallel
        to the third side and has half its length. The three medians are
        concurrent at the centroid $G$, and $G$ divides every median in the ratio
        \[
            2:1
        \]
        measured from the vertex.
    \end{theorem}

    \begin{strategy}[Reverse Engineering in Geometric Constructions]
        First draw a completed configuration satisfying the desired conditions.
        Then identify which points, perpendiculars, parallels, angle bisectors,
        circles, or reflections in the completed picture can be constructed from
        the original data. Familiar configurations involving triangle centers,
        similar triangles, and power of a point are especially useful targets.
    \end{strategy}


    \begin{theorem}[Happy Ending Problem, Five-Point Case]
        Among any five points in the plane, no three collinear, there are four
        points that are the vertices of a convex quadrilateral.
    \end{theorem}

    \begin{idea}
        Examine the convex hull. If it has at least four vertices, choose four
        of them. If it is a triangle, use the line through the two interior
        points and choose two hull vertices lying on the same side of that line.
    \end{idea}

    \begin{definition}[Classical Triangle Centers]
        For a triangle $\triangle ABC$, the most important classical centers
        are:
        \[
            O=\text{circumcenter},
            \qquad
            I=\text{incenter},
            \qquad
            H=\text{orthocenter},
            \qquad
            G=\text{centroid}.
        \]
        Here $O$ is the center of the circumcircle, $I$ is the center of the
        incircle, $H$ is the intersection of the altitudes, and $G$ is the
        intersection of the medians.
    \end{definition}

    \begin{theorem}[Gergonne Point and Nagel Point]
        The cevians joining the vertices of a triangle to the contact points of
        its incircle with the opposite sides are concurrent at the Gergonne point.
        The cevians joining the vertices to the contact points of the opposite
        excircles are concurrent at the Nagel point.
    \end{theorem}

    \begin{theorem}[Symmedian Lemma and Lemoine Point]
        Let the $A$-symmedian of $\triangle ABC$ meet $BC$ at $K$. Then
        \[
            \frac{BK}{KC}=\frac{AB^2}{AC^2}.
        \]
        The three symmedians are concurrent at the Lemoine point, also called
        the symmedian point.
    \end{theorem}

    \begin{theorem}[Fermat Point; Torricelli Point]
        If every angle of $\triangle ABC$ is less than $120^\circ$, there is a
        unique point $P$ minimizing
        \[
            PA+PB+PC.
        \]
        It satisfies
        \[
            \angle APB=\angle BPC=\angle CPA=120^\circ.
        \]
        If one angle is at least $120^\circ$, the minimizing point is that vertex.
    \end{theorem}

    \begin{formula}[Extended Law of Sines]
        For a triangle $\triangle ABC$ with side lengths
        \[
            a=BC,\qquad b=CA,\qquad c=AB,
        \]
        and circumradius $R$, we have
        \[
            \frac{a}{\sin A}
            =
            \frac{b}{\sin B}
            =
            \frac{c}{\sin C}
            =
            2R.
        \]
    \end{formula}

    \begin{formula}[Law of Cosines]
        For a triangle $\triangle ABC$ with side lengths
        \[
            a=BC,\qquad b=CA,\qquad c=AB,
        \]
        we have
        \[
            c^2=a^2+b^2-2ab\cos C.
        \]
        The analogous formulas hold after cyclic permutation of $a,b,c$.
    \end{formula}

    \begin{formula}[Heron's Formula]
        Let $\triangle ABC$ have side lengths $a,b,c$ and semiperimeter
        \[
            s=\frac{a+b+c}{2}.
        \]
        Then its area $K$ is
        \[
            K
            =
            \sqrt{s(s-a)(s-b)(s-c)}.
        \]
        Also,
        \[
            K=rs
            \qquad\text{and}\qquad
            K=\frac{abc}{4R},
        \]
        where $r$ is the inradius and $R$ is the circumradius.
    \end{formula}

    \begin{strategy}
        The formulas
        \[
            K=rs,
            \qquad
            K=\frac{abc}{4R}
        \]
        are often more useful than Heron's formula when incircles or
        circumcircles appear.
    \end{strategy}

    \begin{formula}[Triangle Area Formulas]
        For a triangle with side lengths $a,b,c$, angles $A,B,C$, circumradius
        $R$, and inradius $r$,
        \[
            [ABC]
            =
            \frac12bc\sin A
            =
            \frac12ca\sin B
            =
            \frac12ab\sin C,
        \]
        and
        \[
            [ABC]=rs=\frac{abc}{4R}.
        \]
    \end{formula}


    \begin{theorem}[Weitzenb\"ock's Inequality]
        For a triangle with side lengths $a,b,c$ and area $\Delta$,
        \[
            a^2+b^2+c^2\geq4\sqrt{3}\,\Delta,
        \]
        with equality if and only if the triangle is equilateral.
    \end{theorem}

    \begin{formula}[Inradius and Exradii]
        For a triangle of area $K$ and semiperimeter $s$, the inradius is
        \[
            r=\frac{K}{s}.
        \]
        If $r_a,r_b,r_c$ are the exradii opposite $A,B,C$, respectively, then
        \[
            r_a=\frac{K}{s-a},
            \qquad
            r_b=\frac{K}{s-b},
            \qquad
            r_c=\frac{K}{s-c}.
        \]
    \end{formula}

    \begin{formula}[Median Length Formula; Apollonius' Theorem]
        If $m_a$ is the median from $A$ to side $a=BC$, then
        \[
            m_a^2
            =
            \frac{2b^2+2c^2-a^2}{4}.
        \]
        Equivalently, if $M$ is the midpoint of $BC$, then
        \[
            AB^2+AC^2=2(AM^2+BM^2).
        \]
    \end{formula}

    \begin{theorem}[Stewart's Theorem]
        Let $D$ lie on side $BC$ of a triangle $\triangle ABC$. Write
        \[
            a=BC,
            \qquad
            b=CA,
            \qquad
            c=AB,
            \qquad
            m=BD,
            \qquad
            n=DC,
            \qquad
            d=AD,
        \]
        so that $a=m+n$. Then
        \[
            b^2m+c^2n
            =
            a(d^2+mn).
        \]
        Equivalently,
        \[
            d^2
            =
            \frac{b^2m+c^2n}{a}-mn.
        \]
    \end{theorem}

    \begin{remark}
        Stewart's Theorem is the standard general formula for the length of a
        cevian. The median-length formula is the special case $m=n=a/2$.
    \end{remark}

    \begin{formula}[Internal Angle-Bisector Length]
        Let the internal angle bisector from $A$ meet $BC$ at $D$, and denote
        its length by $\ell_a=AD$. For
        \[
            a=BC,
            \qquad
            b=CA,
            \qquad
            c=AB,
        \]
        we have
        \[
            \ell_a^2
            =
            bc\left(1-\frac{a^2}{(b+c)^2}\right)
            =
            \frac{4bc\,s(s-a)}{(b+c)^2},
        \]
        where $s=(a+b+c)/2$.
    \end{formula}

    \begin{theorem}[Angle Bisector Theorem]
        In a triangle $\triangle ABC$, suppose the internal angle bisector of
        $\angle A$ meets $BC$ at $D$. Then
        \[
            \frac{BD}{DC}
            =
            \frac{AB}{AC}.
        \]
    \end{theorem}

    \begin{theorem}[Ceva's Theorem]
        Let $D,E,F$ lie on the lines $BC,CA,AB$, respectively. Then the cevians
        $AD,BE,CF$ are concurrent if and only if
        \[
            \frac{BD}{DC}
            \cdot
            \frac{CE}{EA}
            \cdot
            \frac{AF}{FB}
            =
            1,
        \]
        using directed lengths.
    \end{theorem}

    \begin{remark}
        Ceva's Theorem is the standard theorem for proving concurrence of three
        lines in a triangle.
    \end{remark}

    \begin{theorem}[Trigonometric Form of Ceva's Theorem]
        Let $D,E,F$ lie on the lines $BC,CA,AB$, respectively. The cevians
        $AD,BE,CF$ are concurrent if and only if
        \[
            \frac{\sin\angle BAD}{\sin\angle DAC}
            \cdot
            \frac{\sin\angle CBE}{\sin\angle EBA}
            \cdot
            \frac{\sin\angle ACF}{\sin\angle FCB}
            =1,
        \]
        with the usual directed-angle convention when points lie on extended
        sides.
    \end{theorem}

    \begin{strategy}[Mass Points]
        In a cevian configuration, assign masses to vertices so that a point
        dividing a segment balances the endpoint masses inversely to the segment
        lengths. For example, if $D\in BC$, then
        \[
            \frac{BD}{DC}=\frac{m_C}{m_B}.
        \]
        Consistent masses convert several ratio computations into weighted
        averages and often recover Ceva-type relations immediately.
    \end{strategy}

    \begin{theorem}[Menelaus' Theorem]
        Let $D,E,F$ lie on the lines $BC,CA,AB$, respectively. Then the points
        $D,E,F$ are collinear if and only if
        \[
            \frac{BD}{DC}
            \cdot
            \frac{CE}{EA}
            \cdot
            \frac{AF}{FB}
            =
            -1,
        \]
        using directed lengths.
    \end{theorem}

    \begin{remark}
        Ceva is for concurrence, while Menelaus is for collinearity. This pair
        is one of the most important toolkits in olympiad geometry.
    \end{remark}

    \begin{theorem}[Euler Line Theorem]
        In any non-equilateral triangle, the circumcenter $O$, centroid $G$,
        and orthocenter $H$ are collinear. Moreover,
        \[
            OG:GH=1:2.
        \]
        The line passing through these points is called the Euler line.
    \end{theorem}

    \begin{theorem}[Nine-Point Circle Theorem; Euler Circle]
        In any triangle, the following nine points lie on one circle:
        the three side midpoints, the three feet of the altitudes, and the
        three midpoints between the orthocenter and the vertices. This circle
        is called the nine-point circle.
    \end{theorem}

    \begin{remark}
        The center of the nine-point circle is the midpoint of the segment
        joining the circumcenter $O$ and the orthocenter $H$.
    \end{remark}

    \begin{theorem}[Feuerbach's Theorem]
        The nine-point circle of a triangle is internally tangent to the incircle
        and externally tangent to each of the three excircles.
    \end{theorem}

    \begin{theorem}[Euler's Formula for the Incenter and Circumcenter]
        For a triangle with circumradius $R$, inradius $r$, circumcenter $O$,
        and incenter $I$, we have
        \[
            OI^2=R^2-2Rr.
        \]
        In particular,
        \[
            R\geq 2r.
        \]
    \end{theorem}

    \begin{remark}
        The inequality $R\geq 2r$ is called Euler's inequality for triangles.
        Equality holds exactly for an equilateral triangle.
    \end{remark}


    \begin{strategy}[Spiral Similarity]
        When two segments must be compared through both an angle and a ratio,
        look for a center of spiral similarity. It often reveals the required
        similar triangles immediately.
    \end{strategy}

    \begin{theorem}[Inscribed Angle Theorem]
        An inscribed angle equals half the central angle subtending the same arc.
        Consequently, angles subtending the same chord are equal.
    \end{theorem}

    \begin{theorem}[Cyclic Quadrilateral Criteria]
        For four distinct points $A,B,C,D$ in the plane, the following standard
        conditions are equivalent whenever the expressions are defined:
        \[
            A,B,C,D\text{ are concyclic},
        \]
        \[
            \angle ABC+\angle ADC=\pi,
        \]
        and
        \[
            \angle BAC=\angle BDC.
        \]
        Directed angles modulo $\pi$ remove case distinctions.
    \end{theorem}

    \begin{theorem}[Tangent--Chord Theorem; Alternate Segment Theorem]
        The angle between a tangent to a circle and a chord through the point of
        tangency equals the inscribed angle subtending that chord in the opposite
        arc.
    \end{theorem}

    \begin{theorem}[Power of a Point Theorem]
        Let $P$ be a point and let a line through $P$ meet a circle at
        $A$ and $B$. The product
        \[
            PA\cdot PB
        \]
        is independent of the choice of the line through $P$. It is called the
        power of $P$ with respect to the circle.
    \end{theorem}

    \begin{formula}[Tangent--Secant Theorem]
        If $PT$ is tangent to a circle at $T$ and a secant through $P$ meets
        the circle at $A$ and $B$, then
        \[
            PT^2=PA\cdot PB.
        \]
    \end{formula}

    \begin{theorem}[Radical Axis Theorem]
        For two nonconcentric circles, the set of points having equal powers
        with respect to the two circles is a line. This line is called the
        radical axis of the two circles.
    \end{theorem}

    \begin{remark}
        Radical axes are powerful because they turn circle configurations into
        collinearity statements.
    \end{remark}

    \begin{theorem}[Ptolemy's Theorem]
        A quadrilateral $ABCD$ is cyclic if and only if
        \[
            AC\cdot BD
            =
            AB\cdot CD+BC\cdot DA
        \]
        for the appropriate cyclic order of the vertices.
    \end{theorem}


    \begin{theorem}[Ptolemy's Inequality]
        For any four points $A,B,C,D$ in the Euclidean plane,
        \[
            AC\cdot BD
            \leq
            AB\cdot CD+BC\cdot DA.
        \]
        Equality holds in the cyclic convex case and in the corresponding
        degenerate collinear case.
    \end{theorem}

    \begin{formula}[Brahmagupta's Formula]
        If a cyclic quadrilateral has side lengths $a,b,c,d$ and semiperimeter
        \[
            s=\frac{a+b+c+d}{2},
        \]
        then its area $K$ is
        \[
            K
            =
            \sqrt{(s-a)(s-b)(s-c)(s-d)}.
        \]
        Heron's formula is the limiting triangular case.
    \end{formula}

    \begin{theorem}[Varignon's Theorem]
        The midpoints of the four sides of any quadrilateral form a
        parallelogram. Its area is half the area of the original quadrilateral.
    \end{theorem}

    \begin{theorem}[Pitot's Theorem]
        If a convex quadrilateral with consecutive side lengths $a,b,c,d$ has an
        incircle tangent to all four sides, then
        \[
            a+c=b+d.
        \]
    \end{theorem}

    \begin{theorem}[Euler's Quadrilateral Theorem]
        If $M$ and $N$ are the midpoints of the diagonals of a quadrilateral with
        side lengths $a,b,c,d$ and diagonal lengths $p,q$, then
        \[
            a^2+b^2+c^2+d^2=p^2+q^2+4MN^2.
        \]
        The parallelogram law is the case $M=N$.
    \end{theorem}

    \begin{theorem}[Simson Line Theorem]
        Let $P$ be a point on the circumcircle of $\triangle ABC$. The feet of
        the perpendiculars from $P$ to the three lines $AB,BC,CA$ are collinear.
        This line is called the Simson line of $P$ with respect to
        $\triangle ABC$.
    \end{theorem}

    \begin{theorem}[Miquel's Theorem]
        In a complete quadrilateral, the circumcircles of the four triangles
        formed by choosing three of the four lines pass through a common point.
        This point is called the Miquel point.
    \end{theorem}

    \begin{theorem}[Butterfly Theorem]
        Let $M$ be the midpoint of a chord $PQ$ of a circle. Through $M$, draw
        two other chords $AB$ and $CD$, and let $AD$ and $BC$ meet $PQ$ at
        $X$ and $Y$, respectively. Then $M$ is also the midpoint of $XY$.
    \end{theorem}

    \begin{theorem}[Morley's Trisector Theorem]
        In any triangle, the adjacent trisectors of the three angles intersect
        in three points that form an equilateral triangle.
    \end{theorem}

    \begin{theorem}[Napoleon's Theorem]
        Construct equilateral triangles externally on the three sides of any
        triangle. Then the centers of these three equilateral triangles form an
        equilateral triangle.
    \end{theorem}


    \begin{theorem}[British Flag Theorem]
        For any point $P$ in the plane of a rectangle $ABCD$,
        \[
            PA^2+PC^2=PB^2+PD^2.
        \]
    \end{theorem}

    \begin{theorem}[Van Aubel's Theorem]
        Construct squares externally on the four sides of a quadrilateral. The
        segments joining the centers of opposite squares are equal in length
        and perpendicular.
    \end{theorem}

    \begin{theorem}[Viviani's Theorem]
        For any point $P$ inside an equilateral triangle of altitude $h$, the sum
        of the perpendicular distances from $P$ to the three sides is $h$.
    \end{theorem}

    \begin{theorem}[Desargues' Theorem]
        If two triangles are perspective from a point, then the three
        intersections of corresponding sides are collinear. Conversely, in the
        projective plane, collinearity of those intersections implies concurrence
        of the three lines joining corresponding vertices.
    \end{theorem}

    \begin{theorem}[Pappus' Hexagon Theorem]
        Let $A,B,C$ lie on one line and $A',B',C'$ lie on another. Then the three
        points
        \[
            AB'\cap A'B,
            \qquad
            AC'\cap A'C,
            \qquad
            BC'\cap B'C
        \]
        are collinear.
    \end{theorem}

    \begin{theorem}[Homothety Centers of Two Circles]
        For two nonconcentric circles, the two external common tangents meet at
        the external homothety center, and the two internal common tangents meet
        at the internal homothety center, when these intersections exist. Each
        homothety center lies on the line joining the two circle centers.
    \end{theorem}

    \begin{strategy}[Inversion]
        An inversion with center $O$ and radius $\rho$ sends a point
        $P\neq O$ to the point $P'$ on ray $OP$ satisfying
        \[
            OP\cdot OP'=\rho^2.
        \]
        It sends a line not through $O$ to a circle through $O$, a circle
        through $O$ to a line not through $O$, and a circle not through $O$ to
        another circle. Inversion is especially effective when many circles are
        tangent or pass through one common point.
    \end{strategy}


    \begin{formula}[Circle Inversion]
        Under inversion with center $O$ and radius $\rho$, a point $P\neq O$
        is sent to the point $P'$ on the ray $OP$ satisfying
        \[
            OP\cdot OP'=\rho^2.
        \]
        Inversion preserves angles in magnitude, sends circles through $O$ to
        lines not through $O$, and sends such lines back to circles through $O$.
    \end{formula}

    \begin{formula}[Regular Polygon Tiling Angle Condition]
        If $q$ regular $p$-gons meet edge-to-edge at each vertex of a Euclidean
        tiling, then
        \[
            q\left(1-\frac{2}{p}\right)\pi=2\pi,
        \]
        equivalently,
        \[
            (p-2)(q-2)=4.
        \]
        Thus the only tilings by one type of regular polygon are
        \[
            (p,q)=(3,6),(4,4),(6,3).
        \]
    \end{formula}

    \begin{formula}[Obtuse Triangles from Equally Spaced Points]
        Let $N$ equally spaced points lie on a circle. The number of obtuse
        triangles determined by three of the points is
        \[
            N\binom{k-1}{2}
            \quad\text{if }N=2k,
        \]
        and
        \[
            N\binom{k}{2}
            \quad\text{if }N=2k+1.
        \]
        For even $N$, triangles containing a diameter are right rather than
        obtuse.
    \end{formula}

    \begin{fact}[Alternating-Area Invariant in a Regular Even-Gon]
        Let $V_1,\dots,V_{2m}$ be the vertices of a regular $2m$-gon in cyclic
        order, and let $P$ be any point inside it. Then
        \[
            [PV_1V_2]+[PV_3V_4]+\cdots+[PV_{2m-1}V_{2m}]
        \]
        equals
        \[
            [PV_2V_3]+[PV_4V_5]+\cdots+[PV_{2m}V_1].
        \]
        The equality follows from the cancellation of alternating directed edge
        vectors.
    \end{fact}

    \begin{theorem}[Classification of Platonic Solids]
        There are exactly five Platonic solids:
        \[
            \text{tetrahedron},\quad
            \text{cube},\quad
            \text{octahedron},\quad
            \text{dodecahedron},\quad
            \text{icosahedron}.
        \]
    \end{theorem}

    \section*{Analytic Geometry}


    \begin{formula}[Coordinate Geometry Table]
        For points $A=(x_1,y_1)$ and $B=(x_2,y_2)$,

        {\small
        \renewcommand{\arraystretch}{1.18}
        \begin{longtable}{@{}L{0.27\textwidth}L{0.63\textwidth}@{}}
            \toprule
            Quantity & Formula \\
            \midrule
            \endfirsthead

            \toprule
            Quantity & Formula \\
            \midrule
            \endhead

            \bottomrule
            \endlastfoot

            Distance & $\displaystyle AB=\sqrt{(x_2-x_1)^2+(y_2-y_1)^2}$ \\
            Midpoint & $\displaystyle\left(\frac{x_1+x_2}{2},\frac{y_1+y_2}{2}\right)$ \\
            Internal division $AP:PB=m:n$ &
            $\displaystyle\left(\frac{nx_1+mx_2}{m+n},
            \frac{ny_1+my_2}{m+n}\right)$ \\
            Slope & $\displaystyle\frac{y_2-y_1}{x_2-x_1}$, when $x_1\neq x_2$ \\
            Point-slope line & $y-y_1=m(x-x_1)$ \\
            General line & $ax+by+c=0$ \\
            Distance from $(x_0,y_0)$ to $ax+by+c=0$ &
            $\displaystyle\frac{|ax_0+by_0+c|}{\sqrt{a^2+b^2}}$ \\
        \end{longtable}
        }
    \end{formula}

    \begin{formula}[Triangle Centers in Coordinates]
        For $A=(x_1,y_1)$, $B=(x_2,y_2)$, and $C=(x_3,y_3)$, the centroid is
        \[
            G
            =
            \left(
                \frac{x_1+x_2+x_3}{3},
                \frac{y_1+y_2+y_3}{3}
            \right).
        \]
        If the side lengths opposite $A,B,C$ are $a,b,c$, respectively, then the
        incenter is
        \[
            I
            =
            \frac{aA+bB+cC}{a+b+c}.
        \]
    \end{formula}

    \begin{formula}[Barycentric Coordinates]
        Let $A,B,C$ be noncollinear. Every point $P$ in their affine plane has
        unique coordinates $(\alpha,\beta,\gamma)$ satisfying
        \[
            \alpha+\beta+\gamma=1,
            \qquad
            \mathbf{p}
            =
            \alpha\mathbf{a}+\beta\mathbf{b}+\gamma\mathbf{c}.
        \]
        For $P$ inside $\triangle ABC$,
        \[
            \alpha=\frac{[PBC]}{[ABC]},
            \qquad
            \beta=\frac{[PCA]}{[ABC]},
            \qquad
            \gamma=\frac{[PAB]}{[ABC]}.
        \]
    \end{formula}

    \begin{formula}[Shoelace Formula; Gauss's Area Formula]
        For a polygon with vertices
        \[
            (x_1,y_1),(x_2,y_2),\dots,(x_n,y_n)
        \]
        in cyclic order, its area is
        \[
            \frac12
            \left|
            \sum_{i=1}^{n}
            (x_i y_{i+1}-x_{i+1}y_i)
            \right|,
        \]
        where
        \[
            (x_{n+1},y_{n+1})=(x_1,y_1).
        \]
    \end{formula}

    \begin{remark}
        The Shoelace Formula is one of the cleanest examples of analytic
        geometry: a geometric area becomes a determinant-like coordinate
        calculation.
    \end{remark}

    \begin{formula}[Oriented Area and Collinearity Test]
        For points $A=(x_1,y_1)$, $B=(x_2,y_2)$, and $C=(x_3,y_3)$, the
        oriented double area is
        \[
            2[ABC]_{\mathrm{or}}
            =
            \det
            \begin{pmatrix}
                x_1 & y_1 & 1 \\
                x_2 & y_2 & 1 \\
                x_3 & y_3 & 1
            \end{pmatrix}.
        \]
        Hence $A,B,C$ are collinear if and only if this determinant is $0$.
    \end{formula}

    \begin{formula}[Plane Rotation]
        Counterclockwise rotation through angle $\theta$ about the origin sends
        \[
            \begin{pmatrix}x\\y\end{pmatrix}
            \longmapsto
            \begin{pmatrix}
                \cos\theta & -\sin\theta \\
                \sin\theta & \cos\theta
            \end{pmatrix}
            \begin{pmatrix}x\\y\end{pmatrix}.
        \]
        In complex notation, the same transformation is simply
        \[
            z\longmapsto e^{i\theta}z.
        \]
    \end{formula}

    \begin{formula}[Reflection across a Line through the Origin]
        If a line through the origin makes angle $\theta$ with the positive
        $x$-axis, reflection across that line is
        \[
            z\longmapsto e^{2i\theta}\overline{z}
        \]
        in complex notation.
    \end{formula}

    \begin{strategy}[Affine Normalization]
        Affine transformations preserve collinearity, parallelism, ratios along
        one line, concurrence, and ratios of areas. Therefore, when a problem
        depends only on these properties, one may send a triangle to convenient
        coordinates such as
        \[
            A=(0,0),
            \qquad
            B=(1,0),
            \qquad
            C=(0,1).
        \]
        Angles, lengths, and circles are not generally preserved.
    \end{strategy}

    \begin{theorem}[Cavalieri's Principle]
        If two plane regions lie between the same pair of parallel lines and
        every line parallel to them cuts the regions in segments whose lengths
        have a constant ratio $c$, then their areas have ratio $c$. The
        three-dimensional analogue replaces segment lengths by cross-sectional
        areas and concludes the same ratio for volumes.
    \end{theorem}

    \begin{strategy}[Unfolding Method for Reflection Problems]
        A broken path reflecting from a line or face becomes a straight segment
        after reflecting the region across each boundary in succession. Thus a
        shortest reflected path is found by measuring a straight-line distance in
        the unfolded picture and then folding the path back.
    \end{strategy}

    \begin{formula}[Lattice Segment and Cell Counts]
        A segment from $(0,0)$ to $(m,n)$, where $m,n\in\N$, passes through
        \[
            m+n-\gcd(m,n)
        \]
        unit squares. The number of lattice points on the segment, including its
        endpoints, is
        \[
            \gcd(m,n)+1.
        \]
        In three dimensions, a segment from $(0,0,0)$ to $(a,b,c)$ passes through
        \[
            a+b+c
            -\gcd(a,b)-\gcd(a,c)-\gcd(b,c)
            +\gcd(a,b,c)
        \]
        unit cubes.
    \end{formula}

    \begin{formula}[Vector Test for Parallel Segments]
        Directed segments $AB$ and $CD$ are parallel if and only if
        \[
            \mathbf{b}-\mathbf{a}
            =
            \lambda(\mathbf{d}-\mathbf{c})
        \]
        for some scalar $\lambda$. Equality with $\lambda=1$ identifies
        equal directed segments and is often the fastest way to detect a hidden
        parallelogram.
    \end{formula}

    \begin{formula}[Vector Projection]
        If $\mathbf{u}\neq\mathbf{0}$, then the projection of $\mathbf{v}$ onto
        the line spanned by $\mathbf{u}$ is
        \[
            \proj_{\mathbf{u}}\mathbf{v}
            =
            \frac{\mathbf{v}\cdot\mathbf{u}}{\mathbf{u}\cdot\mathbf{u}}
            \mathbf{u}.
        \]
        The scalar component of $\mathbf{v}$ in the direction of $\mathbf{u}$ is
        \[
            \frac{\mathbf{v}\cdot\mathbf{u}}{\|\mathbf{u}\|}.
        \]
    \end{formula}

    \begin{formula}[Vector Equations of Lines and Planes]
        A line through a point $\mathbf{p}$ with direction vector
        $\mathbf{v}\neq\mathbf{0}$ is
        \[
            \mathbf{x}=\mathbf{p}+t\mathbf{v}
            \qquad (t\in\R).
        \]
        A plane through a point $\mathbf{p}$ with normal vector
        $\mathbf{n}\neq\mathbf{0}$ is
        \[
            (\mathbf{x}-\mathbf{p})\cdot\mathbf{n}=0.
        \]
    \end{formula}

    \begin{formula}[Distance from a Point to a Line]
        In $\R^3$, the distance from a point $\mathbf{p}_0$ to the line
        \[
            \mathbf{x}=\mathbf{p}+t\mathbf{v}
        \]
        is
        \[
            \frac{\|(\mathbf{p}_0-\mathbf{p})\times\mathbf{v}\|}{\|\mathbf{v}\|}.
        \]
        In $\R^2$, the distance from $(x_0,y_0)$ to the line $ax+by+c=0$ is
        \[
            \frac{|ax_0+by_0+c|}{\sqrt{a^2+b^2}}.
        \]
    \end{formula}

    \begin{formula}[Distance from a Point to a Plane]
        The distance from a point $\mathbf{p}_0=(x_0,y_0,z_0)$ to the plane
        \[
            ax+by+cz+d=0
        \]
        is
        \[
            \frac{|ax_0+by_0+cz_0+d|}{\sqrt{a^2+b^2+c^2}}.
        \]
    \end{formula}

    \begin{formula}[Sphere Equation]
        A sphere with center $\mathbf{c}=(a,b,c)$ and radius $r$ has equation
        \[
            (x-a)^2+(y-b)^2+(z-c)^2=r^2.
        \]
        Equivalently,
        \[
            \|\mathbf{x}-\mathbf{c}\|=r.
        \]
    \end{formula}

    \begin{strategy}[Line--Circle Intersection by a Discriminant]
        Substitute a parametric line into a circle equation. The resulting
        quadratic equation has two, one, or no real solutions according as its
        discriminant is positive, zero, or negative. The zero-discriminant case
        is the tangency condition.
    \end{strategy}

    \begin{formula}[Circle Equation]
        The circle with center $(a,b)$ and radius $r$ has equation
        \[
            (x-a)^2+(y-b)^2=r^2.
        \]
        The general equation
        \[
            x^2+y^2+Dx+Ey+F=0
        \]
        can be converted to center-radius form by completing the square.
    \end{formula}

    \begin{theorem}[Three Perpendiculars Theorem]
        Let a line $\ell$ lie in a plane $\Pi$, and let $P$ be a point outside
        $\Pi$. If $H$ is the foot of the perpendicular from $P$ to $\Pi$, and if
        the line $m$ in $\Pi$ through $H$ is perpendicular to $\ell$, then the
        line through $P$ and the intersection point of $m$ with $\ell$ is also
        perpendicular to $\ell$.
    \end{theorem}

    \begin{remark}
        In coordinate geometry, the three perpendiculars theorem is often
        replaced by dot products and projections. Still, it is a standard
        spatial-geometry fact from the classical school curriculum.
    \end{remark}

    \begin{definition}[Conic Section]
        A conic section is a curve obtained by intersecting a plane with a
        double cone. The nondegenerate conics are:
        \[
            \text{ellipse},\qquad
            \text{parabola},\qquad
            \text{hyperbola}.
        \]
        Equivalently, a conic can be described using a focus, a directrix, and
        an eccentricity $e$.
    \end{definition}

    \begin{formula}[Focus--Directrix Definition of Conics]
        A conic is the set of points $P$ satisfying
        \[
            \frac{\dist(P,\text{focus})}
                 {\dist(P,\text{directrix})}
            =
            e.
        \]
        It is an ellipse if
        \[
            0<e<1,
        \]
        a parabola if
        \[
            e=1,
        \]
        and a hyperbola if
        \[
            e>1.
        \]
    \end{formula}

    \begin{formula}[Standard Forms of Conics]
        The standard form of an ellipse is
        \[
            \frac{x^2}{a^2}+\frac{y^2}{b^2}=1.
        \]
        The standard form of a hyperbola is
        \[
            \frac{x^2}{a^2}-\frac{y^2}{b^2}=1.
        \]
        The standard form of a parabola opening to the right is
        \[
            y^2=4px.
        \]
    \end{formula}

    \begin{formula}[Eccentricity of an Ellipse and a Hyperbola]
        For an ellipse
        \[
            \frac{x^2}{a^2}+\frac{y^2}{b^2}=1
            \qquad (a\geq b>0),
        \]
        the eccentricity is
        \[
            e=\frac{c}{a},
            \qquad
            c^2=a^2-b^2.
        \]
        For a hyperbola
        \[
            \frac{x^2}{a^2}-\frac{y^2}{b^2}=1,
        \]
        the eccentricity is
        \[
            e=\frac{c}{a},
            \qquad
            c^2=a^2+b^2.
        \]
    \end{formula}

    \begin{formula}[Tangent Lines to Standard Conics]
        At a point $(x_1,y_1)$ on the indicated conic, the tangent line is
        \[
            xx_1+yy_1=r^2
            \qquad
            \text{for }x^2+y^2=r^2,
        \]
        \[
            yy_1=2a(x+x_1)
            \qquad
            \text{for }y^2=4ax,
        \]
        \[
            \frac{xx_1}{a^2}+\frac{yy_1}{b^2}=1
            \qquad
            \text{for }\frac{x^2}{a^2}+\frac{y^2}{b^2}=1,
        \]
        and
        \[
            \frac{xx_1}{a^2}-\frac{yy_1}{b^2}=1
            \qquad
            \text{for }\frac{x^2}{a^2}-\frac{y^2}{b^2}=1.
        \]
    \end{formula}

    \begin{theorem}[Reflection Property of Conics]
        An ellipse reflects a ray from one focus to the other focus. A parabola
        reflects rays parallel to its axis through its focus. A hyperbola has
        the analogous reflection property involving its two foci and branches.
    \end{theorem}

    \begin{remark}
        The reflection property gives a geometric reason why conics appear in
        optics, satellite dishes, whispering galleries, and orbital mechanics.
    \end{remark}

    \begin{theorem}[Quadratic Classification of Conics]
        Consider a real quadratic equation
        \[
            Ax^2+Bxy+Cy^2+Dx+Ey+F=0.
        \]
        The discriminant
        \[
            B^2-4AC
        \]
        determines the type of the nondegenerate conic:
        \[
            B^2-4AC<0
            \quad\Rightarrow\quad
            \text{ellipse-type},
        \]
        \[
            B^2-4AC=0
            \quad\Rightarrow\quad
            \text{parabola-type},
        \]
        and
        \[
            B^2-4AC>0
            \quad\Rightarrow\quad
            \text{hyperbola-type}.
        \]
    \end{theorem}

    \begin{remark}
        This is the analytic-geometry counterpart of classifying conic sections
        by their shapes.
    \end{remark}

    \begin{formula}[Polar Equation of a Conic]
        With a focus at the pole, a conic has polar equation
        \[
            r
            =
            \frac{\ell}{1+e\cos\theta}
        \]
        after choosing a suitable axis. Here $e$ is the eccentricity and
        $\ell$ is the semi-latus rectum.
    \end{formula}

    \begin{remark}
        The same equation describes ellipses, parabolas, and hyperbolas
        depending on whether $e<1$, $e=1$, or $e>1$.
    \end{remark}

    \begin{definition}[Apollonius Circle]
        Given two fixed points $A$ and $B$ and a positive constant
        $k\neq 1$, the locus of points $P$ satisfying
        \[
            \frac{PA}{PB}=k
        \]
        is a circle. This circle is called an Apollonius circle.
    \end{definition}

    \begin{remark}
        Apollonius circles turn a ratio of distances into a concrete circle,
        which is a common hidden structure in geometry problems.
    \end{remark}

    \begin{theorem}[Descartes' Circle Theorem; Soddy Circle Theorem]
        Suppose four circles are mutually tangent, and let their curvatures be
        \[
            k_1,k_2,k_3,k_4,
        \]
        where curvature means the reciprocal of the radius, with sign convention
        for an enclosing circle. Then
        \[
            (k_1+k_2+k_3+k_4)^2
            =
            2(k_1^2+k_2^2+k_3^2+k_4^2).
        \]
    \end{theorem}

    \begin{remark}
        Descartes' Circle Theorem is a striking formula: tangency of four
        circles becomes one quadratic equation in their curvatures.
    \end{remark}

    \begin{formula}[Boundary Lattice Points on a Segment]
        If a lattice segment has displacement $(\Delta x,\Delta y)$, then the
        number of lattice points on it, including both endpoints, is
        \[
            \gcd(|\Delta x|,|\Delta y|)+1.
        \]
        Summing the gcd contributions over polygon edges gives the number $B$ of
        boundary lattice points without double-counting vertices.
    \end{formula}

    \begin{theorem}[Pick's Theorem]
        Let $P$ be a simple polygon whose vertices lie on lattice points in the
        plane. If $I$ is the number of lattice points strictly inside $P$ and
        $B$ is the number of lattice points on the boundary of $P$, then
        \[
            \area(P)
            =
            I+\frac{B}{2}-1.
        \]
    \end{theorem}

    \begin{remark}
        Pick's Theorem is a beautiful bridge between geometry and arithmetic:
        area is determined by counting lattice points.
    \end{remark}

    \chapter{Number Theory}

    \relatedcourses{(MATH304) Introduction to Number Theory}

    \relatedbooks{Introduction to Number Theory
        (William W. Adams, Larry Joel Goldstein)}

    \begin{theorem}[Fundamental Theorem of Arithmetic]
        Every integer $n>1$ can be written as a product of prime numbers.
        Moreover, this factorization is unique up to the order of the factors.
        Equivalently,
        \[
            n=p_1^{\alpha_1}p_2^{\alpha_2}\cdots p_k^{\alpha_k},
        \]
        where $p_1,\dots,p_k$ are distinct primes and
        $\alpha_1,\dots,\alpha_k$ are positive integers, and this expression is
        unique up to reordering.
    \end{theorem}

    \begin{remark}
        This theorem is often called unique prime factorization. It is the
        reason prime numbers are treated as the basic building blocks of
        integers.
    \end{remark}

    \begin{theorem}[Euclid's Theorem; Infinitude of Primes]
        There are infinitely many prime numbers.
    \end{theorem}

    \begin{idea}
        If $p_1,\dots,p_n$ are finitely many primes, consider
        \[
            p_1p_2\cdots p_n+1.
        \]
        This number is not divisible by any prime on the list.
    \end{idea}

    \begin{theorem}[B\'ezout's Identity; B\'ezout's Lemma]
        Let $a,b\in\Z$, not both zero. Then there exist integers $x,y\in\Z$
        such that
        \[
            ax+by=\gcd(a,b).
        \]
        In particular, if $\gcd(a,b)=1$, then there exist integers $x,y\in\Z$
        such that
        \[
            ax+by=1.
        \]
    \end{theorem}

    \begin{remark}
        B\'ezout's identity is the algebraic heart of the Euclidean algorithm and
        modular inverses.
    \end{remark}

    \begin{algorithm}[Euclidean Algorithm]
        To compute $\gcd(a,b)$ for integers $a,b\in\Z$, repeatedly divide the
        larger number by the smaller number and replace the pair by
        \[
            (b,r),
        \]
        where $r$ is the remainder. Continue until the remainder becomes $0$.
        The last nonzero remainder is $\gcd(a,b)$.
    \end{algorithm}

    \begin{proposition}[Basic Rules of Congruences]
        If
        \[
            a\equiv b\pmod n,
            \qquad
            c\equiv d\pmod n,
        \]
        then
        \[
            a+c\equiv b+d\pmod n,
            \qquad
            ac\equiv bd\pmod n.
        \]
        If $\gcd(c,n)=1$, then
        \[
            ac\equiv bc\pmod n
            \quad\Longrightarrow\quad
            a\equiv b\pmod n.
        \]
    \end{proposition}

    \begin{theorem}[Linear Congruence Criterion]
        The congruence
        \[
            ax\equiv b\pmod n
        \]
        has a solution if and only if
        \[
            \gcd(a,n)\mid b.
        \]
        If $d=\gcd(a,n)$ and $d\mid b$, then there are exactly $d$ solutions
        modulo $n$.
    \end{theorem}

    \begin{formula}[Least Common Multiple from Prime Factorization]
        If
        \[
            a=\prod_p p^{\alpha_p},
            \qquad
            b=\prod_p p^{\beta_p},
        \]
        then
        \[
            \gcd(a,b)=\prod_p p^{\min(\alpha_p,\beta_p)},
            \qquad
            \lcm(a,b)=\prod_p p^{\max(\alpha_p,\beta_p)}.
        \]
        More generally, the least common multiple of several integers is found
        by taking the largest exponent of each prime appearing in the list.
    \end{formula}

    \begin{formula}[Basic Divisibility Tests]
        For a decimal integer $N$, the following tests are useful.
        \begin{itemize}
            \item Divisibility by $2$, $4$, and $8$ depends on the last $1$, $2$,
                  and $3$ digits, respectively.
            \item Divisibility by $3$ and $9$ depends on the digit sum.
            \item Divisibility by $5$ depends on the last digit.
            \item Divisibility by $11$ depends on the alternating sum of the
                  digits.
        \end{itemize}
    \end{formula}

    \begin{strategy}
        Digit problems often combine an lcm condition, a digit-sum condition,
        and a last-digits condition. First minimize the number of digits, then
        arrange the final few digits to satisfy the strongest divisibility
        conditions.
    \end{strategy}

    \begin{definition}[Complete Residue System]
        A set of $n$ integers is a complete residue system modulo $n$ if its
        elements represent every residue class modulo $n$ exactly once. To prove
        that $a_0,\dots,a_{n-1}$ form such a system, it is enough to prove
        \[
            a_i\equiv a_j\pmod n
            \quad\Longrightarrow\quad
            i=j.
        \]
    \end{definition}

    \begin{strategy}[Modular Casework and Base Representation]
        When a problem repeats division by $N$, rewrite the integer in base $N$.
        The quotient deletes the final digit and the remainder reveals it. For
        divisibility or existence problems, first split integers into residue
        classes and look for a complete residue system rather than checking
        values individually.
    \end{strategy}


    \begin{strategy}[Infinite Descent]
        Assume a positive-integer solution exists and choose one minimizing a
        natural positive quantity. Construct a strictly smaller solution of the
        same type, contradicting well-ordering.
    \end{strategy}

    \begin{strategy}[Vieta Jumping]
        If an integer equation is quadratic in one variable, regard that
        variable as one root of a monic quadratic. Vi\`ete's formulas give the
        other integral root. Prove that it yields a smaller positive solution
        and apply infinite descent.
    \end{strategy}

    \begin{theorem}[Frobenius Coin Theorem for Two Denominations]
        If $a$ and $b$ are coprime positive integers, the largest integer not
        representable as
        \[
            ax+by,
            \qquad
            x,y\in\Nzero,
        \]
        is $ab-a-b$. Exactly
        \[
            \frac{(a-1)(b-1)}{2}
        \]
        positive integers are not representable.
    \end{theorem}

    \begin{theorem}[Chinese Remainder Theorem; CRT]
        Let $m_1,\dots,m_k$ be pairwise relatively prime positive integers.
        Then the system
        \[
            x\equiv a_1\pmod {m_1},
            \qquad
            x\equiv a_2\pmod {m_2},
            \qquad
            \dots,
            \qquad
            x\equiv a_k\pmod {m_k}
        \]
        has a solution. Moreover, the solution is unique modulo
        \[
            M=m_1m_2\cdots m_k.
        \]
    \end{theorem}

    \begin{strategy}
        CRT is the standard theorem for combining several congruence conditions
        into one congruence.
    \end{strategy}

    \begin{formula}[Euler's Phi Function Formula]
        If
        \[
            n=p_1^{\alpha_1}p_2^{\alpha_2}\cdots p_k^{\alpha_k}
        \]
        is the prime factorization of $n$, then
        \[
            \varphi(n)
            =
            n\prod_{i=1}^{k}\left(1-\frac{1}{p_i}\right).
        \]
        In particular, for a prime power,
        \[
            \varphi(p^\alpha)=p^\alpha-p^{\alpha-1}
            =
            p^{\alpha-1}(p-1).
        \]
    \end{formula}

    \begin{formula}[Sum-of-Divisors Function]
        If
        \[
            n=p_1^{\alpha_1}\cdots p_k^{\alpha_k},
        \]
        then the number of positive divisors of $n$ is
        \[
            \tau(n)=(\alpha_1+1)\cdots(\alpha_k+1),
        \]
        and the sum of positive divisors of $n$ is
        \[
            \sigma(n)
            =
            \prod_{i=1}^{k}\frac{p_i^{\alpha_i+1}-1}{p_i-1}.
        \]
    \end{formula}

    \begin{formula}[$p$-Adic Valuation of a Factorial; Legendre's Formula]
        For a prime $p$ and a positive integer $n$,
        \[
            \nu_p(n!)
            =
            \sum_{k=1}^{\infty}
            \left\lfloor\frac{n}{p^k}\right\rfloor.
        \]
        Only finitely many terms are nonzero. This formula gives the exponent
        of $p$ in $n!$ and is especially useful for trailing-zero and
        divisibility questions.
    \end{formula}

    \begin{theorem}[M\"obius Inversion Formula]
        If arithmetic functions $f$ and $g$ satisfy
        \[
            g(n)=\sum_{d\mid n}f(d),
        \]
        then
        \[
            f(n)=\sum_{d\mid n}\mu(d)g\!\left(\frac{n}{d}\right).
        \]
    \end{theorem}

    \begin{theorem}[Fermat's Little Theorem]
        Let $p$ be a prime number. If $p\nmid a$, then
        \[
            a^{p-1}\equiv 1\pmod p.
        \]
        Equivalently, for every integer $a$,
        \[
            a^p\equiv a\pmod p.
        \]
    \end{theorem}

    \begin{theorem}[Euler's Theorem; Euler--Fermat Theorem]
        If $\gcd(a,n)=1$, then
        \[
            a^{\varphi(n)}\equiv 1\pmod n.
        \]
    \end{theorem}

    \begin{remark}
        Fermat's Little Theorem is the special case of Euler's Theorem when
        $n=p$ is prime, since $\varphi(p)=p-1$.
    \end{remark}

    \begin{proposition}[Multiplicative Order]
        If $\gcd(a,n)=1$, then $[a]_n$ lies in the finite group
        $(\Z/n\Z)^\times$. Its multiplicative order $\ord_n(a)$ divides
        \[
            |(\Z/n\Z)^\times|=\varphi(n).
        \]
        Consequently, for every positive integer $k$,
        \[
            a^k\equiv 1\pmod n
            \quad\Longleftrightarrow\quad
            \ord_n(a)\mid k.
        \]
    \end{proposition}

    \begin{theorem}[Primitive Root Theorem; Gauss's Theorem on Primitive Roots]
        The multiplicative group $(\Z/n\Z)^\times$ is cyclic if and only if
        \[
            n\in\{1,2,4,p^k,2p^k\},
        \]
        where $p$ is an odd prime and $k\geq 1$.
    \end{theorem}

    \begin{theorem}[Existence and Uniqueness of Finite Fields]
        Every finite field has $q=p^n$ elements for some prime $p$ and positive
        integer $n$. Conversely, for each prime power $q=p^n$, there exists a
        field $\mathbb{F}_q$ with $q$ elements, unique up to isomorphism. Its
        multiplicative group $\mathbb{F}_q^\times$ is cyclic of order $q-1$.
    \end{theorem}

    \begin{algorithm}[Modular Exponentiation by Period Reduction]
        To find the last digits of $a^n$, compute $a^n$ modulo a suitable power
        of $10$. Look for a short period in the powers of $a$ modulo that
        modulus. For example,
        \[
            7^4=2401\equiv 1\pmod {100},
        \]
        so exponents of $7$ may be reduced modulo $4$ when finding the last two
        digits.
    \end{algorithm}

    \begin{strategy}
        Last-digit and last-two-digit questions are modular arithmetic problems.
        Work modulo $10$ or $100$, not with the full number.
    \end{strategy}

    \begin{theorem}[Wilson's Theorem]
        A positive integer $p>1$ is prime if and only if
        \[
            (p-1)!\equiv -1\pmod p.
        \]
    \end{theorem}

    \begin{remark}
        Wilson's theorem is famous because it gives a clean factorial
        characterization of primes, even though it is not efficient as a
        primality test.
    \end{remark}


    \begin{theorem}[Lucas's Theorem]
        Let $p$ be prime and write
        \[
            n=\sum_{i=0}^{r}n_ip^i,
            \qquad
            k=\sum_{i=0}^{r}k_ip^i
        \]
        in base $p$. Then
        \[
            \binom{n}{k}
            \equiv
            \prod_{i=0}^{r}\binom{n_i}{k_i}
            \pmod p.
        \]
    \end{theorem}

    \begin{theorem}[Kummer's Theorem]
        The valuation
        \[
            \nu_p\!\left(\binom{n}{k}\right)
        \]
        equals the number of carries when $k$ and $n-k$ are added in base $p$.
    \end{theorem}

    \begin{theorem}[Lifting the Exponent Lemma; LTE]
        Let $p$ be an odd prime, and suppose that $p\mid x-y$ while
        $p\nmid xy$. Then, for every positive integer $n$,
        \[
            \nu_p(x^n-y^n)
            =
            \nu_p(x-y)+\nu_p(n).
        \]
        A standard companion form is: if $n$ is odd, $p\mid x+y$, and
        $p\nmid xy$, then
        \[
            \nu_p(x^n+y^n)
            =
            \nu_p(x+y)+\nu_p(n).
        \]
    \end{theorem}

    \begin{fact}[Quadratic Residues Modulo Small Moduli]
        Squares have very restricted residues modulo small integers:
        \[
            x^2\equiv 0,1\pmod 4,
            \qquad
            x^2\equiv 0,1,4\pmod 8.
        \]
        In particular, an integer congruent to $2$ or $3$ modulo $4$ cannot be a
        square, and an integer congruent to $2,3,5,6,7$ modulo $8$ cannot be a
        square.
    \end{fact}

    \begin{tactic}
        For Diophantine equations, first reduce modulo $4$, $8$, $3$, $5$, or $9$. Many impossible equations are eliminated by
        the small list of possible square residues.
    \end{tactic}

    \begin{definition}[Legendre Symbol]
        Let $p$ be an odd prime. The Legendre symbol is defined by
        \[
            \left(\frac{a}{p}\right)
            =
            \begin{cases}
                1, & \text{if }a\text{ is a quadratic residue modulo }p, \\
                -1, & \text{if }a\text{ is a quadratic nonresidue modulo }p, \\
                0, & \text{if }p\mid a.
            \end{cases}
        \]
    \end{definition}

    \begin{theorem}[Euler's Criterion]
        Let $p$ be an odd prime. For any integer $a$,
        \[
            \left(\frac{a}{p}\right)
            \equiv
            a^{(p-1)/2}
            \pmod p.
        \]
        Equivalently, if $p\nmid a$, then
        \[
            a^{(p-1)/2}
            \equiv
            \begin{cases}
                1\pmod p, & \text{if }a\text{ is a quadratic residue modulo }p, \\
                -1\pmod p, & \text{if }a\text{ is a quadratic nonresidue modulo }p.
            \end{cases}
        \]
    \end{theorem}

    \begin{theorem}[Law of Quadratic Reciprocity; Gauss's Theorem of Quadratic Reciprocity]
        Let $p$ and $q$ be distinct odd primes. Then
        \[
            \left(\frac{p}{q}\right)
            \left(\frac{q}{p}\right)
            =
            (-1)^{\frac{p-1}{2}\frac{q-1}{2}}.
        \]
        Equivalently,
        \[
            \left(\frac{p}{q}\right)
            =
            \left(\frac{q}{p}\right)
        \]
        unless both $p$ and $q$ are congruent to $3$ modulo $4$.
    \end{theorem}

    \begin{remark}
        Quadratic reciprocity is one of the crown jewels of elementary number
        theory. Gauss called it the golden theorem.
    \end{remark}

    \begin{theorem}[Fermat's Two-Square Theorem; Fermat's Theorem on Sums of Two Squares]
        An odd prime $p$ can be written as a sum of two squares,
        \[
            p=a^2+b^2
        \]
        for some integers $a,b\in\Z$, if and only if
        \[
            p\equiv 1\pmod 4.
        \]
        Also,
        \[
            2=1^2+1^2.
        \]
    \end{theorem}


    \begin{formula}[Brahmagupta--Fibonacci Identity]
        \[
            (a^2+b^2)(c^2+d^2)
            =
            (ac-bd)^2+(ad+bc)^2
            =
            (ac+bd)^2+(ad-bc)^2.
        \]
    \end{formula}

    \begin{formula}[Sophie Germain's Identity]
        \[
            a^4+4b^4
            =
            (a^2-2ab+2b^2)(a^2+2ab+2b^2).
        \]
    \end{formula}

    \begin{theorem}[Lagrange's Four-Square Theorem]
        Every nonnegative integer can be written as a sum of four integer
        squares. That is, for every $n\in\Nzero$, there exist integers
        $a,b,c,d\in\Z$ such that
        \[
            n=a^2+b^2+c^2+d^2.
        \]
    \end{theorem}

    \begin{remark}
        The two-square theorem has a congruence condition. The four-square
        theorem has no exception at all.
    \end{remark}

    \begin{theorem}[Continued-Fraction Best Approximation Theorem]
        Let $p_k/q_k$ be a convergent of the simple continued fraction of an
        irrational number $\alpha$. Then
        \[
            \left|\alpha-\frac{p_k}{q_k}\right|
            <
            \frac{1}{q_kq_{k+1}}
            <
            \frac{1}{q_k^2}.
        \]
        Moreover, convergents provide best rational approximations in the sense
        that a fraction with a smaller denominator cannot improve the error in
        the standard best-approximation comparison.
    \end{theorem}

    \begin{strategy}
        Problems asking for an unusually accurate fraction with the smallest
        possible denominator often signal a continued-fraction expansion.
    \end{strategy}


    \begin{definition}[Pell Equation]
        A Pell equation has the form
        \[
            x^2-Dy^2=1,
        \]
        where $D$ is a positive nonsquare integer. If $(x_1,y_1)$ is the
        fundamental positive solution, then all positive solutions satisfy
        \[
            x_n+y_n\sqrt{D}
            =
            (x_1+y_1\sqrt{D})^n.
        \]
        Continued fractions of $\sqrt{D}$ produce the fundamental solution.
    \end{definition}

    \begin{theorem}[Dirichlet's Theorem on Arithmetic Progressions]
        If $a$ and $d$ are relatively prime positive integers, then the
        arithmetic progression
        \[
            a,\ a+d,\ a+2d,\ a+3d,\dots
        \]
        contains infinitely many primes.
    \end{theorem}

    \begin{remark}
        Dirichlet's theorem is a major result connecting primes and arithmetic
        progressions. The condition $\gcd(a,d)=1$ is necessary.
    \end{remark}

    \begin{theorem}[Bertrand's Postulate; Chebyshev's Theorem]
        For every integer $n>1$, there exists a prime $p$ such that
        \[
            n<p<2n.
        \]
    \end{theorem}

    \begin{theorem}[Prime Number Theorem]
        Let $\pi(x)$ be the number of primes less than or equal to $x$. Then
        \[
            \pi(x)\sim \frac{x}{\ln x}
            \qquad
            \text{as }x\to\infty.
        \]
    \end{theorem}

    \begin{idea}
        The Prime Number Theorem says that the ``probability'' that a large
        integer near $x$ is prime is roughly $1/\ln x$. This is a mental model,
        not a literal probability statement.
    \end{idea}

    \begin{definition}[Riemann Zeta Function]
        For $\Real(s)>1$, the Riemann zeta function is defined by
        \[
            \zeta(s)
            =
            \sum_{n=1}^{\infty}\frac{1}{n^s}.
        \]
        It also has the Euler product
        \[
            \zeta(s)
            =
            \prod_{p\text{ prime}}
            \frac{1}{1-p^{-s}}.
        \]
    \end{definition}

    \begin{remark}
        The Euler product is the bridge between the zeta function and prime
        numbers. The Riemann Hypothesis asserts that the nontrivial zeros of
        $\zeta(s)$ all have real part $1/2$.
    \end{remark}

    \begin{theorem}[Euclid--Euler Theorem on Perfect Numbers]
        An even positive integer $N$ is perfect if and only if
        \[
            N=2^{p-1}(2^p-1),
        \]
        where $2^p-1$ is a Mersenne prime.
    \end{theorem}

    \begin{remark}
        A perfect number is a positive integer equal to the sum of its positive
        proper divisors. For example,
        \[
            6=1+2+3.
        \]
    \end{remark}

    \begin{theorem}[Fermat's Last Theorem]
        For every integer $n\geq 3$, there do not exist positive integers
        $a,b,c$ such that
        \[
            a^n+b^n=c^n.
        \]
    \end{theorem}

    \begin{remark}
        Fermat's Last Theorem was proved by Andrew Wiles. It is one of the most
        famous theorems in the history of mathematics.
    \end{remark}

    \chapter{Algebra}

    \relatedcourses{(MATH301) Modern Algebra I, (MATH302) Modern Algebra II}

    \relatedbooks{Abstract Algebra (David S. Dummit, Richard M. Foote)}

    \begin{definition}[Group]
        A group is a set $G$ with a binary operation $\ast:G\times G\to G$ s.t.
        \[
            \forall a,b,c\in G,
            \qquad
            (a\ast b)\ast c=a\ast(b\ast c),
        \]
        \[
            \exists e\in G\ \text{s.t.}\quad
            \forall a\in G,
            \qquad
            e\ast a=a\ast e=a,
        \]
        and
        \[
            \forall a\in G,
            \qquad
            \exists a^{-1}\in G\ \text{s.t.}\quad
            a\ast a^{-1}=a^{-1}\ast a=e.
        \]
        The group is abelian if
        \[
            \forall a,b\in G,
            \qquad
            a\ast b=b\ast a.
        \]
    \end{definition}

    \begin{theorem}[Disjoint-Cycle Decomposition of a Permutation]
        Every permutation $\sigma\in S_n$ can be written as a product of
        disjoint cycles, uniquely up to the order of the cycles and cyclic
        rotation within each cycle. If $c(\sigma)$ is the number of cycles,
        including fixed points, then the minimum number of transpositions needed
        to express $\sigma$ is
        \[
            n-c(\sigma).
        \]
    \end{theorem}

    \begin{formula}[Inversions and the Sign of a Permutation]
        The inversion number is
        \[
            \inv(\sigma)
            =
            \left|
                \{(i,j)\mid i<j,\ \sigma(i)>\sigma(j)\}
            \right|.
        \]
        The sign satisfies
        \[
            \sgn(\sigma)
            =
            (-1)^{\inv(\sigma)}
            =
            (-1)^{n-c(\sigma)}.
        \]
        Thus any transposition reverses parity.
    \end{formula}

    \begin{fact}[Dihedral Group]
        The symmetries of a regular $n$-gon form the dihedral group
        \[
            D_{2n}
            =
            \langle r,s\mid r^n=e,\ s^2=e,\ srs=r^{-1}\rangle.
        \]
        Its elements are
        \[
            e,r,r^2,\dots,r^{n-1},
            \qquad
            s,rs,r^2s,\dots,r^{n-1}s.
        \]
        This presentation is the standard algebraic model for rotation and
        reflection counting with Burnside's Lemma.
    \end{fact}

    \begin{definition}[Cyclic Group]
        A group $G$ is cyclic if there exists an element $g\in G$ such that every
        element of $G$ is a power of $g$. In this case, we write
        \[
            G=\langle g\rangle.
        \]
    \end{definition}

    \begin{theorem}[Lagrange's Theorem]
        If $G$ is a finite group and $H\leq G$, then
        \[
            |H|\mid |G|.
        \]
        Equivalently, the order of every subgroup divides the order of the group.
    \end{theorem}

    \begin{corollary}[Order of an Element]
        If $G$ is a finite group and $g\in G$, then the order of $g$ divides
        $|G|$. In particular,
        \[
            g^{|G|}=e.
        \]
    \end{corollary}

    \begin{theorem}[Cauchy's Theorem]
        If $G$ is a finite group and a prime $p$ divides $|G|$, then $G$ contains
        an element of order $p$.
    \end{theorem}

    \begin{theorem}[First Isomorphism Theorem for Groups]
        If $\varphi:G\to H$ is a group homomorphism, then
        \[
            G/\ker\varphi\cong \image\varphi.
        \]
    \end{theorem}

    \begin{theorem}[Cayley's Theorem]
        Every group is isomorphic to a subgroup of a symmetric group. In
        particular, every finite group of order $n$ is isomorphic to a subgroup
        of $S_n$.
    \end{theorem}

    \begin{theorem}[Orbit--Stabilizer Theorem]
        If a finite group $G$ acts on a set $X$ and $x\in X$, then
        \[
            |G\cdot x|
            =
            [G:G_x]
            =
            \frac{|G|}{|G_x|},
        \]
        where $G\cdot x$ is the orbit of $x$ and $G_x$ is its stabilizer.
    \end{theorem}

    \begin{theorem}[Class Equation]
        Let $G$ be a finite group, and choose one representative $x_i$ from
        each conjugacy class not contained in the center $Z(G)$. Then
        \[
            |G|
            =
            |Z(G)|
            +
            \sum_i [G:C_G(x_i)].
        \]
    \end{theorem}

    \begin{theorem}[Sylow Theorems]
        Let $G$ be a finite group and suppose
        \[
            |G|=p^k m,
            \qquad
            p\nmid m,
        \]
        where $p$ is prime. Then $G$ has a subgroup of order $p^k$, called a
        Sylow $p$-subgroup. Moreover, all Sylow $p$-subgroups are conjugate, and
        the number $n_p$ of Sylow $p$-subgroups satisfies
        \[
            n_p\equiv 1\pmod p,
            \qquad
            n_p\mid m.
        \]
    \end{theorem}

    \begin{strategy}
        For a finite group, subgroup orders and element orders must divide the group
        order. When a
        prime divides $|G|$, Cauchy's theorem guarantees an element of that prime
        order.
    \end{strategy}

    \begin{fact}[Quaternion Group]
        The quaternion group is
        \[
            Q_8=\{\pm 1,\pm i,\pm j,\pm k\},
        \]
        with
        \[
            i^2=j^2=k^2=ijk=-1.
        \]
        It is nonabelian, but every subgroup of $Q_8$ is normal.
    \end{fact}

    \begin{remark}
        The quaternion group is a standard small group with behavior that is
        less intuitive than cyclic or dihedral groups, so it is a natural
        source of quick algebra questions.
    \end{remark}

    \begin{theorem}[Fundamental Theorem of Finite Abelian Groups; Classification Theorem for Finite Abelian Groups]
        Every finite abelian group is isomorphic to a direct product of cyclic
        groups of prime-power order. In particular, finite abelian groups are
        built from cyclic groups.
    \end{theorem}

    \begin{definition}[Ring and Field]
        A ring is a set equipped with addition and multiplication satisfying the
        usual distributive laws. A field is a commutative ring in which every
        nonzero element has a multiplicative inverse.
    \end{definition}

    \begin{definition}[Integral Domain]
        An integral domain is a commutative ring with identity in which there are
        no zero divisors. In other words, if $ab=0$, then $a=0$ or $b=0$.
    \end{definition}

    \begin{formula}[Basic Factorization Identities]
        The most frequently used polynomial factorizations include
        \[
            a^2-b^2=(a-b)(a+b),
        \]
        \[
            a^3-b^3=(a-b)(a^2+ab+b^2),
            \qquad
            a^3+b^3=(a+b)(a^2-ab+b^2).
        \]
        More generally,
        \[
            a^n-b^n
            =
            (a-b)\sum_{k=0}^{n-1}a^{n-1-k}b^k,
        \]
        and, when $n$ is odd,
        \[
            a^n+b^n
            =
            (a+b)\sum_{k=0}^{n-1}(-1)^k a^{n-1-k}b^k.
        \]
    \end{formula}

    \begin{formula}[Cubic Factorization Identity]
        For all $a,b,c$,
        \[
            a^3+b^3+c^3-3abc
            =
            (a+b+c)(a^2+b^2+c^2-ab-bc-ca),
        \]
        or equivalently,
        \[
            a^3+b^3+c^3-3abc
            =
            \frac12(a+b+c)
            \left((a-b)^2+(b-c)^2+(c-a)^2\right).
        \]
        In particular, if $a+b+c=0$, then
        \[
            a^3+b^3+c^3=3abc.
        \]
    \end{formula}

    \begin{algorithm}[Polynomial Division Algorithm]
        Let $F$ be a field. If $f(x),g(x)\in F[x]$ and $g(x)\neq 0$, then there
        exist unique polynomials $q(x),r(x)\in F[x]$ such that
        \[
            f(x)=q(x)g(x)+r(x),
            \qquad
            r(x)=0\lor \deg r<\deg g.
        \]
    \end{algorithm}

    \begin{theorem}[Factor Theorem]
        For a polynomial $f(x)\in F[x]$ and an element $a\in F$,
        \[
            f(a)=0
            \quad\Longleftrightarrow\quad
            x-a \text{ divides } f(x).
        \]
    \end{theorem}


    \begin{proposition}[Root-Counting Principle for Polynomials]
        A nonzero polynomial of degree $n$ over a field has at most $n$ roots.
        Hence two polynomials of degree at most $n$ that agree at $n+1$
        distinct points are identical.
    \end{proposition}

    \begin{formula}[Lagrange Interpolation Formula]
        For distinct $x_0,\dots,x_n$, the unique polynomial of degree at most
        $n$ satisfying $P(x_i)=y_i$ is
        \[
            P(x)
            =
            \sum_{i=0}^{n}y_i
            \prod_{\substack{0\leq j\leq n\\j\neq i}}
            \frac{x-x_j}{x_i-x_j}.
        \]
    \end{formula}

    \begin{formula}[Vi\`ete's Formulas]
        If
        \[
            a_nx^n+a_{n-1}x^{n-1}+\cdots+a_0
            =
            a_n\prod_{i=1}^{n}(x-r_i),
        \]
        where the roots are counted with multiplicity, then
        \[
            \sum_{i=1}^{n}r_i=-\frac{a_{n-1}}{a_n},
            \qquad
            \prod_{i=1}^{n}r_i=(-1)^n\frac{a_0}{a_n}.
        \]
        More generally, the elementary symmetric sums of the roots are the
        coefficients with alternating signs, divided by $a_n$.
    \end{formula}


    \begin{formula}[Newton's Identities; Newton--Girard Formulas]
        Let $e_k$ be the elementary symmetric polynomials in roots
        $r_1,\dots,r_n$, and let
        \[
            p_k=r_1^k+\cdots+r_n^k.
        \]
        With $e_0=1$, for $1\leq k\leq n$,
        \[
            p_k-e_1p_{k-1}+e_2p_{k-2}-\cdots
            +(-1)^{k-1}e_{k-1}p_1+(-1)^kke_k=0.
        \]
    \end{formula}

    \begin{formula}[Polynomial Discriminant]
        If
        \[
            f(x)=a_n\prod_{i=1}^{n}(x-r_i),
        \]
        then
        \[
            \Disc(f)
            =
            a_n^{2n-2}\prod_{1\leq i<j\leq n}(r_i-r_j)^2.
        \]
        It vanishes if and only if $f$ has a repeated root.
    \end{formula}

    \begin{theorem}[Rational Root Theorem]
        Let
        \[
            f(x)=a_nx^n+\cdots+a_0\in\Z[x].
        \]
        If a rational number $p/q$ in lowest terms is a root of $f$, then
        \[
            p\mid a_0,
            \qquad
            q\mid a_n.
        \]
    \end{theorem}

    \begin{theorem}[Eisenstein's Irreducibility Criterion]
        Let
        \[
            f(x)=a_nx^n+\cdots+a_0\in\Z[x].
        \]
        If there exists a prime $p$ such that
        \[
            p\nmid a_n,
            \qquad
            p\mid a_i\quad(0\leq i<n),
            \qquad
            p^2\nmid a_0,
        \]
        then $f$ is irreducible over $\Q$.
    \end{theorem}


    \begin{theorem}[Fundamental Theorem of Symmetric Polynomials]
        Every symmetric polynomial can be expressed uniquely as a polynomial in
        the elementary symmetric polynomials.
    \end{theorem}

    \begin{theorem}[Combinatorial Nullstellensatz, Coefficient Form]
        Let $P\in F[x_1,\dots,x_n]$ have total degree
        $t_1+\cdots+t_n$, and suppose the coefficient of
        $x_1^{t_1}\cdots x_n^{t_n}$ is nonzero. If
        \[
            |S_i|>t_i
        \]
        for every finite set $S_i\subseteq F$, then some
        $(s_1,\dots,s_n)\in S_1\times\cdots\times S_n$ satisfies
        \[
            P(s_1,\dots,s_n)\neq0.
        \]
    \end{theorem}

    \begin{strategy}[Polynomial Method]
        Encode a forbidden configuration as zeros of a polynomial, then use a
        degree bound, interpolation, factorization, or a nonzero coefficient to
        prove that the polynomial cannot vanish everywhere.
    \end{strategy}

    \begin{theorem}[Fundamental Theorem of Algebra]
        Every nonconstant complex polynomial has at least one complex root.
        Equivalently, every polynomial of degree $n\geq 1$ over $\C$ splits into
        $n$ linear factors over $\C$, counted with multiplicity.
    \end{theorem}

    \begin{definition}[Algebraically Closed Field]
        A field $F$ is algebraically closed if every nonconstant polynomial in
        $F[x]$ has a root in $F$.
    \end{definition}

    \begin{theorem}[Fundamental Theorem of Galois Theory; Galois Correspondence]
        For a finite Galois extension $L/K$, there is an inclusion-reversing
        correspondence between intermediate fields $K\subseteq E\subseteq L$ and
        subgroups $H\leq \Gal(L/K)$, given by
        \[
            E\longmapsto \Gal(L/E),
            \qquad
            H\longmapsto L^H.
        \]
    \end{theorem}

    \begin{theorem}[Abel--Ruffini Theorem]
        There is no formula for the roots of a general polynomial of degree $5$
        or higher using only the field operations and radicals.
    \end{theorem}

    \begin{remark}
        The Abel--Ruffini theorem does not say that every quintic equation is
        unsolvable. It says that no universal radical formula exists for the
        general quintic.
    \end{remark}

    \chapter{Topology}

    \relatedcourses{(MATH321) General Topology}

    \relatedbooks{Topology (James R. Munkres)}

    \begin{definition}[Topology]
        A topology on a set $X$ is a collection $\mathcal{T}$ of subsets of $X$
        whose elements are called open sets, satisfying:
        \begin{enumerate}
            \item[(i)] $\varnothing,X\in\mathcal{T}$;
            \item[(ii)] arbitrary unions of sets in $\mathcal{T}$ belong to
            $\mathcal{T}$;
            \item[(iii)] finite intersections of sets in $\mathcal{T}$ belong to
            $\mathcal{T}$.
        \end{enumerate}
    \end{definition}

    \begin{definition}[Open and Closed Sets]
        If $(X,\mathcal{T})$ is a topological space, a subset $U\subseteq X$ is
        open if $U\in\mathcal{T}$. A subset $F\subseteq X$ is closed if
        $X\setminus F\in\mathcal{T}$.
    \end{definition}

    \begin{definition}[Compactness]
        A topological space $X$ is compact if every open cover has a finite
        subcover. Equivalently, for every family $\mathcal{U}\subseteq\mathcal{T}$,
        \[
            X\subseteq\bigcup_{U\in\mathcal{U}}U
            \implies
            \exists U_1,\dots,U_n\in\mathcal{U}
            \text{ such that }
            X\subseteq U_1\cup\cdots\cup U_n.
        \]
    \end{definition}

    \begin{definition}[Connectedness]
        A topological space $X$ is connected if it cannot be written as the union
        of two disjoint nonempty open sets. Equivalently, the only subsets of
        $X$ that are both open and closed are $\varnothing$ and $X$.
    \end{definition}


    \begin{theorem}[Jordan Curve Theorem]
        Every simple closed curve in the plane separates the plane into exactly
        two connected components, an interior and an exterior, and is the
        common boundary of both.
    \end{theorem}

    \begin{definition}[Continuity in Topological Spaces]
        A map $f:(X,\mathcal{T}_X)\to(Y,\mathcal{T}_Y)$ is continuous if the
        inverse image of every open set is open:
        \[
            \forall V\in\mathcal{T}_Y,
            \qquad
            f^{-1}(V)\in\mathcal{T}_X.
        \]
    \end{definition}

    \begin{definition}[Homeomorphism]
        A homeomorphism is a bijective continuous map whose inverse is also
        continuous. Equivalently, $f:X\xrightarrow{\sim}Y$ is a homeomorphism
        if $f$ and $f^{-1}$ are continuous. Two spaces are homeomorphic, written
        $X\cong Y$, if such a map exists.
    \end{definition}

    \begin{theorem}[Brouwer Fixed-Point Theorem]
        Every continuous map from a closed ball in $\R^n$ to itself has at least
        one fixed point.
    \end{theorem}


    \begin{theorem}[Borsuk--Ulam Theorem]
        Every continuous map
        \[
            f:S^n\to\R^n
        \]
        takes some pair of antipodal points to the same value.
    \end{theorem}

    \begin{corollary}[Ham Sandwich Theorem]
        Any $n$ bounded measurable subsets of $\R^n$ can be bisected
        simultaneously by one hyperplane.
    \end{corollary}

    \begin{theorem}[Hairy Ball Theorem]
        Every continuous tangent vector field on an even-dimensional sphere
        $S^{2m}$ vanishes at some point.
    \end{theorem}

    \begin{theorem}[Invariance of Domain]
        If $U\subseteq\R^n$ is open and $f:U\to\R^n$ is continuous and injective,
        then $f(U)$ is open and $f$ is a homeomorphism from $U$ onto $f(U)$.
    \end{theorem}

    \begin{definition}[Manifold]
        An $n$-dimensional topological manifold is a space that locally looks like
        $\R^n$. More precisely, each point has a neighborhood homeomorphic to an
        open subset of $\R^n$.
    \end{definition}

    \begin{definition}[Fundamental Group]
        The fundamental group $\pi_1(X,x_0)$ records loops in $X$ based at
        $x_0$, up to continuous deformation. It is one of the basic algebraic
        invariants of a topological space.
    \end{definition}

    \begin{fact}[Fundamental Groups of Familiar Spaces]
        The following basic examples are worth recognizing:
        \[
            \pi_1(S^1)\cong\Z,
            \qquad
            \pi_1(S^n)\cong 0\quad(n\geq 2),
            \qquad
            \pi_1(\mathbb{T}^2)\cong\Z^2.
        \]
    \end{fact}

    \begin{strategy}
        For introductory topology questions, the most common quick checks are:
        compactness means every open cover has a finite subcover, connectedness
        means the space cannot be separated into two disjoint nonempty open
        sets, and a homeomorphism is a continuous bijection with continuous
        inverse.
    \end{strategy}

    \begin{formula}[Euler Characteristic of a Surface]
        For a triangulated closed surface, the Euler characteristic is
        \[
            \chi=V-E+F.
        \]
        For an orientable closed surface of genus $g$,
        \[
            \chi=2-2g.
        \]
    \end{formula}

    \begin{theorem}[Classification of Closed Surfaces]
        Every connected closed surface is homeomorphic to either an orientable
        surface of genus $g\geq 0$ or a nonorientable surface obtained as a
        connected sum of $k\geq 1$ projective planes.
    \end{theorem}

    \begin{theorem}[Poincaré Conjecture]
        Every closed simply connected $3$-manifold is homeomorphic to the
        $3$-sphere $S^3$.
    \end{theorem}

    \begin{remark}
        The Poincaré conjecture was proved by Grigori Perelman using Ricci flow.
        The fundamental correspondence is
        \[
            \text{Poincaré conjecture} \longleftrightarrow
            \text{Perelman} \longleftrightarrow
            \text{Ricci flow}.
        \]
    \end{remark}

    \begin{theorem}[Thurston Geometrization Theorem]
        Every closed orientable $3$-manifold can be decomposed into pieces that
        admit one of the eight Thurston geometries.
    \end{theorem}

    \begin{fact}[The Eight Thurston Geometries]
        The eight model geometries are
        \[
            S^3,
            \quad \mathrm{E}^3,
            \quad \mathbb{H}^3,
            \quad S^2\times\R,
            \quad \mathbb{H}^2\times\R,
            \quad \widetilde{\SL}_2(\R),
            \quad \mathrm{Nil},
            \quad \mathrm{Sol}.
        \]
    \end{fact}

    \chapter{Miscellaneous Topics}

    \relatedcourses{(MATH201) Introduction to Mathematics,
        (MATH261) Discrete Mathematics}

    \relatedbooks{Problem-Solving Strategies (Arthur Engel),
        The Art and Craft of Problem Solving (Paul Zeitz)}

    \begin{strategy}[Transform the Given Objects]
        Replace the original configuration by an equivalent one whose structure
        is easier to see: partition a region, introduce prefix sums, encode a
        movement by vectors, pass to residues, diagonalize a matrix, or complete
        a partial object. The transformed problem should preserve the quantity or
        property being asked about.
    \end{strategy}

    \begin{strategy}[Invariant, Monovariant, and Extremal Principles]
        An invariant remains unchanged under every allowed move. A monovariant
        moves strictly in one direction and therefore prevents cycles or proves
        termination. The extremal principle chooses a largest, smallest, first,
        or last object and exploits the extra restrictions forced by extremality.
        These are among the most common contest ideas hidden behind elementary
        statements.
    \end{strategy}

    \begin{strategy}[Complete the Configuration and Work Backward]
        First imagine a completed solution, then determine which parts are forced
        and which choices remain free. This is effective in geometric
        constructions, permutation completions, lattice paths, tilings, and
        assignment problems. Fixing one object under a symmetry often removes
        redundant cases before the remaining count begins.
    \end{strategy}

    \begin{strategy}[Rotate, Reflect, and Superimpose]
        A rotation, reflection, or translated copy can convert a broken path into
        a straight one, expose equal vectors, pair terms, or turn a local pattern
        into a global invariant. When two configurations overlap, compare what
        cancels and what remains rather than recomputing each separately.
    \end{strategy}


    \begin{strategy}[Homogenization and Normalization]
        If an inequality is homogeneous, scale the variables to impose a simple
        condition such as $a+b+c=1$, $abc=1$, or $x^2+y^2=1$.
    \end{strategy}

    \begin{strategy}[Equality Cases and Smoothing]
        Identify the expected equality case before selecting an inequality.
        For a symmetric expression, replacing two variables by their average
        often reveals the extremal configuration.
    \end{strategy}

    \begin{theorem}[Triangle and Reverse Triangle Inequalities]
        For real or complex numbers $x,y$,
        \[
            \bigl||x|-|y|\bigr|
            \leq
            |x-y|
            \leq
            |x|+|y|.
        \]
        More generally, every norm satisfies
        \[
            \|\mathbf{x}+\mathbf{y}\|
            \leq
            \|\mathbf{x}\|+\|\mathbf{y}\|.
        \]
    \end{theorem}

    \begin{formula}[Basic Sum-of-Squares Identities]
        For real $a,b,c$,
        \[
            a^2+b^2+c^2-ab-bc-ca
            =
            \frac12\left((a-b)^2+(b-c)^2+(c-a)^2\right)
            \geq0.
        \]
        Consequently,
        \[
            a^2+b^2+c^2\geq ab+bc+ca
        \]
        and
        \[
            (a+b+c)^2\leq3(a^2+b^2+c^2).
        \]
    \end{formula}

    \begin{formula}[Frequently Used Two-Variable Inequalities]
        For real $x,y$,
        \[
            x^2+y^2\geq2xy.
        \]
        For positive $x,y$,
        \[
            \frac{x}{y}+\frac{y}{x}\geq2,
            \qquad
            x+\frac1x\geq2.
        \]
        Equality holds exactly when the two compared positive quantities are
        equal.
    \end{formula}

    \begin{theorem}[Lagrange's Identity]
        For real numbers $a_1,\dots,a_n$ and $b_1,\dots,b_n$,
        \[
            \left(\sum_{i=1}^{n}a_i^2\right)
            \left(\sum_{i=1}^{n}b_i^2\right)
            -
            \left(\sum_{i=1}^{n}a_ib_i\right)^2
            =
            \sum_{1\leq i<j\leq n}(a_ib_j-a_jb_i)^2.
        \]
        The nonnegativity of the right-hand side gives the
        Cauchy--Schwarz inequality immediately.
    \end{theorem}

    \begin{theorem}[QM--AM--GM--HM Inequalities]
        For positive $x_1,\dots,x_n$,
        \[
            \sqrt{\frac{x_1^2+\cdots+x_n^2}{n}}
            \geq
            \frac{x_1+\cdots+x_n}{n}
            \geq
            \sqrt[n]{x_1\cdots x_n}
            \geq
            \frac{n}{\frac1{x_1}+\cdots+\frac1{x_n}}.
        \]
        Equality holds if and only if all variables are equal.
    \end{theorem}

    \begin{theorem}[Weighted AM--GM Inequality]
        If $w_i\geq0$, $\sum_{i=1}^{n}w_i=1$, and $x_i>0$, then
        \[
            \sum_{i=1}^{n}w_ix_i
            \geq
            \prod_{i=1}^{n}x_i^{w_i}.
        \]
        Equality holds when all $x_i$ with positive weight are equal.
    \end{theorem}

    \begin{theorem}[Power Mean Inequality]
        For positive $x_1,\dots,x_n$ and real $r\neq0$, define
        \[
            M_r
            =
            \left(
                \frac{x_1^r+\cdots+x_n^r}{n}
            \right)^{1/r},
        \]
        and define $M_0=(x_1\cdots x_n)^{1/n}$. If $r<s$, then
        \[
            M_r\leq M_s.
        \]
        The harmonic, geometric, arithmetic, and quadratic means correspond to
        $r=-1,0,1,2$, respectively.
    \end{theorem}

    \begin{theorem}[Young's Inequality]
        If $x,y\geq0$, $p,q>1$, and
        \[
            \frac1p+\frac1q=1,
        \]
        then
        \[
            xy\leq\frac{x^p}{p}+\frac{y^q}{q}.
        \]
    \end{theorem}

    \begin{theorem}[Cauchy--Schwarz Inequality, Engel Form]
        For real $a_i$ and positive $b_i$,
        \[
            \sum_{i=1}^{n}\frac{a_i^2}{b_i}
            \geq
            \frac{(a_1+\cdots+a_n)^2}{b_1+\cdots+b_n}.
        \]
    \end{theorem}

    \begin{theorem}[H\"older's Inequality]
        If $p,q>1$ and $1/p+1/q=1$, then
        \[
            \sum_{i=1}^{n}|a_ib_i|
            \leq
            \left(\sum_{i=1}^{n}|a_i|^p\right)^{1/p}
            \left(\sum_{i=1}^{n}|b_i|^q\right)^{1/q}.
        \]
    \end{theorem}

    \begin{theorem}[Minkowski's Inequality]
        For $p\geq1$,
        \[
            \left(\sum_{i=1}^{n}|a_i+b_i|^p\right)^{1/p}
            \leq
            \left(\sum_{i=1}^{n}|a_i|^p\right)^{1/p}
            +
            \left(\sum_{i=1}^{n}|b_i|^p\right)^{1/p}.
        \]
    \end{theorem}

    \begin{theorem}[Rearrangement Inequality]
        If
        \[
            a_1\leq\cdots\leq a_n,
            \qquad
            b_1\leq\cdots\leq b_n,
        \]
        then for every permutation $\sigma\in S_n$,
        \[
            \sum_{i=1}^{n}a_ib_{n+1-i}
            \leq
            \sum_{i=1}^{n}a_ib_{\sigma(i)}
            \leq
            \sum_{i=1}^{n}a_ib_i.
        \]
    \end{theorem}

    \begin{theorem}[Chebyshev's Sum Inequality]
        If $a_1\leq\cdots\leq a_n$ and $b_1\leq\cdots\leq b_n$, then
        \[
            \frac1n\sum_{i=1}^{n}a_ib_i
            \geq
            \left(\frac1n\sum_{i=1}^{n}a_i\right)
            \left(\frac1n\sum_{i=1}^{n}b_i\right).
        \]
    \end{theorem}

    \begin{theorem}[Schur's Inequality]
        For $a,b,c\geq0$ and $r\geq0$,
        \[
            \sum_{\mathrm{cyc}}a^r(a-b)(a-c)\geq0.
        \]
    \end{theorem}

    \begin{formula}[Nesbitt's Inequality]
        For positive $a,b,c$,
        \[
            \frac{a}{b+c}+\frac{b}{c+a}+\frac{c}{a+b}
            \geq
            \frac32.
        \]
    \end{formula}

    \begin{theorem}[Bernoulli's Inequality]
        If $x\geq-1$ and $r\geq1$, then
        \[
            (1+x)^r\geq1+rx.
        \]
    \end{theorem}

    \begin{strategy}[Comparing Powers by Logarithms]
        To compare two positive numbers such as $a^b$ and $b^a$, take logarithms:
        \[
            a^b>b^a
            \quad\Longleftrightarrow\quad
            b\ln a>a\ln b
            \quad\Longleftrightarrow\quad
            \frac{\ln a}{a}>\frac{\ln b}{b}.
        \]
        The function $\ln x/x$ is decreasing for $x>e$. This quickly handles
        comparisons such as $11^9$ versus $9^{11}$ or $e^\pi$ versus $\pi^e$.
    \end{strategy}

    \begin{formula}[Mercator Series for $\ln 2$]
        The alternating harmonic series satisfies
        \[
            \sum_{n=1}^{\infty}(-1)^{n+1}\frac{1}{n}
            =
            1-\frac12+\frac13-\frac14+\cdots
            =
            \ln 2.
        \]
    \end{formula}

    \begin{remark}
        This follows from the Mercator series
        \[
            \ln(1+x)
            =
            x-\frac{x^2}{2}+\frac{x^3}{3}-\frac{x^4}{4}+\cdots
        \]
        by substituting $x=1$.
    \end{remark}

    \begin{formula}[Basel Sum; Euler's Basel Formula]
        The Basel sum is
        \[
            \sum_{n=1}^{\infty}\frac{1}{n^2}
            =
            1+\frac{1}{2^2}+\frac{1}{3^2}+\cdots
            =
            \frac{\pi^2}{6}.
        \]
        Equivalently,
        \[
            \zeta(2)=\frac{\pi^2}{6}.
        \]
    \end{formula}

    \begin{remark}
        The problem of evaluating
        \[
            \sum_{n=1}^{\infty}\frac{1}{n^2}
        \]
        is called the Basel problem. Its famous solution was found by Euler.
    \end{remark}

    \begin{formula}[Odd Part of the Basel Sum]
        The sum of reciprocal squares over odd positive integers is
        \[
            \sum_{n=1}^{\infty}\frac{1}{(2n-1)^2}
            =
            1+\frac{1}{3^2}+\frac{1}{5^2}+\cdots
            =
            \frac{\pi^2}{8}.
        \]
    \end{formula}

    \begin{idea}
        This is obtained by subtracting the even part from the Basel sum:
        \[
            \sum_{n=1}^{\infty}\frac{1}{(2n-1)^2}
            =
            \sum_{n=1}^{\infty}\frac{1}{n^2}
            -
            \sum_{n=1}^{\infty}\frac{1}{(2n)^2}
            =
            \left(1-\frac14\right)\frac{\pi^2}{6}
            =
            \frac{\pi^2}{8}.
        \]
    \end{idea}

    \begin{formula}[Alternating Basel Series; Dirichlet Eta Value at $2$]
        The alternating Basel series satisfies
        \[
            \sum_{n=1}^{\infty}(-1)^{n+1}\frac{1}{n^2}
            =
            1-\frac{1}{2^2}+\frac{1}{3^2}-\frac{1}{4^2}+\cdots
            =
            \frac{\pi^2}{12}.
        \]
        Equivalently,
        \[
            \eta(2)=\frac{\pi^2}{12},
        \]
        where $\eta(s)$ is the Dirichlet eta function.
    \end{formula}

    \begin{idea}
        The positive terms are the odd reciprocal squares and the negative
        terms are the even reciprocal squares. Hence
        \[
            \sum_{n=1}^{\infty}(-1)^{n+1}\frac{1}{n^2}
            =
            \sum_{n=1}^{\infty}\frac{1}{(2n-1)^2}
            -
            \sum_{n=1}^{\infty}\frac{1}{(2n)^2}
            =
            \frac{\pi^2}{8}-\frac{\pi^2}{24}
            =
            \frac{\pi^2}{12}.
        \]
    \end{idea}

    \begin{formula}[Gregory--Leibniz Series; Leibniz Formula for $\pi$]
        The Gregory--Leibniz series is
        \[
            \sum_{n=1}^{\infty}(-1)^{n+1}\frac{1}{2n-1}
            =
            1-\frac13+\frac15-\frac17+\cdots
            =
            \frac{\pi}{4}.
        \]
        Equivalently,
        \[
            \pi
            =
            4\left(
                1-\frac13+\frac15-\frac17+\cdots
            \right).
        \]
    \end{formula}

    \begin{remark}
        This follows from the Taylor series
        \[
            \arctan x
            =
            x-\frac{x^3}{3}+\frac{x^5}{5}-\frac{x^7}{7}+\cdots
        \]
        by substituting $x=1$, since $\arctan 1=\pi/4$.
    \end{remark}

    \begin{formula}[More Useful Series Values]
        The following values are common in fast calculation problems:
        \[
            \sum_{n=1}^{\infty}\frac{1}{n(n+1)}=1,
            \qquad
            \sum_{n=1}^{\infty}\frac{1}{n(n+2)}=\frac34,
        \]
        \[
            \sum_{n=1}^{\infty}\frac{1}{n^4}=\frac{\pi^4}{90},
            \qquad
            \sum_{n=1}^{\infty}\frac{1}{(2n-1)^4}=\frac{\pi^4}{96}.
        \]
    \end{formula}

    \begin{formula}[Power-Series Sums]
        For $|x|<1$,
        \[
            \sum_{n=0}^{\infty}x^n=\frac{1}{1-x},
            \qquad
            \sum_{n=1}^{\infty}nx^{n}=\frac{x}{(1-x)^2},
        \]
        and
        \[
            \sum_{n=1}^{\infty}n^2x^{n}
            =
            \frac{x(1+x)}{(1-x)^3}.
        \]
    \end{formula}

    \begin{strategy}
        These values often appear after a simple transformation. For example,
        \[
            \frac12-\frac14+\frac16-\frac18+\cdots
            =
            \frac12
            \left(
                1-\frac12+\frac13-\frac14+\cdots
            \right)
            =
            \frac{\ln 2}{2}.
        \]
    \end{strategy}
